\documentclass[11pt]{report}
\usepackage[a4paper,margin=1.05in]{geometry}
\usepackage[T1]{fontenc}
\usepackage{amsmath,amssymb,amsthm,mathtools}
\usepackage{stmaryrd}
\usepackage{mathpartir}
\usepackage{enumitem}
\usepackage{array}
\usepackage{xcolor}
\usepackage[colorlinks=true,linkcolor=blue!55!black,citecolor=blue!55!black,urlcolor=blue!55!black]{hyperref}

%%% Theorem environments
\theoremstyle{plain}
\newtheorem{theorem}{Theorem}[chapter]
\newtheorem{lemma}[theorem]{Lemma}
\newtheorem{corollary}[theorem]{Corollary}
\theoremstyle{definition}
\newtheorem{definition}[theorem]{Definition}
\newtheorem{example}[theorem]{Example}
\newtheorem{choice}{Choice}[chapter]
\theoremstyle{remark}
\newtheorem{remark}[theorem]{Remark}

%%% Roles and basics
\newcommand{\kw}[1]{\mathsf{#1}}
\newcommand{\pp}{\mathsf{p}}
\newcommand{\qq}{\mathsf{q}}
\newcommand{\rr}{\mathsf{r}}
\newcommand{\sr}{\mathsf{s}}
\newcommand{\tr}{\mathsf{t}}
\newcommand{\Roles}{\mathcal{R}}
\newcommand{\ttrue}{\kw{tt}}
\newcommand{\ffalse}{\kw{ff}}
\newcommand{\fv}{\mathrm{fv}}

%%% Choreographies
\newcommand{\bind}[3]{#1 \leftarrow #2\,;\;#3}
\newcommand{\cret}[1]{\kw{ret}\;#1}
\newcommand{\com}[3]{#1.#2 \to #3}
\newcommand{\cif}[4]{\kw{if}\;#1.#2\;\kw{then}\;#3\;\kw{else}\;#4}
\newcommand{\mbind}[3]{#1 \mathbin{>\!\!>\!\!=} #2.\,#3}
\newcommand{\rt}[3]{\langle #1 : #2\rangle\,#3}
\newcommand{\unfold}[1]{\mathrm{unfold}(#1)}
\newcommand{\roles}[1]{\mathrm{roles}(#1)}
\newcommand{\rcv}[1]{\mathrm{rcv}(#1)}
\newcommand{\loc}[1]{\mathrm{loc}(#1)}
\newcommand{\WF}{\mathrm{WF}}
\newcommand{\pay}[1]{\lvert #1\rvert}

%%% Labels and steps
\newcommand{\tauat}[1]{\tau\mathord{@}#1}
\newcommand{\castlbl}[3]{#1.#2 \rightsquigarrow #3}
\newcommand{\lto}[1]{\xrightarrow{\;#1\;}}
\newcommand{\seqto}{\to_{\mathrm{seq}}}

%%% Networks
\newcommand{\nsend}[2]{\kw{send}\;#1\;#2}
\newcommand{\nrecv}[1]{\kw{recv}\;#1}
\newcommand{\sendbr}[4]{#1!#2 \mathbin{\triangleright} (#3,\,#4)}
\newcommand{\recvbr}[3]{#1?(#2,\,#3)}
\newcommand{\upd}[3]{#1[#2 \mapsto #3]}
\newcommand{\nfun}[3]{\kw{fun}\;#1(#2).\;#3}

%%% Projection
\newcommand{\proj}[2]{\llbracket #1 \rrbracket_{#2}}
\newcommand{\projn}[1]{\llbracket #1 \rrbracket}
\newcommand{\pre}[3]{\llbracket #1 \rrbracket^{#2}_{#3}}

%%% Functions
\newcommand{\Sig}{\Sigma}
\newcommand{\Defs}{\mathcal{D}}
\newcommand{\farr}[3]{(#1) \xrightarrow{#2} #3}
\newcommand{\cfun}[5]{\kw{fun}\;#1(#2) \xrightarrow{#3} #4.\;#5}
\newcommand{\clam}[3]{\lambda #1{:}#2.\;#3}
\newcommand{\cfix}[3]{\kw{fix}\;#1{:}#2.\;#3}
\newcommand{\nlam}[2]{\lambda(#1).\;#2}
\newcommand{\nfix}[2]{\kw{fix}\;#1.\;#2}
\newcommand{\unroll}[1]{\mathrm{unroll}(#1)}

%%% Proof layout
\newcommand{\rn}[1]{{\normalfont\scshape #1}}
\newcommand{\case}[1]{\medskip\noindent\textbf{Case #1.}\;}
\newcommand{\subcase}[1]{\smallskip\noindent\textit{Subcase #1.}\;}

\title{Monadic Choreographies}

\begin{document}
\maketitle

\begin{abstract}
  We develop four choreographic languages of increasing expressiveness:
  local computation and communication; conditionals; first-order recursive
  functions; and first-class, higher-order functions.  A fifth language pares
  the last down to a minimal core, with one sort of values, single-argument
  anonymous functions and a fixpoint.  All five are written
  in the normal form of a free monad, as a sequence of actions $\bind{x}{A}{C}$
  ending in a return, and each of the first four adds its feature as a new
  kind of action whose result the continuation binds.  For each language, endpoint projection is a
  strong bisimulation with identical labels whose network side is the
  projection of the choreography itself: one choreography step matches one
  network step, with no pruning relation, weak steps, or equivalence up to the
  unfolding of calls.  We chose the languages to keep this proof short, and
  we justify each design decision by the proof obligation it removes.
\end{abstract}

\tableofcontents

%======================================================================
\chapter*{Introduction}
\addcontentsline{toc}{chapter}{Introduction}
%======================================================================

A choreography describes a distributed computation from a global point of
view, as a single program in which each communication is an atomic
interaction between two roles.  \emph{Endpoint projection} (EPP) compiles a
choreography to a \emph{network}: one local program per role, communicating
by message passing.  The central theorem of choreographic programming is that
EPP is correct, meaning that a choreography and its projection behave the
same.

In most presentations the correspondence is looser than that.  One
choreography step may correspond to several network steps, and a network may
only be related to the projection of a choreography up to a \emph{pruning}
relation, which discards branches that can no longer be taken, or up to the
unfolding of procedure calls.  Each of these gaps needs its own auxiliary
relation and lemmas about it, and the resulting proofs are hard to check.

This document develops five languages without such gaps.  For each of them
the relation
\[
  \{\,(C, \projn{C}) \mid C \text{ is well-typed (and, from Chapter~\ref{ch:fun} on, well-formed)}\,\}
\]
is a strong bisimulation with identical labels: one choreography step
corresponds to one network step with the same label, and the network reached
is the projection of the choreography reached.  We started from this goal and
chose the languages to meet it.  Each chapter explains its design choices by
the proof obligations they remove, and numbers them so that later chapters
can refer back.

\paragraph{The monadic style.}
A choreography is a sequence of actions, each binding its result, ending in a
return:
\[
  C \;::=\; \cret{\pp.a} \;\mid\; \bind{x}{A}{C}.
\]
This is the normal form of the free monad over the actions.  Chapter~\ref{ch:core}
has two actions, local computation and communication.  Each later chapter
adds a new kind of action: a conditional (Chapter~\ref{ch:cond}), a call to a
named function (Chapter~\ref{ch:fun}), and the application of a first-class
function (Chapter~\ref{ch:ho}).  Each new action returns a value that the
continuation binds, so the language gains conditionals and functions that
return values, and functions that are passed as arguments and returned as
results.  The bind and the \rn{Delay} rule, which lets independent actions run
out of order, stay the same throughout; Chapter~\ref{ch:fun} adds a companion
rule, \rn{Inside}, for calls in progress.  Chapter~\ref{ch:min} adds no
action.  It keeps those of Chapter~\ref{ch:ho} and removes everything they
do not need.

\paragraph{Proof architecture.}
Every chapter proves the same theorem with the same skeleton.
\begin{itemize}[nosep]
\item \emph{Completeness}: if $C\lto{\ell}C'$ then $\projn{C}\lto{\ell}\projn{C'}$, by induction on the step.
\item \emph{Soundness}: if $\projn{C}\lto{\ell}N$ then $C\lto{\ell}C'$ for some $C'$ with $N=\projn{C'}$, by induction on $C$.
\end{itemize}
Both rest on the same supporting lemmas: substitution and subject reduction;
\emph{projection is local} (a role's program mentions only that role's
variables); \emph{projection commutes with substitution}; a
\emph{choreographic frame} lemma (a step does not change the projection at a
role it does not involve); a \emph{network frame} lemma (the same for
networks); and \emph{inversion} for network steps.  Each later chapter
re-proves these lemmas by adding cases for its new constructs, and adds a few
new lemmas of its own.

\paragraph{Roadmap.}
Chapter~\ref{ch:core} treats local computation and synchronous
communication.  Chapter~\ref{ch:cond} adds conditionals, implemented by
multicasting the guard to every participant of either branch; there are no
selections and no merging.  Chapter~\ref{ch:fun} adds first-order recursive
functions.  Calls are entered one role at a time, which requires one runtime
term and a well-formedness invariant.  Chapter~\ref{ch:ho} removes named
functions altogether and replaces them with first-class recursive functions,
each located at the set of roles that run it.  Chapter~\ref{ch:min} reduces
that language to a minimal core: one sort of values, functions of one
argument without names, and a fixpoint that is unrolled when a role enters an
application, so that recursion costs no step of its own.

%======================================================================
\chapter{Core Choreographies}
\label{ch:core}
%======================================================================

The first language has two actions, local computation and synchronous
communication, sequenced by a monadic bind.  Endpoint projection is a strong
bisimulation with identical labels: every network reachable from the
projection of a choreography is the projection of a choreography reachable
from it (Theorem~\ref{thm:c1-epp}).  With confluence
(Theorem~\ref{thm:c1-agree}), it follows that every run of the network
computes the result of the sequential reading of the choreography.

\section{Design choices}
\label{sec:c1-choices}

A step of the network is taken by one role (a local computation) or by two
roles together (a communication), and roles not involved in it are untouched.
For a choreography to match the network step for step, three things must
hold: the choreography must be able to perform any action its roles are ready
for, regardless of program order; an action must take exactly one step; and a
role that is not involved in an action must not be required to do anything for
it.  The following choices secure these properties.

\begin{choice}[Actions are fused with bind]\label{cho:c1-fused}
  The syntax is the free monad's normal form $\bind{x}{A}{C}$: every action is
  immediately followed by the continuation that binds its result.  General
  bind is a derived operation on syntax (Definition~\ref{def:c1-bind}), not a
  term former.

  \emph{Proof saved.}  An action step resumes its continuation directly.
  Neither semantics has a step that reduces a bind, so neither has steps that
  are administrative relative to the other.  Projection maps
  $\bind{x}{A}{C}$ to a fixed \emph{prefix} followed by the projection of $C$
  (Definition~\ref{def:c1-epp}); most proof cases for \rn{Delay} are then the
  observation that both sides have the same prefix.
\end{choice}

\begin{choice}[Projection erases uninvolved actions]\label{cho:c1-erase}
  A role that is not involved in an action gets nothing for it, rather than a
  bind of a dummy value such as $\bind{x}{\kw{ret}\;\bot}{\cdots}$.

  \emph{Proof saved.}  Uninvolved roles take no administrative steps, so one
  choreography step is matched by one network step.  The frame lemmas (Lemmas~\ref{lem:c1-cframe}
  and~\ref{lem:c1-nframe}) become one-line arguments.  The cost is that the
  bound variable disappears from the programs of uninvolved roles, which is
  safe because typing guarantees they never mention it
  (Lemma~\ref{lem:c1-local}).
\end{choice}

\begin{choice}[Substitution, not stores]\label{cho:c1-subst}
  Values flow by substitution: \rn{Local} and \rn{Com} substitute the
  computed or received value for the bound variable, at both levels.  There
  is no store.

  \emph{Proof saved.}  Configurations are just terms.  There is no
  store-locality lemma; locality is instead a typing fact
  (Lemma~\ref{lem:c1-local}), and a single lemma relates substitution at the
  two levels (Lemma~\ref{lem:c1-psubst}).
\end{choice}

\begin{choice}[Values are required only where they appear in labels]\label{cho:c1-values}
  Local computation takes an expression and evaluates it.  Communication
  takes an \emph{atom} (a variable or a value) and fires only when the atom
  is a value; that value appears in the label.

  \emph{Proof saved.}  \rn{Com} has no premise, and its label determines its
  effect, which gives label determinacy (Lemma~\ref{lem:c1-det}).  \rn{Delay}
  needs no side condition on variables (Remark~\ref{rem:c1-delay}): an action
  that uses a variable bound by a pending action cannot fire, because its
  expression is not closed or its atom is not a value.  Later chapters
  extend this principle: the guard of a conditional must be a value because
  it appears in the label, but nothing in a function call has to be.
\end{choice}

\begin{choice}[Synchronous communication]\label{cho:c1-sync}
  A communication is a single rendezvous between sender and receiver.

  \emph{Proof saved.}  There are no message queues in networks and no runtime
  term in choreographies for messages that are sent but not yet received.
  One global communication is one network step.
\end{choice}

\begin{choice}[Out-of-order execution by a single rule]\label{cho:c1-delay}
  The rule \rn{Delay} lets an action step under pending binds when its roles
  are disjoint from theirs.  The sequential semantics, which only fires the
  head action, is kept as the specification.

  \emph{Proof saved.}  Roles of a network run independently, so without
  \rn{Delay} the choreography could not match a network step taken by a role
  whose actions come late in program order, and soundness would fail.  With
  it, the choreography is as concurrent as the network, and by confluence
  every run still computes the sequential result
  (Theorem~\ref{thm:c1-agree}).
\end{choice}

\section{The local language}
\label{sec:c1-local}

The local language is a parameter, used both inside choreographies and inside
local programs.  It provides:
\begin{itemize}[nosep]
\item expressions $e$, built from variables $x,y,z$ and closed \emph{values} $v$;
\item capture-avoiding substitution $e[a/x]$ of an atom (defined below) for a
  variable, with $e[a/x]=e$ when $x\notin\fv(e)$, and
  $\fv(e[a/x])\subseteq(\fv(e)\setminus\{x\})\cup\fv(a)$;
\item evaluation $e\Downarrow v$, defined only on closed expressions.
\end{itemize}
\emph{Atoms} are $a ::= x \mid v$, and every atom is an expression.  Two
assumptions are used, each only where noted.
\begin{description}[nosep,font=\normalfont\itshape]
\item[(L1) Determinism.] If $e\Downarrow v$ and $e\Downarrow v'$ then
  $v=v'$.  Used only for confluence (\S\ref{sec:c1-confluence}).
\item[(L2) Totality.] Every closed $e$ evaluates to some $v$.  Used only for
  deadlock freedom (\S\ref{sec:c1-df}) and for agreement with the sequential
  semantics (Theorem~\ref{thm:c1-agree}).
\end{description}
The symbol $\bot$ is not a value.  It appears only in local programs, as the
result a role returns when the result of the choreography is held elsewhere.

\section{Syntax and typing}
\label{sec:c1-syntax}

Roles $\pp,\qq,\rr,\sr,\tr$ are drawn from a fixed finite set $\Roles$.

\begin{definition}[Choreographies]\label{def:c1-syntax}
  Actions $A$ and choreographies $C$ are given by
  \[
  \begin{array}{r@{\;\;}c@{\;\;}l@{\qquad}l}
    A & ::= & \pp.e & \text{local computation at } \pp\\
      & \mid & \com{\pp}{a}{\qq} & \text{communication from $\pp$ to $\qq$, with } \pp\neq\qq\\[1ex]
    C & ::= & \cret{\pp.a} & \text{return a value held at } \pp\\
      & \mid & \bind{x}{A}{C} & \text{run $A$, bind its result to $x$, continue as } C
  \end{array}
  \]
  Each action has a set of roles and a location for its result:
  \[
  \begin{array}{l@{\qquad}l}
    \roles{\pp.e} = \{\pp\} & \loc{\pp.e} = \pp\\
    \roles{\com{\pp}{a}{\qq}} = \{\pp,\qq\} & \loc{\com{\pp}{a}{\qq}} = \qq
  \end{array}
  \]
\end{definition}

In $\bind{x}{A}{C}$ the variable $x$ is bound in $C$.  Terms are taken up to
$\alpha$-equivalence, and $C[a/x]$ is capture-avoiding substitution of an
atom for a variable, which acts on the expressions and atoms inside $C$.
Since values are closed, substituting a value never captures.

\begin{definition}[Bind]\label{def:c1-bind}
  The derived bind $\mbind{C_1}{x}{C_2}$ is defined by recursion on $C_1$:
  \[
  \begin{array}{r@{\;\;=\;\;}l@{\qquad}l}
    \mbind{(\cret{\pp.a})}{x}{C_2} & C_2[a/x]\\
    \mbind{(\bind{y}{A}{C_1})}{x}{C_2} & \bind{y}{A}{(\mbind{C_1}{x}{C_2})} & (y\notin\fv(C_2)\cup\{x\})
  \end{array}
  \]
\end{definition}

Bind walks only the spine of $C_1$, the sequence of actions leading to its
return.  It is associative,
$\mbind{(\mbind{C_1}{x}{C_2})}{y}{C_3} = \mbind{C_1}{x}{(\mbind{C_2}{y}{C_3})}$,
by induction on $C_1$, and $\cret{}$ is a left unit by definition.  It plays no
role in the semantics of this chapter; from Chapter~\ref{ch:cond} on, it is
how the result of a compound action is passed to its continuation.

\paragraph{Location typing.}
Typing checks one thing: each role uses only data it holds.  A context
$\Gamma$ maps variables to roles.  Write $\Gamma\vdash_\pp e$ when $\Gamma(x)=\pp$
for every $x\in\fv(e)$; this applies to atoms too.

\begin{definition}[Typing]\label{def:c1-typing}
  The judgement $\Gamma\vdash C$ is given by
  \begin{mathpar}
    \inferrule{\Gamma\vdash_\pp a}{\Gamma\vdash\cret{\pp.a}}
    \and
    \inferrule{\Gamma\vdash_\pp e \\ \Gamma,x{:}\pp\vdash C}{\Gamma\vdash\bind{x}{\pp.e}{C}}
    \and
    \inferrule{\pp\neq\qq \\ \Gamma\vdash_\pp a \\ \Gamma,x{:}\qq\vdash C}{\Gamma\vdash\bind{x}{\com{\pp}{a}{\qq}}{C}}
  \end{mathpar}
  A \emph{program} is a choreography $C$ with $\emptyset\vdash C$.
\end{definition}

In both action rules the bound variable is held at $\loc{A}$, the role that
holds the action's result.  The metatheory is stated for open terms,
$\Gamma\vdash C$, because the inductions go under binders.  If the local
language is typed, roles in $\Gamma$ can be replaced by located types
$t\mathord{@}\pp$; nothing below changes.

\section{Choreographic semantics}
\label{sec:c1-sem}

The semantics is a labelled transition system on choreographies.  Labels
record which roles acted:
\[
  \ell \;::=\; \tauat{\pp} \;\mid\; \com{\pp}{v}{\qq}
  \qquad
  \roles{\tauat{\pp}} = \{\pp\}
  \qquad
  \roles{\com{\pp}{v}{\qq}} = \{\pp,\qq\}.
\]

\begin{figure}[t]
\begin{mathpar}
  \inferrule*[right=\rn{Local}]
    {e\Downarrow v}
    {\bind{x}{\pp.e}{C} \lto{\tauat{\pp}} C[v/x]}
  \and
  \inferrule*[right=\rn{Com}]
    { }
    {\bind{x}{\com{\pp}{v}{\qq}}{C} \lto{\com{\pp}{v}{\qq}} C[v/x]}
  \and
  \inferrule*[right=\rn{Delay}]
    {C\lto{\ell}C' \\ \roles{\ell}\cap\roles{A}=\emptyset}
    {\bind{x}{A}{C} \lto{\ell} \bind{x}{A}{C'}}
\end{mathpar}
\caption{Semantics of core choreographies.}
\label{fig:c1-sem}
\end{figure}

\begin{definition}[Semantics]\label{def:c1-sem}
  The \emph{sequential semantics} $\seqto$ is the least relation closed under
  \rn{Local} and \rn{Com} of Figure~\ref{fig:c1-sem}; only the head action can
  fire, and it is the specification.  The \emph{out-of-order semantics}
  $\lto{\ell}$ is the least relation closed under all three rules, and it is
  the one projection is proved against.
\end{definition}

\begin{remark}[\rn{Delay} needs no side condition on $x$]\label{rem:c1-delay}
  \rn{Local} fires only on a closed expression and \rn{Com} only when its atom
  is a value, so an action under pending binds steps only if it uses none of
  their variables.  For well-typed terms, role disjointness already implies
  this: a variable bound by $A$ is held at $\loc{A}\in\roles{A}$, and only
  actions at that role may use it.
\end{remark}

\section{Networks}
\label{sec:c1-net}

Each role runs a local program, and a network runs them all, with synchronous
communication.

\begin{definition}[Local programs and networks]\label{def:c1-net}
  Local programs $P$ and prefixes $\pi$ are given by
  \[
  \begin{array}{r@{\;\;}c@{\;\;}l@{\qquad}l}
    P & ::= & \kw{ret}\;a \;\mid\; \kw{ret}\;\bot & \text{return}\\
      & \mid & \pi\,P & \text{prefixed program}\\[1ex]
    \pi & ::= & x\leftarrow e\,; & \text{local computation}\\
      & \mid & \nsend{\qq}{a}\,; & \text{send to } \qq\\
      & \mid & x\leftarrow\nrecv{\pp}\,; & \text{receive from } \pp
  \end{array}
  \]
  We also write $\varepsilon$ for the empty prefix, with $\varepsilon\,P = P$.
  A \emph{network} $N$ is a function mapping every role in $\Roles$ to a local
  program.  $\upd{N}{\pp}{P}$ updates one entry.  Two networks \emph{agree on}
  a set of roles $S$ if they map every role in $S$ to the same program.  A
  network is \emph{terminated} if every $N(\rr)$ is a return.
\end{definition}

\begin{figure}[t]
\begin{mathpar}
  \inferrule*[right=\rn{N-Local}]
    {N(\pp) = x\leftarrow e\,;\,P \\ e\Downarrow v}
    {N \lto{\tauat{\pp}} \upd{N}{\pp}{P[v/x]}}
  \and
  \inferrule*[right=\rn{N-Com}]
    {N(\pp) = \nsend{\qq}{v}\,;\,P \\ N(\qq) = x\leftarrow\nrecv{\pp}\,;\,P'}
    {N \lto{\com{\pp}{v}{\qq}} N[\pp\mapsto P,\ \qq\mapsto P'[v/x]]}
\end{mathpar}
\caption{Semantics of networks.}
\label{fig:c1-net}
\end{figure}

The semantics of networks is given in Figure~\ref{fig:c1-net}.  A step reads
and changes only the programs of the roles in its label
(Lemma~\ref{lem:c1-nframe}).

\section{Endpoint projection}
\label{sec:c1-epp}

Projection is total and needs no merging: each role keeps exactly the actions
it takes part in.

\begin{definition}[Endpoint projection]\label{def:c1-epp}
  The projection of a choreography onto a role $\rr$ is
  \[
  \proj{\cret{\pp.a}}{\rr} =
    \begin{cases}
      \kw{ret}\;a & \rr = \pp\\
      \kw{ret}\;\bot & \text{otherwise}
    \end{cases}
  \qquad\qquad
  \proj{\bind{x}{A}{C}}{\rr} = \pre{A}{x}{\rr}\,\proj{C}{\rr}
  \]
  where the prefix $\pre{A}{x}{\rr}$ contributed by an action is
  \[
  \pre{\pp.e}{x}{\rr} =
    \begin{cases}
      x\leftarrow e\,; & \rr=\pp\\
      \varepsilon & \text{otherwise}
    \end{cases}
  \qquad
  \pre{\com{\pp}{a}{\qq}}{x}{\rr} =
    \begin{cases}
      \nsend{\qq}{a}\,; & \rr=\pp\\
      x\leftarrow\nrecv{\pp}\,; & \rr=\qq\\
      \varepsilon & \text{otherwise.}
    \end{cases}
  \]
  The projection $\projn{C}$ is the network mapping each $\rr\in\Roles$ to
  $\proj{C}{\rr}$.
\end{definition}

When $\rr$ is not involved in an action, the variable bound by the action
disappears from $\rr$'s program along with the action.
Lemma~\ref{lem:c1-local} shows that typing makes this safe.  The following
facts are immediate from the definition.

\begin{lemma}[Prefixes]\label{lem:c1-prefix}
  For every action $A$, variable $x$ and role $\rr$:
  \begin{enumerate}[label=(\roman*),nosep]
  \item\label{it:c1-prefix-eps} $\pre{A}{x}{\rr}=\varepsilon$ if and only if $\rr\notin\roles{A}$.
  \item\label{it:c1-prefix-bind} If $\pre{A}{x}{\rr}$ does not bind $x$, then $\rr\neq\loc{A}$.
  \end{enumerate}
\end{lemma}
\begin{proof}
  \ref{it:c1-prefix-eps} is read off the definition.  For
  \ref{it:c1-prefix-bind}, the prefixes that do not bind $x$ are $\varepsilon$,
  which occurs only when $\rr\notin\roles{A}\ni\loc{A}$, and the send at $\pp$
  in a communication, where $\pp\neq\qq=\loc{A}$.
\end{proof}

\section{Metatheory}
\label{sec:c1-meta}

\subsection{Supporting lemmas}

\begin{lemma}[Substitution]\label{lem:c1-subst}
  If $\Gamma,x{:}\pp\vdash C$ and $\Gamma\vdash_\pp a$, then $\Gamma\vdash C[a/x]$.
\end{lemma}
\begin{proof}
  By induction on $C$, renaming bound variables apart from $x$ and $\fv(a)$.
  Every expression or atom $e$ in $C$ is checked as $\Gamma,x{:}\pp\vdash_\qq e$
  for some $\qq$.  If $\qq=\pp$, then $\Gamma\vdash_\pp e[a/x]$ by the property of
  substitution on free variables.  If $\qq\neq\pp$, then $x\notin\fv(e)$ and
  $e[a/x]=e$.
\end{proof}

\begin{lemma}[Subject reduction]\label{lem:c1-sr}
  If $\Gamma\vdash C$ and $C\lto{\ell}C'$, then $\Gamma\vdash C'$.
\end{lemma}
\begin{proof}
  By induction on the step.  For \rn{Local}, $\Gamma,x{:}\pp\vdash C_0$ and
  $\Gamma\vdash_\pp v$ since $v$ is closed, so Lemma~\ref{lem:c1-subst}
  applies; \rn{Com} is the same with $\qq$ in place of $\pp$.  For \rn{Delay},
  $\Gamma,x{:}\loc{A}\vdash C_0$, and the induction hypothesis applies.
\end{proof}

\begin{lemma}[Projection is local]\label{lem:c1-local}
  If $\Gamma\vdash C$, then every free variable of $\proj{C}{\rr}$ is mapped to
  $\rr$ by $\Gamma$.
\end{lemma}
\begin{proof}
  By induction on $C$.  For $\cret{\pp.a}$ the projection is $\kw{ret}\;a$ at
  $\pp$, and $\Gamma\vdash_\pp a$; elsewhere it is closed.  For
  $\bind{x}{A}{C_0}$, the prefix $\pre{A}{x}{\rr}$ is nonempty only when
  $\rr\in\roles{A}$, and then its free variables are those of the expression
  or atom that $\rr$ evaluates or sends, which typing places at $\rr$.  By the
  induction hypothesis, every free variable of $\proj{C_0}{\rr}$ is mapped to
  $\rr$ by $\Gamma,x{:}\loc{A}$.  If the prefix binds $x$ we are done.
  Otherwise $\rr\neq\loc{A}$ by Lemma~\ref{lem:c1-prefix}\ref{it:c1-prefix-bind},
  so $x$ is not free in $\proj{C_0}{\rr}$.
\end{proof}

\begin{lemma}[Projection commutes with substitution]\label{lem:c1-psubst}
  $\proj{C[a/x]}{\rr} = \proj{C}{\rr}[a/x]$.
\end{lemma}
\begin{proof}
  By induction on $C$.  The shape of a projection depends only on the roles
  of the actions, which substitution does not change, and
  $\pre{A[a/x]}{y}{\rr} = \pre{A}{y}{\rr}[a/x]$ by inspection.  No typing is
  needed.
\end{proof}

\begin{corollary}\label{cor:c1-psubst}
  If $\Gamma,x{:}\pp\vdash C$ and $\rr\neq\pp$, then
  $\proj{C[a/x]}{\rr}=\proj{C}{\rr}$.
\end{corollary}
\begin{proof}
  By Lemma~\ref{lem:c1-local}, $x$ is not free in $\proj{C}{\rr}$; apply
  Lemma~\ref{lem:c1-psubst}.
\end{proof}

\begin{lemma}[Choreographic frame]\label{lem:c1-cframe}
  If $\Gamma\vdash C$, $C\lto{\ell}C'$ and $\rr\notin\roles{\ell}$, then
  $\proj{C'}{\rr}=\proj{C}{\rr}$.
\end{lemma}
\begin{proof}
  By induction on the step.

  \case{\rn{Local}}
  $C=\bind{x}{\pp.e}{C_0}$ and $C'=C_0[v/x]$, with $\rr\neq\pp$.  Then
  $\proj{C}{\rr}=\proj{C_0}{\rr}$, which equals $\proj{C_0[v/x]}{\rr}$ by
  Corollary~\ref{cor:c1-psubst}, since $x$ is held at $\pp\neq\rr$.

  \case{\rn{Com}}
  The same, with $\rr\notin\{\pp,\qq\}$ and $x$ held at $\qq\neq\rr$.

  \case{\rn{Delay}}
  $C=\bind{x}{A}{C_0}$ and $C'=\bind{x}{A}{C_0'}$.  The induction hypothesis
  gives $\proj{C_0'}{\rr}=\proj{C_0}{\rr}$, and both projections put the same
  prefix $\pre{A}{x}{\rr}$ in front.
\end{proof}

\begin{lemma}[Network frame]\label{lem:c1-nframe}
  Suppose $N\lto{\ell}N'$.
  \begin{enumerate}[label=(\roman*),nosep]
  \item\label{it:c1-nframe-loc} $N'(\rr)=N(\rr)$ for every $\rr\notin\roles{\ell}$.
  \item\label{it:c1-nframe-move} If $M$ agrees with $N$ on $\roles{\ell}$, then
    $M\lto{\ell}M'$, where $M'$ agrees with $N'$ on $\roles{\ell}$ and with
    $M$ elsewhere.
  \end{enumerate}
\end{lemma}
\begin{proof}
  By inspection of Figure~\ref{fig:c1-net}.  The premises of each rule
  mention only the programs of the roles in $\roles{\ell}$, and the
  conclusion updates only those.
\end{proof}

\begin{corollary}[Lifting]\label{cor:c1-lift}
  Suppose $\projn{C_0}\lto{\ell}\projn{C_0'}$.  If
  $\proj{C}{\rr}=\proj{C_0}{\rr}$ and $\proj{C'}{\rr}=\proj{C_0'}{\rr}$ for
  every $\rr\in\roles{\ell}$, and $\proj{C'}{\rr}=\proj{C}{\rr}$ for every
  $\rr\notin\roles{\ell}$, then $\projn{C}\lto{\ell}\projn{C'}$.
\end{corollary}
\begin{proof}
  By Lemma~\ref{lem:c1-nframe}\ref{it:c1-nframe-move} with $M=\projn{C}$,
  $\projn{C}\lto{\ell}M'$, where $M'$ agrees with $\projn{C_0'}$, and hence
  with $\projn{C'}$, on $\roles{\ell}$, and agrees with $\projn{C}$, and
  hence with $\projn{C'}$, elsewhere.
\end{proof}

\begin{lemma}[Network inversion]\label{lem:c1-inv}
  Suppose $N\lto{\ell}N'$ and $\rr\in\roles{\ell}$.  Then $N(\rr)$ is not a
  return, and:
  \begin{enumerate}[label=(\alph*),nosep]
  \item\label{it:c1-inv-local} if $N(\rr)=x\leftarrow e\,;\,P$, then the rule
    is \rn{N-Local}, $\ell=\tauat{\rr}$, $e\Downarrow v$ and $N'=\upd{N}{\rr}{P[v/x]}$;
  \item\label{it:c1-inv-send} if $N(\rr)=\nsend{\qq}{a}\,;\,P$, then the rule
    is \rn{N-Com}, $a$ is a value $v$, $\ell=\com{\rr}{v}{\qq}$,
    $N(\qq)=y\leftarrow\nrecv{\rr}\,;\,P'$ and
    $N'=N[\rr\mapsto P,\ \qq\mapsto P'[v/y]]$;
  \item\label{it:c1-inv-recv} if $N(\rr)=x\leftarrow\nrecv{\pp}\,;\,P$, then the
    rule is \rn{N-Com}, $N(\pp)=\nsend{\rr}{v}\,;\,P'$ for some value $v$,
    $\ell=\com{\pp}{v}{\rr}$ and $N'=N[\pp\mapsto P',\ \rr\mapsto P[v/x]]$.
  \end{enumerate}
\end{lemma}
\begin{proof}
  By inspection of Figure~\ref{fig:c1-net}.  Each role in $\roles{\ell}$
  appears in the rule's premises with a program of one of the three prefixed
  shapes.  Each shape occurs in exactly one position of exactly one rule, so
  the shape of $N(\rr)$ determines the rule and $\rr$'s part in it.
\end{proof}

\subsection{Endpoint projection is lockstep}

\begin{theorem}[EPP is lockstep]\label{thm:c1-epp}
  Suppose $\Gamma\vdash C$.
  \begin{enumerate}[label=(\arabic*),nosep]
  \item\label{it:c1-complete} \emph{Completeness.} If $C\lto{\ell}C'$, then $\projn{C}\lto{\ell}\projn{C'}$.
  \item\label{it:c1-sound} \emph{Soundness.} If $\projn{C}\lto{\ell}N$, then
    $C\lto{\ell}C'$ for some $C'$ with $N=\projn{C'}$.
  \end{enumerate}
\end{theorem}

\begin{proof}[Proof of completeness]
  By induction on the step.

  \case{\rn{Local}}
  $C=\bind{x}{\pp.e}{C_0}$, $e\Downarrow v$ and $C'=C_0[v/x]$.  At $\pp$ the
  projection is $x\leftarrow e\,;\,\proj{C_0}{\pp}$, so \rn{N-Local} gives
  $\projn{C}\lto{\tauat\pp}\upd{\projn{C}}{\pp}{\proj{C_0}{\pp}[v/x]}$.  At
  $\pp$, $\proj{C_0}{\pp}[v/x]=\proj{C_0[v/x]}{\pp}$ by
  Lemma~\ref{lem:c1-psubst}.  At $\rr\neq\pp$,
  $\proj{C}{\rr}=\proj{C_0}{\rr}=\proj{C_0[v/x]}{\rr}$ by
  Corollary~\ref{cor:c1-psubst}.

  \case{\rn{Com}}
  $C=\bind{x}{\com{\pp}{v}{\qq}}{C_0}$ and $C'=C_0[v/x]$.  The projection is
  $\nsend{\qq}{v}\,;\,\proj{C_0}{\pp}$ at $\pp$ and
  $x\leftarrow\nrecv{\pp}\,;\,\proj{C_0}{\qq}$ at $\qq$, so \rn{N-Com} fires.
  The target is $\proj{C_0}{\pp}$ at $\pp$, which equals
  $\proj{C_0[v/x]}{\pp}$ by Corollary~\ref{cor:c1-psubst} since $x$ is held
  at $\qq\neq\pp$; it is $\proj{C_0}{\qq}[v/x]=\proj{C_0[v/x]}{\qq}$ at $\qq$
  by Lemma~\ref{lem:c1-psubst}; and other roles are unchanged on both sides,
  by Corollary~\ref{cor:c1-psubst}.

  \case{\rn{Delay}}
  $C=\bind{x}{A}{C_0}$, $C'=\bind{x}{A}{C_0'}$, $C_0\lto{\ell}C_0'$ and
  $\roles{\ell}\cap\roles{A}=\emptyset$.  Since $\Gamma,x{:}\loc{A}\vdash C_0$,
  the induction hypothesis gives $\projn{C_0}\lto{\ell}\projn{C_0'}$.  We
  apply Corollary~\ref{cor:c1-lift}.  For $\rr\in\roles{\ell}$ we have
  $\rr\notin\roles{A}$, so the prefix is empty, $\proj{C}{\rr}=\proj{C_0}{\rr}$
  and $\proj{C'}{\rr}=\proj{C_0'}{\rr}$.  For $\rr\notin\roles{\ell}$,
  $\proj{C'}{\rr}=\proj{C}{\rr}$ by Lemma~\ref{lem:c1-cframe}.
\end{proof}

\begin{proof}[Proof of soundness]
  By induction on $C$.  Write $N_C$ for $\projn{C}$ and $N'$ for the target
  of the network step.  If $C$ is a return, every role's projection is a
  return, and by Lemma~\ref{lem:c1-inv} there is no step.  Otherwise
  $C=\bind{x}{A}{C_0}$, and we distinguish whether the step involves a role
  of $A$.

  \subcase{$\roles{\ell}\cap\roles{A}=\emptyset$}
  For $\rr\in\roles{\ell}$ the prefix is empty, so $N_C$ agrees with
  $\projn{C_0}$ on $\roles{\ell}$.  By
  Lemma~\ref{lem:c1-nframe}\ref{it:c1-nframe-move},
  $\projn{C_0}\lto{\ell}N''$, where $N''$ agrees with $N'$ on $\roles{\ell}$ and
  with $\projn{C_0}$ elsewhere.  The induction hypothesis gives
  $C_0\lto{\ell}C_0'$ with $\projn{C_0'}=N''$, and \rn{Delay} gives
  $C\lto{\ell}C'=\bind{x}{A}{C_0'}$.  For $\rr\in\roles{\ell}$,
  $\proj{C'}{\rr}=\proj{C_0'}{\rr}=N''(\rr)=N'(\rr)$.  For
  $\rr\notin\roles{\ell}$, $N'(\rr)=N_C(\rr)$ by
  Lemma~\ref{lem:c1-nframe}\ref{it:c1-nframe-loc}, and
  $N_C(\rr)=\pre{A}{x}{\rr}\,\proj{C_0}{\rr}=\pre{A}{x}{\rr}\,\proj{C_0'}{\rr}=\proj{C'}{\rr}$
  by Lemma~\ref{lem:c1-cframe}.

  \subcase{some $\rr\in\roles{\ell}\cap\roles{A}$}
  The program $N_C(\rr)$ begins with $\rr$'s part of $A$, and
  Lemma~\ref{lem:c1-inv} shows that the step is $A$ itself.
  \begin{itemize}
  \item If $A=\pp.e$, then $\rr=\pp$ and $N_C(\pp)=x\leftarrow e\,;\,\proj{C_0}{\pp}$.
    By Lemma~\ref{lem:c1-inv}\ref{it:c1-inv-local}, $\ell=\tauat{\pp}$,
    $e\Downarrow v$ and $N'=\upd{N_C}{\pp}{\proj{C_0}{\pp}[v/x]}$.  Rule
    \rn{Local} gives $C\lto{\tauat\pp}C_0[v/x]$ with the same $v$, and
    $\projn{C_0[v/x]}=N'$ as in the completeness proof.
  \item If $A=\com{\pp}{a}{\qq}$, then $N_C(\pp)=\nsend{\qq}{a}\,;\,\proj{C_0}{\pp}$
    and $N_C(\qq)=x\leftarrow\nrecv{\pp}\,;\,\proj{C_0}{\qq}$.  If $\rr=\pp$,
    Lemma~\ref{lem:c1-inv}\ref{it:c1-inv-send} applies; if $\rr=\qq$,
    Lemma~\ref{lem:c1-inv}\ref{it:c1-inv-recv} applies, and the sender it
    finds is $\pp$ with program $N_C(\pp)$.  Either way $a$ is a value $v$,
    $\ell=\com{\pp}{v}{\qq}$, and
    $N'=N_C[\pp\mapsto\proj{C_0}{\pp},\ \qq\mapsto\proj{C_0}{\qq}[v/x]]$.  Rule
    \rn{Com} gives $C\lto{\ell}C_0[v/x]$, and $\projn{C_0[v/x]}=N'$ as in the
    completeness proof.\qedhere
  \end{itemize}
\end{proof}

\begin{corollary}[Bisimulation]\label{cor:c1-bisim}
  The relation $\{(C,\projn{C})\mid\Gamma\vdash C\}$ is a strong bisimulation
  with identical labels between choreographies and networks.
\end{corollary}
\begin{proof}
  Theorem~\ref{thm:c1-epp}, with Lemma~\ref{lem:c1-sr} to keep the
  choreography well-typed.
\end{proof}

In particular, every network reachable from $\projn{C}$ is $\projn{C'}$ for
some $C'$ reachable from $C$.

\subsection{Deadlock freedom}
\label{sec:c1-df}

\begin{theorem}[Deadlock freedom]\label{thm:c1-df}
  Assume (L2) and $\emptyset\vdash C$.  Every network reachable from
  $\projn{C}$ is either terminated or can step.
\end{theorem}
\begin{proof}
  By Corollary~\ref{cor:c1-bisim}, every reachable network is $\projn{C'}$ for
  some $C'$ with $\emptyset\vdash C'$.  If $C'$ is a return, $\projn{C'}$ is
  terminated.  Otherwise its head action is closed: \rn{Local} fires by (L2),
  or \rn{Com} fires because a closed atom is a value.  By
  Theorem~\ref{thm:c1-epp}\ref{it:c1-complete}, $\projn{C'}$ steps.
\end{proof}

\subsection{Agreement with the sequential semantics}
\label{sec:c1-confluence}

We prove confluence on networks rather than on choreographies, and transfer it
along the bisimulation.  The network proof uses only the fact that each local
program has a single head, so it holds verbatim in every later chapter,
including those where the corresponding statement about choreographies fails
syntactically (Remark~\ref{rem:c3-diamond}).

\begin{lemma}[Network determinacy]\label{lem:c1-det}
  Assume (L1).  If $N\lto{\ell}N_1$ and $N\lto{\ell}N_2$, then $N_1=N_2$.
\end{lemma}
\begin{proof}
  The label determines the roles involved.  The heads of their programs
  determine the rule and the target, using (L1) for \rn{N-Local}.
\end{proof}

\begin{lemma}[Network diamond]\label{lem:c1-diamond}
  If $N\lto{\ell_1}N_1$, $N\lto{\ell_2}N_2$ and $\ell_1\neq\ell_2$, then
  $\roles{\ell_1}\cap\roles{\ell_2}=\emptyset$, and $N_1\lto{\ell_2}N_3$ and
  $N_2\lto{\ell_1}N_3$ for some $N_3$.
\end{lemma}
\begin{proof}
  Suppose $\rr$ is in both $\roles{\ell_1}$ and $\roles{\ell_2}$.  By
  Lemma~\ref{lem:c1-inv}, the head of $N(\rr)$ determines the label of any
  step involving $\rr$: $\tauat{\rr}$ for a local computation,
  $\com{\rr}{v}{\qq}$ for a send of $v$, and $\com{\pp}{v}{\rr}$ for a receive
  from $\pp$, where $v$ is read off $N(\pp)$.  So $\ell_1=\ell_2$, a
  contradiction.  Since $N_1$ agrees with $N$ outside $\roles{\ell_1}$
  (Lemma~\ref{lem:c1-nframe}\ref{it:c1-nframe-loc}), it agrees with $N$ on
  $\roles{\ell_2}$, so Lemma~\ref{lem:c1-nframe}\ref{it:c1-nframe-move} gives
  $N_1\lto{\ell_2}N_3$, where $N_3$ is $N_2$ on $\roles{\ell_2}$ and $N_1$
  elsewhere.  Symmetrically $N_2\lto{\ell_1}N_3'$, where $N_3'$ is $N_1$ on
  $\roles{\ell_1}$ and $N_2$ elsewhere.  Both $N_3$ and $N_3'$ are $N_1$ on
  $\roles{\ell_1}$, $N_2$ on $\roles{\ell_2}$, and $N$ elsewhere, so
  $N_3=N_3'$.
\end{proof}

For a run $\sigma$ and a role $\rr$, write $\sigma|_\rr$ for the sequence of
labels in $\sigma$ that involve $\rr$.  A run is \emph{maximal} if it is
infinite or ends in a network with no step.

\begin{lemma}[Network confluence]\label{lem:c1-conf}
  Assume (L1).  If $N$ has a maximal run of finite length $n$, then every
  maximal run from $N$ has length $n$, ends in the same network, and has the
  same $\sigma|_\rr$ for every $\rr$.
\end{lemma}
\begin{proof}
  By induction on $n$.  If $n=0$, $N$ has no step.  Otherwise let
  $N\lto{\ell_1}N_1$ begin a maximal run $\sigma_1$ of length $n$, and let
  $N\lto{\ell_2}N_2$ begin any maximal run $\sigma_2$.  If $\ell_1=\ell_2$,
  then $N_1=N_2$ by Lemma~\ref{lem:c1-det}, and the induction hypothesis for
  $N_1$ applies to both tails.  Otherwise Lemma~\ref{lem:c1-diamond} gives
  $N_3$ with $N_1\lto{\ell_2}N_3$ and $N_2\lto{\ell_1}N_3$, and disjoint
  $\roles{\ell_1}$ and $\roles{\ell_2}$.  By the induction hypothesis for
  $N_1$, every maximal run from $N_1$ has length $n-1$ and the same final
  network and per-role sequences as the tail of $\sigma_1$.  Since every run
  from $N_3$ extends to a run from $N_1$, $N_3$ has a maximal run $\tau$, and
  $\ell_2\tau$ is a maximal run from $N_1$, so $\tau$ has length $n-2$.
  Then $\ell_1\tau$ is a maximal run from $N_2$ of length $n-1$, and the
  induction hypothesis for $N_2$ shows that the tail of $\sigma_2$ has the
  same length, final network and per-role sequences as $\ell_1\tau$.  Finally,
  for every $\rr$, $(\ell_1\ell_2\tau)|_\rr=(\ell_2\ell_1\tau)|_\rr$ because at
  most one of $\ell_1$ and $\ell_2$ involves $\rr$.
\end{proof}

\begin{theorem}[Agreement with the sequential semantics]\label{thm:c1-agree}
  Assume (L1), (L2) and $\emptyset\vdash C$, and let $n$ be the number of
  actions in $C$.
  \begin{enumerate}[label=(\arabic*),nosep]
  \item Every maximal run of $\lto{}$ from $C$ has exactly $n$ steps.
  \item All maximal runs end in the same return $\cret{\pp.v}$, which is the
    result of the $\seqto$ run.
  \item For each role $\rr$, every maximal run shows the same sequence of
    labels involving $\rr$.
  \end{enumerate}
\end{theorem}
\begin{proof}
  (1) Each step removes exactly one action, and by the progress argument in
  the proof of Theorem~\ref{thm:c1-df} a run stops only at a return.  (2)
  and (3) By completeness, every maximal run of $C$ projects, state by state
  and with the same labels, to a run of $\projn{C}$, and that run is maximal
  by soundness.  These network runs are finite by (1), so
  Lemma~\ref{lem:c1-conf} shows that they all end in the same network and
  have the same per-role label sequences.  The final
  choreographies are returns, and $\projn{\cret{\pp.v}}$ determines $\pp$ (the
  only role not returning $\bot$) and $v$.  The $\seqto$ run is one of the
  maximal runs.
\end{proof}

\begin{corollary}
  Every maximal run of $\projn{C}$ ends in $\projn{\cret{\pp.v}}$, where $v$
  is the sequential result, and each role performs the same interactions, in
  the same order, as in the sequential run.
\end{corollary}

\section{Worked example}
\label{sec:c1-example}

Two independent exchanges run concurrently.  In the choreographic semantics,
$\rr$ can act first even though its actions come later in the program.
\[
\begin{array}{r@{\;\;}l}
  C \;=\; & \bind{x}{\pp.(2+3)}{}\\
          & \bind{y}{\com{\pp}{x}{\qq}}{}\\
          & \bind{z}{\rr.10}{}\\
          & \bind{w}{\com{\rr}{z}{\sr}}{}\\
          & \bind{u}{\qq.(y\cdot 2)}{}\\
          & \cret{\qq.u}
\end{array}
\]
Its projection is
\[
\begin{array}{r@{\;\;}l}
  \pp: & x\leftarrow 2+3\,;\;\nsend{\qq}{x}\,;\;\kw{ret}\;\bot\\
  \qq: & y\leftarrow\nrecv{\pp}\,;\;u\leftarrow y\cdot 2\,;\;\kw{ret}\;u\\
  \rr: & z\leftarrow 10\,;\;\nsend{\sr}{z}\,;\;\kw{ret}\;\bot\\
  \sr: & w\leftarrow\nrecv{\rr}\,;\;\kw{ret}\;\bot
\end{array}
\]
From $C$, two steps are enabled: $\tauat{\pp}$ at the head, and $\tauat{\rr}$
by \rn{Delay} past the two actions involving $\pp$ and $\qq$.  Firing
$\tauat{\rr}$ gives
\[
  C' \;=\; \bind{x}{\pp.(2+3)}{\bind{y}{\com{\pp}{x}{\qq}}{\bind{w}{\com{\rr}{10}{\sr}}{\bind{u}{\qq.(y\cdot 2)}{\cret{\qq.u}}}}}.
\]
In the network, $\rr$'s program becomes $\nsend{\sr}{10}\,;\;\kw{ret}\;\bot$,
which is $\proj{C'}{\rr}$.  Every maximal run takes five steps, and
each ends with $\qq$ holding $\kw{ret}\;10$ and every other role holding
$\kw{ret}\;\bot$.

%======================================================================
\chapter{Conditionals by Multicast}
\label{ch:cond}
%======================================================================

This chapter adds a conditional.  Instead of the selections used by the
language CC of Cruz-Filipe and Montesi~\cite{cc}, knowledge of choice is
propagated by multicast: the role $\pp$ that holds the guard sends its value
to every other role occurring in either branch, and in the same step every
participant, $\pp$ included, continues with the selected branch.  Because the
language is monadic, the conditional is an action like any other: it returns
a value, which the continuation binds.
\[
  \bind{y}{\cif{\pp}{b}{C_1}{C_2}}{C}
\]
Both branches share the continuation $C$, so roles that take part only in
$C$ are not involved in the conditional.

\section{Design choices}
\label{sec:c2-choices}

\begin{choice}[Knowledge of choice by multicast]\label{cho:c2-multicast}
  The guard holder $\pp$ sends the guard's value to the \emph{receivers}: every
  role other than $\pp$ that occurs in either branch, whether or not its
  behaviour differs between them.

  \emph{Proof saved.}  There are no selections and no merging, and
  projection is a total function.  Every participant learns the choice in the
  step that resolves it, so the projection never contains code for a branch
  that has been ruled out, and no pruning relation is needed.  The cost is
  that a role whose behaviour is the same in both branches still receives a
  message.
\end{choice}

\begin{choice}[The sender branches in the same step]\label{cho:c2-sendbranch}
  At the network level, the guard holder executes a \emph{send-and-branch}
  and each receiver a \emph{receive-and-branch}.  Rule \rn{N-Cast} is a
  single synchronous one-to-many rendezvous in which all participants select
  their branch together.

  \emph{Proof saved.}  There is no intermediate network state in which the
  guard has been announced to some participants but not to others, or in
  which the sender has not yet committed.  The choreography therefore needs
  no runtime term for a partially executed conditional, and a conditional is
  one step at both levels.
\end{choice}

\begin{choice}[The conditional binds its result]\label{cho:c2-bind}
  The conditional is an action, used as $\bind{y}{\cif{\pp}{a}{C_1}{C_2}}{C}$,
  rather than a term in tail position.  When it fires, the selected branch is
  spliced in front of the continuation with the derived bind,
  $\mbind{C_v}{y}{C}$, which is again in normal form.

  \emph{Proof saved.}  The alternative is the free monad's own normal form
  for a branching operation, which pushes the continuation into both branches.
  That would make every role of $C$ a receiver, add a message for each of
  them, and block them until the guard is known.  With a shared
  continuation, roles that take part only in $C$ are not involved in the
  conditional, and the existing \rn{Delay} rule lets them run past it
  unchanged.  The only new fact the proofs need is that projection commutes
  with bind (Lemma~\ref{lem:c2-pbind}).
\end{choice}

\begin{choice}[The guard is an atom, and the label carries its value]\label{cho:c2-guard}
  Following Choice~\ref{cho:c1-values}, the guard is an atom, not an
  expression, and the conditional fires only when the guard is a value.  The
  label $\castlbl{\pp}{v}{Q}$ records the value and the receivers.

  \emph{Proof saved.}  \rn{Cast} has no evaluation premise.  A label
  determines the step it labels, so network determinacy and the diamond
  property (\S\ref{sec:c1-confluence}) carry over with one extra case each.
\end{choice}

\begin{choice}[The result holder is a participant]\label{cho:c2-ret}
  A return $\cret{\sr.a}$ counts $\sr$ among its roles.

  \emph{Proof saved.}  The role $\sr$ that will hold the conditional's result
  occurs in both branches, so it is $\pp$ or a receiver, and it learns which
  branch was taken.  A role outside $\roles{C}$ projects to
  $\kw{ret}\;\bot$ (Lemma~\ref{lem:c2-erase}), so the single network equation
  $\mbind{(\kw{ret}\;\bot)}{x}{P}=P$ is enough.
\end{choice}

\begin{choice}[Nothing steps inside a branch]\label{cho:c2-nobranch}
  There is no rule for stepping inside a branch before the guard is
  announced.  CC has such a rule, which delays a step past a
  conditional when both branches can perform it, which is the situation
  that merging handles.

  \emph{Proof saved.}  Every receiver is waiting at its receive-and-branch
  and cannot act in a branch before the multicast, and any other role does
  not occur in either branch, so the network has no such steps to match.
  Branches therefore remain source terms until they are selected, and bind is
  only ever applied with a source term on its left.  This is used in
  Chapter~\ref{ch:fun}.
\end{choice}

\section{Syntax and typing}
\label{sec:c2-syntax}

\begin{definition}[Choreographies with conditionals]\label{def:c2-syntax}
  \[
  \begin{array}{r@{\;\;}c@{\;\;}l@{\qquad}l}
    A & ::= & \pp.e \;\mid\; \com{\pp}{a}{\qq} & \text{as in Chapter~\ref{ch:core}}\\
      & \mid & \cif{\pp}{a}{C_1}{C_2} & \text{conditional on a guard held at } \pp\\[1ex]
    C & ::= & \cret{\pp.a} \;\mid\; \bind{x}{A}{C} & \text{as in Chapter~\ref{ch:core}}
  \end{array}
  \]
  We write $C_\ttrue=C_1$ and $C_\ffalse=C_2$, so that $C_v$ is the branch
  selected by the boolean $v$.  Derived bind (Definition~\ref{def:c1-bind}) is
  unchanged; it walks only the spine and never enters a branch.
\end{definition}

\begin{definition}[Roles and receivers]\label{def:c2-roles}
  \[
  \begin{array}{r@{\;\;=\;\;}l}
    \roles{\cret{\pp.a}} & \{\pp\}\\
    \roles{\bind{x}{A}{C}} & \roles{A}\cup\roles{C}\\
    \roles{\cif{\pp}{a}{C_1}{C_2}} & \{\pp\}\cup\roles{C_1}\cup\roles{C_2}\\
    \rcv{\cif{\pp}{a}{C_1}{C_2}} & (\roles{C_1}\cup\roles{C_2})\setminus\{\pp\}
  \end{array}
  \]
  with $\roles{A}$ for local computation and communication as in
  Definition~\ref{def:c1-syntax}.
\end{definition}

The receivers may be empty, in which case the conditional is local to $\pp$.
Substitution does not change roles or receivers:
$\roles{C[a/x]}=\roles{C}$ and $\rcv{A[a/x]}=\rcv{A}$.

\paragraph{Typing.}
Since a conditional binds the result of a sub-choreography, typing must now
say where a choreography's result is held.  The judgements are
$\Gamma\vdash A:\sr$ and $\Gamma\vdash C:\sr$, meaning that the result is held
at $\sr$:
\begin{mathpar}
  \inferrule{\Gamma\vdash_\pp a}{\Gamma\vdash\cret{\pp.a}:\pp}
  \and
  \inferrule{\Gamma\vdash A:\sr' \\ \Gamma,x{:}\sr'\vdash C:\sr}{\Gamma\vdash\bind{x}{A}{C}:\sr}
  \\
  \inferrule{\Gamma\vdash_\pp e}{\Gamma\vdash\pp.e:\pp}
  \and
  \inferrule{\pp\neq\qq \\ \Gamma\vdash_\pp a}{\Gamma\vdash\com{\pp}{a}{\qq}:\qq}
  \and
  \inferrule{\Gamma\vdash_\pp a \\ \Gamma\vdash C_1:\sr \\ \Gamma\vdash C_2:\sr}
            {\Gamma\vdash\cif{\pp}{a}{C_1}{C_2}:\sr}
\end{mathpar}
The result location of a well-typed action or choreography is unique, and we
continue to write $\loc{A}$ for it.  For local computation and communication
it is as in Chapter~\ref{ch:core}.  A program is a choreography $C$ with
$\emptyset\vdash C:\sr$ for some $\sr$.

\begin{lemma}[Results are held by participants]\label{lem:c2-locroles}
  If $\Gamma\vdash C:\sr$ then $\sr\in\roles{C}$, and if $\Gamma\vdash A:\sr$
  then $\sr\in\roles{A}$.
\end{lemma}
\begin{proof}
  The first claim is by induction on $C$: a return $\cret{\sr.a}$ has role
  $\sr$, and a bind has the result location of its continuation, whose roles
  it includes.  The second follows: local computation and communication by
  definition, and a conditional because its branches have result location
  $\sr$.
\end{proof}

From Chapter~\ref{ch:fun} on, $\roles{\cdot}$ is defined on source
choreographies only, and the first claim is stated for them; the second
claim then also covers calls (\S\ref{sec:c3-syntax}).

One further assumption on the local language is used, for deadlock freedom
only.
\begin{description}[nosep,font=\normalfont\itshape]
\item[(L3) Boolean guards.] Every value that occurs as the guard of a
  conditional in a program, or that is substituted for its guard during
  execution, is $\ttrue$ or $\ffalse$.
\end{description}
Any standard type system for the local language enforces (L3).  Without it,
a conditional with a non-boolean guard is stuck in both semantics, so the
bisimulation is unaffected.

\section{Choreographic semantics}
\label{sec:c2-sem}

Labels gain a multicast:
\[
  \ell \;::=\; \tauat{\pp} \;\mid\; \com{\pp}{v}{\qq} \;\mid\; \castlbl{\pp}{v}{Q}
  \qquad
  \roles{\castlbl{\pp}{v}{Q}} = \{\pp\}\cup Q.
\]
The semantics consists of the rules of Figure~\ref{fig:c1-sem} and the rule
\begin{mathpar}
  \inferrule*[right=\rn{Cast}]
    {v\in\{\ttrue,\ffalse\} \\ A=\cif{\pp}{v}{C_\ttrue}{C_\ffalse} \\ Q=\rcv{A}}
    {\bind{x}{A}{C} \lto{\castlbl{\pp}{v}{Q}} \mbind{C_v}{x}{C}}
\end{mathpar}
The sequential semantics $\seqto$ consists of \rn{Local}, \rn{Com} and
\rn{Cast}.

\rn{Cast} announces the guard to the receivers and continues with the selected
branch followed by the shared continuation.  \rn{Delay} is unchanged: since
$\roles{A}$ for a conditional includes the roles of both branches, a step of
the continuation may be taken before the conditional exactly when it involves
none of the conditional's participants.  Nothing steps inside a branch
(Choice~\ref{cho:c2-nobranch}).

\section{Networks}
\label{sec:c2-net}

\begin{definition}[Networks with branching]\label{def:c2-net}
  Prefixes gain a send-and-branch and a receive-and-branch:
  \[
  \begin{array}{r@{\;\;}c@{\;\;}l@{\qquad}l}
    \pi & ::= & \cdots & \text{as in Definition~\ref{def:c1-net}}\\
        & \mid & x\leftarrow\sendbr{Q}{a}{P_1}{P_2}\,; & \text{multicast $a$ to $Q$ and branch on it}\\
        & \mid & x\leftarrow\recvbr{\pp}{P_1}{P_2}\,; & \text{receive a boolean from $\pp$ and branch on it}
  \end{array}
  \]
  Network bind $\mbind{P_1}{x}{P_2}$ is defined by
  \[
  \mbind{(\kw{ret}\;a)}{x}{P} = P[a/x]
  \qquad
  \mbind{(\kw{ret}\;\bot)}{x}{P} = P
  \qquad
  \mbind{(\pi\,P_1)}{x}{P_2} = \pi\,(\mbind{P_1}{x}{P_2})
  \]
  where, in the last equation, the variable bound by $\pi$ is renamed apart
  from $\fv(P_2)\cup\{x\}$.
\end{definition}

The network semantics consists of the rules of Figure~\ref{fig:c1-net} and
\begin{mathpar}
  \inferrule*[right=\rn{N-Cast}]
    {N(\pp) = x\leftarrow\sendbr{Q}{v}{P_\ttrue}{P_\ffalse}\,;\,P \\
     v\in\{\ttrue,\ffalse\} \\\\
     N(\qq) = x_\qq\leftarrow\recvbr{\pp}{P^\qq_\ttrue}{P^\qq_\ffalse}\,;\,P^\qq \text{ for every } \qq\in Q}
    {N \lto{\castlbl{\pp}{v}{Q}}
     N[\pp\mapsto\mbind{P_v}{x}{P}][\qq\mapsto\mbind{P^\qq_v}{x_\qq}{P^\qq}]_{\qq\in Q}}
\end{mathpar}
When $Q=\emptyset$ the rule is a local conditional at $\pp$, so the network
language needs no separate local branching construct.  The sender and every
receiver bind the branch's result to their own copy of the variable, and a
participant whose branch returns $\bot$ goes straight on to its
continuation.

\section{Endpoint projection}
\label{sec:c2-epp}

\begin{definition}[Projection of a conditional]\label{def:c2-epp}
  Projection is as in Definition~\ref{def:c1-epp}, with the prefix for a
  conditional $A=\cif{\pp}{a}{C_1}{C_2}$ with receivers $Q=\rcv{A}$ given by
  \[
  \pre{A}{x}{\rr} =
    \begin{cases}
      x\leftarrow\sendbr{Q}{a}{\proj{C_1}{\pp}}{\proj{C_2}{\pp}}\,; & \rr=\pp\\
      x\leftarrow\recvbr{\pp}{\proj{C_1}{\rr}}{\proj{C_2}{\rr}}\,; & \rr\in Q\\
      \varepsilon & \text{otherwise.}
    \end{cases}
  \]
\end{definition}

A role that is neither $\pp$ nor a receiver does not occur in either branch,
and the conditional is erased from its program, as any other action would be
(Choice~\ref{cho:c1-erase}).  This differs from tail-position conditionals,
where such a role would be given a terminated program.  A receiver that does
not occur in one of the branches gets $\kw{ret}\;\bot$ there
(Lemma~\ref{lem:c2-erase}), and after the multicast it continues with its
own projection of the continuation.

Lemma~\ref{lem:c1-prefix} still holds.  Part \ref{it:c1-prefix-eps} is
immediate, and for part \ref{it:c1-prefix-bind} the conditional's prefixes
bind $x$, while its empty prefix occurs only outside $\roles{A}$, which
contains $\loc{A}$ by Lemma~\ref{lem:c2-locroles}.

\section{Correctness}
\label{sec:c2-correct}

All lemmas of \S\ref{sec:c1-meta} hold for the extended language, with
$\Gamma\vdash C$ replaced by $\Gamma\vdash C:\sr$ and $\loc{A}$ given by
typing.  Their proofs are by the same inductions, and we give only the new
cases, after three new lemmas.

\begin{lemma}[Erasure]\label{lem:c2-erase}
  If $\rr\notin\roles{C}$, then $\proj{C}{\rr}=\kw{ret}\;\bot$.
\end{lemma}
\begin{proof}
  By induction on $C$.  For $\cret{\pp.a}$, $\rr\neq\pp$.  For
  $\bind{x}{A}{C_0}$, the prefix is empty by
  Lemma~\ref{lem:c1-prefix}\ref{it:c1-prefix-eps}, and the induction
  hypothesis applies to $C_0$.
\end{proof}

\begin{lemma}[Typing of bind]\label{lem:c2-tbind}
  If $\Gamma\vdash C_1:\sr'$ and $\Gamma,x{:}\sr'\vdash C_2:\sr$, then
  $\Gamma\vdash\mbind{C_1}{x}{C_2}:\sr$.
\end{lemma}
\begin{proof}
  By induction on $C_1$.  The return case is Lemma~\ref{lem:c1-subst}, and the
  bind case follows from the induction hypothesis after weakening $C_2$ with
  the fresh bound variable.
\end{proof}

\begin{lemma}[Projection commutes with bind]\label{lem:c2-pbind}
  If $\Gamma\vdash C_1:\sr'$ and $\Gamma,x{:}\sr'\vdash C_2:\sr$, then for
  every role $\rr$,
  \[
    \proj{\mbind{C_1}{x}{C_2}}{\rr} = \mbind{\proj{C_1}{\rr}}{x}{\proj{C_2}{\rr}}.
  \]
\end{lemma}
\begin{proof}
  By induction on $C_1$.

  \case{$C_1=\cret{\sr'.a}$}
  The left side is $\proj{C_2[a/x]}{\rr}=\proj{C_2}{\rr}[a/x]$, by
  Lemma~\ref{lem:c1-psubst}.  If $\rr=\sr'$, the right side is
  $\mbind{(\kw{ret}\;a)}{x}{\proj{C_2}{\rr}}=\proj{C_2}{\rr}[a/x]$.  If
  $\rr\neq\sr'$, the right side is
  $\mbind{(\kw{ret}\;\bot)}{x}{\proj{C_2}{\rr}}=\proj{C_2}{\rr}$, and $x$ is not
  free in $\proj{C_2}{\rr}$ by Lemma~\ref{lem:c1-local}, since $x$ is held at
  $\sr'\neq\rr$.  So both sides are $\proj{C_2}{\rr}$.

  \case{$C_1=\bind{y}{A}{C_1'}$}
  The left side is
  $\pre{A}{y}{\rr}\,\proj{\mbind{C_1'}{x}{C_2}}{\rr}$, which by the induction
  hypothesis is
  $\pre{A}{y}{\rr}\,(\mbind{\proj{C_1'}{\rr}}{x}{\proj{C_2}{\rr}})$.  By the
  definition of network bind, this is
  $\mbind{(\pre{A}{y}{\rr}\,\proj{C_1'}{\rr})}{x}{\proj{C_2}{\rr}}$, the right
  side.  The bound variable $y$ is fresh for $C_2$ on both sides.
\end{proof}

The ret case is the only one that needs typing.  A role that does not hold
the result continues directly, by the equation
$\mbind{(\kw{ret}\;\bot)}{x}{P}=P$ of Choice~\ref{cho:c2-ret}, because typing
guarantees that it never mentions the result.

\paragraph{New cases of the supporting lemmas.}
\begin{itemize}
\item \emph{Substitution} (Lemma~\ref{lem:c1-subst}): the guard is checked at
  $\pp$ like any atom, and the branches by the induction hypothesis.
\item \emph{Subject reduction} (Lemma~\ref{lem:c1-sr}), case \rn{Cast}: from
  $\Gamma\vdash C_v:\sr'$ and $\Gamma,x{:}\sr'\vdash C:\sr$,
  Lemma~\ref{lem:c2-tbind} gives $\Gamma\vdash\mbind{C_v}{x}{C}:\sr$.
\item \emph{Projection is local} (Lemma~\ref{lem:c1-local}): at $\pp$ the
  prefix mentions the guard, held at $\pp$, and $\proj{C_i}{\pp}$, whose free
  variables are held at $\pp$ by the induction hypothesis; at a receiver
  $\rr$ it mentions only $\proj{C_i}{\rr}$.  The bound variable is handled by
  Lemma~\ref{lem:c1-prefix}\ref{it:c1-prefix-bind} as before.
\item \emph{Projection commutes with substitution} (Lemma~\ref{lem:c1-psubst}):
  substitution does not change $\rcv{A}$, so the prefix of $A[a/x]$ is the
  prefix of $A$ with $a$ substituted.
\item \emph{Choreographic frame} (Lemma~\ref{lem:c1-cframe}), case \rn{Cast}:
  let $C=\bind{x}{A}{C_0}$ and $\rr\notin\{\pp\}\cup Q=\roles{A}$.  Then
  $\proj{C}{\rr}=\proj{C_0}{\rr}$, and by Lemma~\ref{lem:c2-pbind} and
  Lemma~\ref{lem:c2-erase}, since $\rr\notin\roles{C_v}\subseteq\roles{A}$,
  \[
    \proj{\mbind{C_v}{x}{C_0}}{\rr}
    = \mbind{\proj{C_v}{\rr}}{x}{\proj{C_0}{\rr}}
    = \mbind{(\kw{ret}\;\bot)}{x}{\proj{C_0}{\rr}}
    = \proj{C_0}{\rr}.
  \]
\item \emph{Network frame} (Lemma~\ref{lem:c1-nframe}): in \rn{N-Cast} the
  roles whose programs are read or changed are exactly $\{\pp\}\cup Q$.
\item \emph{Network inversion} (Lemma~\ref{lem:c1-inv}) gains two items:
  \begin{enumerate}[label=(\alph*),start=4,nosep]
  \item\label{it:c2-inv-cast} if $N(\rr)=x\leftarrow\sendbr{Q}{a}{P_\ttrue}{P_\ffalse}\,;\,P$,
    then the rule is \rn{N-Cast}, $a$ is a value $v\in\{\ttrue,\ffalse\}$,
    $\ell=\castlbl{\rr}{v}{Q}$, every $\qq\in Q$ has a program
    $N(\qq)=x_\qq\leftarrow\recvbr{\rr}{P^\qq_\ttrue}{P^\qq_\ffalse}\,;\,P^\qq$,
    and $N'$ is as in the conclusion of \rn{N-Cast};
  \item\label{it:c2-inv-brecv} if $N(\rr)=x\leftarrow\recvbr{\pp}{P_\ttrue}{P_\ffalse}\,;\,P$,
    then the rule is \rn{N-Cast}, $N(\pp)=y\leftarrow\sendbr{Q}{v}{P'_\ttrue}{P'_\ffalse}\,;\,P'$
    for some $Q\ni\rr$ and $v\in\{\ttrue,\ffalse\}$, and item
    \ref{it:c2-inv-cast} applies to $\pp$.
  \end{enumerate}
  Each of the two new shapes occurs in exactly one position of one rule.
\end{itemize}

\begin{theorem}[EPP is lockstep]\label{thm:c2-epp}
  Suppose $\Gamma\vdash C:\sr$.
  \begin{enumerate}[label=(\arabic*),nosep]
  \item\label{it:c2-complete} If $C\lto{\ell}C'$, then $\projn{C}\lto{\ell}\projn{C'}$.
  \item\label{it:c2-sound} If $\projn{C}\lto{\ell}N$, then $C\lto{\ell}C'$
    for some $C'$ with $N=\projn{C'}$.
  \end{enumerate}
\end{theorem}
\begin{proof}
  As for Theorem~\ref{thm:c1-epp}, with one new case in each direction.

  \case{\rn{Cast}, completeness}
  $C=\bind{x}{A}{C_0}$ with $A=\cif{\pp}{v}{C_\ttrue}{C_\ffalse}$, $Q=\rcv{A}$,
  $\ell=\castlbl{\pp}{v}{Q}$ and $C'=\mbind{C_v}{x}{C_0}$.  By
  Definition~\ref{def:c2-epp}, $\proj{C}{\pp}$ is a send-and-branch to $Q$
  with guard $v$ and branches $\proj{C_\ttrue}{\pp}$ and $\proj{C_\ffalse}{\pp}$,
  and for every $\qq\in Q$, $\proj{C}{\qq}$ is a receive-and-branch from $\pp$
  with branches $\proj{C_\ttrue}{\qq}$ and $\proj{C_\ffalse}{\qq}$, each
  followed by the projection of $C_0$.  Rule \rn{N-Cast} fires.  Its target
  at each $\rr\in\{\pp\}\cup Q$ is
  $\mbind{\proj{C_v}{\rr}}{x}{\proj{C_0}{\rr}}=\proj{C'}{\rr}$, by
  Lemma~\ref{lem:c2-pbind}.  At every other role the target is unchanged,
  and so is the projection, by the \rn{Cast} case of
  Lemma~\ref{lem:c1-cframe}.

  \case{Conditional, soundness}
  $C=\bind{x}{A}{C_0}$ with $A=\cif{\pp}{a}{C_\ttrue}{C_\ffalse}$, and the
  step involves some $\rr\in\roles{A}=\{\pp\}\cup Q$.  If $\rr=\pp$, then
  $\projn{C}(\pp)$ is a send-and-branch to $Q$ with guard $a$, and
  Lemma~\ref{lem:c1-inv}\ref{it:c2-inv-cast} (as extended above) gives that $a$ is a value
  $v\in\{\ttrue,\ffalse\}$, $\ell=\castlbl{\pp}{v}{Q}$, and that the target is
  that of \rn{N-Cast}.  If $\rr\in Q$, then $\projn{C}(\rr)$ is a
  receive-and-branch from $\pp$, and
  Lemma~\ref{lem:c1-inv}\ref{it:c2-inv-brecv} (as extended above) reduces to the previous case,
  since the sender it names is $\pp$ and $\projn{C}(\pp)$ is a send-and-branch
  to $Q$.  Rule \rn{Cast} gives $C\lto{\castlbl{\pp}{v}{Q}}\mbind{C_v}{x}{C_0}$,
  and the projection of the target is the network target, as in the
  completeness case.

  The disjoint subcase of soundness, for any action, and the \rn{Delay} case
  of completeness are unchanged: they use only the empty prefix at roles
  outside $\roles{A}$ and the frame lemmas.
\end{proof}

\begin{corollary}[Bisimulation]\label{cor:c2-bisim}
  $\{(C,\projn{C})\mid\Gamma\vdash C:\sr\}$ is a strong bisimulation with
  identical labels.
\end{corollary}

\begin{theorem}[Deadlock freedom]\label{thm:c2-df}
  Assume (L2), (L3) and $\emptyset\vdash C:\sr$.  Every network reachable from
  $\projn{C}$ is either terminated or can step.
\end{theorem}
\begin{proof}
  As for Theorem~\ref{thm:c1-df}.  The new case is a head conditional, whose
  guard is closed and hence a value, which is a boolean by (L3), so \rn{Cast}
  fires.
\end{proof}

\paragraph{Confluence.}
Lemmas~\ref{lem:c1-det}, \ref{lem:c1-diamond} and~\ref{lem:c1-conf} hold for
networks with \rn{N-Cast}.  The label $\castlbl{\pp}{v}{Q}$ determines the
rule, the roles and the value, and the head of $N(\rr)$ still determines the
label of any step involving $\rr$: a send-and-branch at $\rr$ forces
$\castlbl{\rr}{v}{Q}$ with $v$ its guard, and a receive-and-branch from $\pp$
forces the label read off $N(\pp)$.  Theorem~\ref{thm:c1-agree} therefore
holds, assuming (L3) in addition so that a run can only stop at a return, with
one change: since the two branches of a conditional need not have
the same number of actions, statement~(1) becomes ``every maximal run is
finite, and all maximal runs have the same length''.  Runs are finite
because every step strictly decreases the total number of actions, counting
those inside branches.

\section{Worked example}
\label{sec:c2-example}

Role $\pp$ decides whether to forward a value to $\qq$, while $\rr$ and $\sr$
exchange a value independently.
\[
\begin{array}{r@{\;\;}l}
  C \;=\; & \bind{n}{\pp.7}{}\\
          & \bind{b}{\pp.(n>5)}{}\\
          & \bind{y}{\cif{\pp}{b}{(\bind{z}{\com{\pp}{n}{\qq}}{\cret{\qq.z}})}{(\cret{\qq.0})}}{}\\
          & \bind{w}{\rr.1}{}\\
          & \bind{u}{\com{\rr}{w}{\sr}}{}\\
          & \cret{\qq.y}
\end{array}
\]
The conditional has roles $\{\pp,\qq\}$ and receivers $\{\qq\}$.  Role $\qq$
takes part in the communication of the first branch, and it holds the result
of both, so it learns the outcome even when the second branch, where it does
nothing else, is taken (Choice~\ref{cho:c2-ret}).  Its projection is
\[
\begin{array}{r@{\;\;}l}
  \pp: & n\leftarrow 7\,;\; b\leftarrow n>5\,;\;
         y\leftarrow\sendbr{\{\qq\}}{b}{(\nsend{\qq}{n}\,;\,\kw{ret}\;\bot)}{\kw{ret}\;\bot}\,;\;\kw{ret}\;\bot\\
  \qq: & y\leftarrow\recvbr{\pp}{(z\leftarrow\nrecv{\pp}\,;\,\kw{ret}\;z)}{\kw{ret}\;0}\,;\;\kw{ret}\;y\\
  \rr: & w\leftarrow 1\,;\;\nsend{\sr}{w}\,;\;\kw{ret}\;\bot\\
  \sr: & u\leftarrow\nrecv{\rr}\,;\;\kw{ret}\;\bot
\end{array}
\]
Neither $\rr$ nor $\sr$ is involved in the conditional.  They receive no
message, and \rn{Delay} lets them run before, during or after it.  Had the
continuation been copied into both branches, both would have been receivers.

After $\pp$'s two local steps, the guard is $\ttrue$ and \rn{Cast} fires with
label $\castlbl{\pp}{\ttrue}{\{\qq\}}$.  Suppose $\rr$ and $\sr$ have not yet
acted.  The choreography becomes
\begin{align*}
  &\mbind{(\bind{z}{\com{\pp}{7}{\qq}}{\cret{\qq.z}})}{y}{(\bind{w}{\rr.1}{\bind{u}{\com{\rr}{w}{\sr}}{\cret{\qq.y}}})}\\
  &\qquad=\;
  \bind{z}{\com{\pp}{7}{\qq}}{\bind{w}{\rr.1}{\bind{u}{\com{\rr}{w}{\sr}}{\cret{\qq.z}}}}.
\end{align*}
In the network, $\pp$ continues as
$\mbind{(\nsend{\qq}{7}\,;\,\kw{ret}\;\bot)}{y}{\kw{ret}\;\bot}=\nsend{\qq}{7}\,;\,\kw{ret}\;\bot$
and $\qq$ as
$\mbind{(z\leftarrow\nrecv{\pp}\,;\,\kw{ret}\;z)}{y}{\kw{ret}\;y}=z\leftarrow\nrecv{\pp}\,;\,\kw{ret}\;z$,
which are the projections of the new choreography.  Every maximal run has six
steps and ends in $\cret{\qq.7}$.

%======================================================================
\chapter{First-Order Functions and Recursion}
\label{ch:fun}
%======================================================================

This chapter adds named functions with located parameters, recursion, and
calls that return a value:
\[
  \bind{x}{f(\pp_1.a_1,\ldots,\pp_n.a_n)}{C}.
\]
Unlike in CC, calls need not be in tail position: the result of a call is
held at a role and bound by the continuation, like the result of any other
action.

At the network level each participant of a function runs its own local
version of it and enters a call independently of the others.  We follow the
approach~\cite{montesi-book,cmp-jar23,ozone24} of making every such entry a
visible step at both levels and
recording, in the choreography, a call that some participants have entered
and others have not, with a runtime term.  In the monadic setting the
runtime term must also absorb the continuation of the call
(Choice~\ref{cho:c3-absorb}).

\section{Design choices}
\label{sec:c3-choices}

\begin{choice}[Entering a call is a visible step for each role]\label{cho:c3-enter}
  A role enters a call on its own, labelled $\tauat{\rr}$, at both levels.  In
  the choreography, the first participant to enter unfolds the body (rule
  \rn{Call}), and the remaining participants enter one at a time (rule
  \rn{Enter}).

  \emph{Proof saved.}  If the choreography unfolded a call for all
  participants at once, as when calls are unfolded by a structural
  congruence, one choreography step would correspond to several network
  steps, or to none, and the bisimulation could only be stated up to
  unfolding.  Making all participants synchronise at a call would restore
  lockstep, but only by adding synchronisation that no implementation needs.
  With per-role entry, every unfolding is a labelled step at both levels and
  the correspondence stays exact.
\end{choice}

\begin{choice}[The runtime term absorbs the continuation]\label{cho:c3-absorb}
  When the first role enters $K=\bind{x}{f(\vec\alpha)}{C}$, the choreography
  becomes $\rt{S}{K}{D}$, where $S$ is the set of participants still to
  enter, $K$ is kept as it was, and $D$ starts as
  $\mbind{C_f[\vec a/\vec x]}{x}{C}$: the body followed by the continuation.
  The runtime term is a choreography in tail position, not an action.

  \emph{Proof saved.}  Recording the call as an action,
  $\bind{x}{\rt{S}{f(\vec\alpha)}{C'}}{C}$, with the continuation outside,
  breaks the bisimulation (Remark~\ref{rem:c3-action}): a role that holds
  the result can continue in the network as soon as its part of the body is
  done, while in the choreography $x$ is not substituted until every
  participant has entered.  Absorbing the continuation applies monad
  associativity once, at entry, so the result reaches the continuation
  exactly when the body produces it, as in the network.
\end{choice}

\begin{choice}[Role sets are declared, not computed]\label{cho:c3-rho}
  Each function declares a set $\rho_f$ of participants, which is checked to
  contain $\roles{C_f}$ (counting the declared sets of the functions it
  calls).

  \emph{Proof saved.}  The proofs use only the inclusion
  $\roles{C_f}\subseteq\rho_f$, never its minimality, so checking replaces
  computing a fixed point, and $\roles{\cdot}$ remains a simple structural
  function.  The least valid declaration exists and can be inferred
  (Remark~\ref{rem:c3-lfp}).  A role in $\rho_f$ that has nothing to do in the
  body still enters the call: one $\tauat{\rr}$ step at both levels, and
  nothing else.  Chapter~\ref{ch:ho} relies on the same device.
\end{choice}

\begin{choice}[Calls are erased at non-participants]\label{cho:c3-erase}
  A role outside $\rho_f$ gets nothing for a call, as for any action
  (Choice~\ref{cho:c1-erase}).

  \emph{Proof saved.}  Non-participants take no steps for the call, and
  \rn{Delay} lets them run past it.  This is also why the network can run a
  role's continuation as soon as that role is done with the body.
\end{choice}

\begin{choice}[Nothing in a call needs to be a value]\label{cho:c3-novalue}
  The label of an entry, $\tauat{\rr}$, carries no value, so following
  Choice~\ref{cho:c1-values}, neither \rn{Call}, \rn{Enter} nor \rn{N-Call}
  requires the arguments to be values.  Arguments are substituted into the
  body at the first entry, as atoms.  An argument of a role that has not yet
  entered may still be a variable bound by a pending action, and the
  substitution of its value later reaches into the runtime term.

  \emph{Proof saved.}  There is no premise to match between the two
  semantics.  For closed well-typed terms a value premise would be
  redundant: an argument held at $\rr$ that is still a variable is bound by a
  pending action involving $\rr$, so \rn{Delay} already prevents $\rr$ from
  entering.  For open terms it would break soundness, since \rn{N-Call}
  fires whatever the arguments are, and so would requiring values in
  Chapter~\ref{ch:ho}, where some arguments are functions that may contain
  free variables.
\end{choice}

\begin{choice}[One well-formedness invariant]\label{cho:c3-wf}
  A runtime term $\rt{S}{K}{D}$ is well-formed if the roles still to enter
  participate in the call, and the current state $D$ still expects of each
  of them exactly what they will do when they enter.

  \emph{Proof saved.}  This is the only invariant in the development.  It
  holds trivially for source choreographies, and \rn{Call} establishes it.
  \rn{Local}, \rn{Com} and \rn{Cast} preserve it because it is closed under
  substitution and bind, and \rn{Inside} by the choreographic frame lemma.
\end{choice}

\begin{choice}[\rn{Call} is \rn{Enter} on a virtual runtime term]\label{cho:c3-view}
  Before anyone has entered, the call $K$ has the same projection as the
  runtime term $\rt{\rho_f}{K}{\unfold{K}}$, in which every participant is
  still to enter (Lemma~\ref{lem:c3-view}).  \rn{Call} is then exactly
  \rn{Enter} applied to this term.

  \emph{Proof saved.}  Every proof case for \rn{Call} is an instance of the
  corresponding case for \rn{Enter}.
\end{choice}

\section{Syntax and typing}
\label{sec:c3-syntax}

We fix a finite set of function names $f,g$.  A \emph{signature} $\Sig$
assigns each $f$ a type $\farr{\pp_1,\ldots,\pp_n}{\rho_f}{\sr_f}$: the roles
holding its arguments, its declared set of participants, and the role holding
its result.  A set of \emph{definitions} $\Defs$ assigns each $f$ a body
$f(x_1,\ldots,x_n)=C_f$.

\begin{definition}[Choreographies with calls]\label{def:c3-syntax}
  \[
  \begin{array}{r@{\;\;}c@{\;\;}l@{\qquad}l}
    A & ::= & \cdots & \text{as in Definition~\ref{def:c2-syntax}}\\
      & \mid & f(\pp_1.a_1,\ldots,\pp_n.a_n) & \text{call, written } f(\vec\alpha)\\[1ex]
    C & ::= & \cret{\pp.a} \;\mid\; \bind{x}{A}{C} & \text{as before}\\
      & \mid & \rt{S}{K}{D} & \text{runtime term, } \emptyset\neq S\subseteq\Roles\\[1ex]
    K & ::= & \bind{x}{f(\vec\alpha)}{C} & \text{a call with its continuation}
  \end{array}
  \]
  A choreography is a \emph{source} choreography if it contains no runtime
  term.  The branches of a conditional and the bodies $C_f$ are source
  choreographies.  We write $\vec a$ for the atoms $a_1,\ldots,a_n$ of the
  arguments $\vec\alpha$, and $\vec\alpha|_\rr$ for the subsequence of those
  $a_i$ with $\pp_i=\rr$; likewise $\vec x|_\rr$ for the parameters $x_i$ with
  $\pp_i=\rr$.
\end{definition}

The runtime term $\rt{S}{K}{D}$ is not written by programmers.  It records a
call $K$ that some, but not all, of its participants have entered: the roles
in $S$ have yet to enter, and $D$ is the current state of the body followed by
the continuation.  So that the rules can be stated uniformly, the notation
$\rt{\emptyset}{K}{D}$ stands for $D$: once every participant has entered,
there is nothing left to record.  Substitution acts on both components,
$(\rt{S}{K}{D})[a/y]=\rt{S}{K[a/y]}{D[a/y]}$.

\paragraph{Roles.}
$\roles{f(\vec\alpha)}=\rho_f$, and $\roles{\cdot}$ is otherwise as in
Definition~\ref{def:c2-roles}.  It is only ever applied to source terms:
actions, branches and bodies.

\paragraph{Typing.}
Two rules are added to those of \S\ref{sec:c2-syntax}:
\begin{mathpar}
  \inferrule{\Sig(f)=\farr{\pp_1,\ldots,\pp_n}{\rho_f}{\sr_f} \\ \Gamma\vdash_{\pp_i}a_i\ \text{for each } i}
            {\Gamma\vdash f(\pp_1.a_1,\ldots,\pp_n.a_n):\sr_f}
  \and
  \inferrule{\Gamma\vdash K:\sr \\ \Gamma\vdash D:\sr}{\Gamma\vdash\rt{S}{K}{D}:\sr}
\end{mathpar}
The definitions are \emph{well-typed} if, for every $f$ with
$\Sig(f)=\farr{\pp_1,\ldots,\pp_n}{\rho_f}{\sr_f}$,
\[
  x_1{:}\pp_1,\ldots,x_n{:}\pp_n\vdash C_f:\sr_f
  \qquad\text{and}\qquad
  \roles{C_f}\subseteq\rho_f.
\]
We assume this throughout.  By the first claim of
Lemma~\ref{lem:c2-locroles}, applied to the source choreography $C_f$,
$\sr_f\in\roles{C_f}\subseteq\rho_f$, so every function has at least one
participant, and the second claim holds for calls.

\begin{remark}[Inferring role sets]\label{rem:c3-lfp}
  Since $\roles{C_f}$ depends on the declared sets of the functions $C_f$
  calls, the valid declarations are the functions $\rho$ with
  $\roles{C_f}\subseteq\rho_f$ for all $f$.  The map
  $\rho\mapsto(f\mapsto\roles{C_f})$ is monotone on a finite lattice, so the
  least valid declaration is its least fixed point, computed by iteration from
  the everywhere-empty function.  Using it minimises the receivers of every
  conditional.
\end{remark}

\section{Choreographic semantics}
\label{sec:c3-sem}

\begin{definition}[Unfolding]\label{def:c3-unfold}
  For $K=\bind{x}{f(\vec\alpha)}{C}$, with the bound variables of $C_f$
  renamed apart,
  \[
    \unfold{K} \;=\; \mbind{C_f[\vec a/\vec x]}{x}{C}.
  \]
\end{definition}

\begin{figure}[t]
\begin{mathpar}
  \inferrule*[right=\rn{Call}]
    {K=\bind{x}{f(\vec\alpha)}{C} \\ \rr\in\rho_f}
    {K \lto{\tauat{\rr}} \rt{\rho_f\setminus\{\rr\}}{K}{\unfold{K}}}
  \and
  \inferrule*[right=\rn{Enter}]
    {\rr\in S}
    {\rt{S}{K}{D} \lto{\tauat{\rr}} \rt{S\setminus\{\rr\}}{K}{D}}
  \and
  \inferrule*[right=\rn{Inside}]
    {D\lto{\ell}D' \\ S\cap\roles{\ell}=\emptyset}
    {\rt{S}{K}{D} \lto{\ell} \rt{S}{K}{D'}}
\end{mathpar}
\caption{Rules for calls.  A target $\rt{\emptyset}{K}{D}$ stands for $D$.}
\label{fig:c3-sem}
\end{figure}

The semantics consists of \rn{Local}, \rn{Com}, \rn{Delay}, \rn{Cast}, and
the rules of Figure~\ref{fig:c3-sem}.  Rule \rn{Call} lets a participant
$\rr$ of $f$ enter: the body, followed by the continuation, becomes active,
and every other participant still has to enter.  Rule \rn{Enter} lets a
further participant enter; when the last one does, the runtime term
disappears.  Rule \rn{Inside} lets the roles that have entered proceed, as
long as the step involves no role that has not.  It plays the role of
\rn{Delay} for runtime terms, with $S$ in place of an action.  Since $D$
contains the continuation, \rn{Inside} also lets non-participants of $f$ run
past the call, and lets participants that have finished their part of the
body continue with the result.

\rn{Call} fires on a call that is not at the head through \rn{Delay}, like any
other action.  The term $K$ in $\rt{S}{K}{D}$ never steps.  It is only used
to project the roles in $S$, which have not moved since the first entry.

\begin{remark}[Why the runtime term absorbs the continuation]\label{rem:c3-action}
  Let $\Sig(f)=\farr{}{\{\qq,\rr\}}{\rr}$ and
  $f()=\bind{y}{\qq.e}{\cret{\rr.5}}$, and consider
  \[
    \bind{x}{f()}{\bind{z}{\rr.(x+1)}{\cret{\rr.z}}}.
  \]
  In the network, $\rr$'s program is
  $x\leftarrow f()\,;\;z\leftarrow x+1\,;\;\kw{ret}\;z$ and $\rr$'s local
  version of $f$ is $\kw{ret}\;5$.  After $\rr$ enters, its program is
  $z\leftarrow 5+1\,;\;\kw{ret}\;z$, and it can compute $5+1$ while $\qq$ has
  not yet entered.  If the runtime term were an action, the choreography would
  be $\bind{x}{\rt{\{\qq\}}{f()}{(\bind{y}{\qq.e}{\cret{\rr.5}})}}{\bind{z}{\rr.(x+1)}{\cret{\rr.z}}}$,
  in which $x$ is still bound and $\rr.(x+1)$ cannot fire.  With the runtime
  term of Figure~\ref{fig:c3-sem}, the choreography after $\rr$ enters is
  \[
    \rt{\{\qq\}}{K}{(\bind{y}{\qq.e}{\bind{z}{\rr.(5+1)}{\cret{\rr.z}}})},
  \]
  where $K$ is the program above,
  and $\rr$'s step is available through \rn{Inside} and \rn{Delay}.
\end{remark}

\begin{remark}[Nesting]\label{rem:c3-nesting}
  A role may fall arbitrarily far behind the others, and runtime terms may
  then nest.  With $\Sig(g)=\farr{}{\{\pp,\qq,\rr\}}{\pp}$ and
  $g()=\bind{x}{\com{\pp}{1}{\qq}}{\bind{y}{\rr.1}{\bind{z}{g()}{\cret{\pp.z}}}}$,
  roles $\pp$ and $\qq$ can iterate while $\rr$ has not entered the first
  call, building a term of the form
  $\rt{\{\rr\}}{K_1}{(\bind{y}{\rr.1}{\rt{\{\rr\}}{K_2}{(\cdots)}})}$ that
  records the actions $\rr$ still owes.  The network state stays small:
  $\rr$ is still at its first call.
\end{remark}

\section{Networks}
\label{sec:c3-net}

\begin{definition}[Networks with functions]\label{def:c3-net}
  Prefixes gain a local call, $\pi ::= \cdots \mid x\leftarrow f(\vec a)\,;$.
  The network semantics is parameterised by \emph{local definitions}
  $\Defs_\rr(f)=(\vec y).\,P$, giving, for each role $\rr$ and function $f$ that
  $\rr$ participates in, a parameter list and a body.  It consists of the rules
  of \S\ref{sec:c2-net} and
  \begin{mathpar}
    \inferrule*[right=\rn{N-Call}]
      {N(\rr)=x\leftarrow f(\vec a)\,;\,P \\ \Defs_\rr(f)=(\vec y).\,P_f}
      {N \lto{\tauat{\rr}} \upd{N}{\rr}{\mbind{P_f[\vec a/\vec y]}{x}{P}}}
  \end{mathpar}
\end{definition}

A role unfolds its own call without synchronising with anyone.  No network
counterpart of the runtime term is needed: a call that some participants have
entered is a network in which some roles are at the call and the
others are running their local bodies.

\section{Endpoint projection}
\label{sec:c3-epp}

\begin{definition}[Projection of calls and runtime terms]\label{def:c3-epp}
  Projection extends Definition~\ref{def:c2-epp} with
  \[
  \pre{f(\vec\alpha)}{x}{\rr} =
    \begin{cases}
      x\leftarrow f(\vec\alpha|_\rr)\,; & \rr\in\rho_f\\
      \varepsilon & \text{otherwise}
    \end{cases}
  \qquad
  \proj{\rt{S}{K}{D}}{\rr} =
    \begin{cases}
      \proj{K}{\rr} & \rr\in S\\
      \proj{D}{\rr} & \text{otherwise.}
    \end{cases}
  \]
  The projected local definitions are
  $\projn{\Defs}_\rr(f)=(\vec x|_\rr).\,\proj{C_f}{\rr}$ for $\rr\in\rho_f$.
  Throughout, the network semantics is taken with respect to
  $\projn{\Defs}$.
\end{definition}

A role that has not entered a call is, in the projection, still at the call.
Every other role runs its projection of the current state.
Lemma~\ref{lem:c1-prefix} still holds: a call's prefix binds $x$ at every
participant, and is empty only outside $\rho_f$, which contains $\sr_f$.

\section{Well-formedness}
\label{sec:c3-wf}

A role $\qq\in S$ that has not yet entered $\rt{S}{K}{D}$ is, in the
projection, at the call $K$.  When it enters it will continue as its
projection of $\unfold{K}$.  For the choreography and the network to stay in
step, $D$ must still expect exactly that behaviour from $\qq$.

\begin{definition}[Well-formedness]\label{def:c3-wf}
  $\WF(C)$ is defined by recursion on $C$: $\WF(\cret{\pp.a})$ holds;
  $\WF(\bind{x}{A}{C})$ holds if $\WF(C)$ does; and
  $\WF(\rt{S}{K}{D})$ holds, where $K=\bind{x}{f(\vec\alpha)}{C}$, if
  \begin{enumerate}[label=(\roman*),nosep]
  \item\label{it:c3-wf-sub} $\WF(K)$ and $\WF(D)$,
  \item\label{it:c3-wf-part} $S\subseteq\rho_f$, and
  \item\label{it:c3-wf-agree} $\proj{D}{\qq}=\proj{\unfold{K}}{\qq}$ for every $\qq\in S$.
  \end{enumerate}
\end{definition}

Every source choreography is well-formed, since the conditions concern only
runtime terms.  Both non-trivial conditions are needed.

\begin{example}[Ill-formed runtime terms]\label{ex:c3-notwf}
  Let $\Sig(f)=\farr{}{\{\pp,\qq\}}{\qq}$ with
  $f()=\bind{y}{\com{\pp}{1}{\qq}}{\cret{\qq.y}}$, and
  $K=\bind{x}{f()}{\cret{\qq.x}}$.
  \begin{itemize}
  \item The term $\rt{\{\qq\}}{K}{(\cret{\qq.5})}$ violates
    \ref{it:c3-wf-agree}: $\proj{\cret{\qq.5}}{\qq}=\kw{ret}\;5$, but
    $\proj{\unfold{K}}{\qq}=y\leftarrow\nrecv{\pp}\,;\,\kw{ret}\;y$.  By
    \rn{Enter} it steps to $\cret{\qq.5}$, which is terminated.  Its
    projection maps $\qq$ to $x\leftarrow f()\,;\,\kw{ret}\;x$, and once $\qq$
    unfolds its call it waits forever for a message from $\pp$.
  \item With $D=\unfold{K}$, the term $\rt{\{\tr\}}{K}{D}$ for a role
    $\tr\notin\{\pp,\qq\}$ satisfies \ref{it:c3-wf-agree}, since both sides
    are $\kw{ret}\;\bot$ at $\tr$, but violates \ref{it:c3-wf-part}.  By
    \rn{Enter} it steps with $\tauat{\tr}$, but in its projection $\tr$'s
    program is $\kw{ret}\;\bot$, which cannot step.
  \end{itemize}
\end{example}

\section{Correctness}
\label{sec:c3-correct}

All lemmas of Chapters~\ref{ch:core} and~\ref{ch:cond} hold for the extended
language.  Lemma~\ref{lem:c2-erase} (erasure) and Lemma~\ref{lem:c2-pbind}
(projection commutes with bind) are only ever applied with a source term in
the position where they inspect roles, namely $C$ in the first and $C_1$ in
the second, and we state them for that case.  The new lemmas come first,
followed by the new cases of the old ones.

\begin{lemma}[Substitution and well-formedness]\label{lem:c3-wfsubst}
  Substitution commutes with bind and unfolding, and preserves
  well-formedness:
  \begin{enumerate}[label=(\roman*),nosep]
  \item\label{it:c3-ws-bind} If $x\neq y$ and $x\notin\fv(a)$, then
    $(\mbind{C_1}{x}{C_2})[a/y]=\mbind{C_1[a/y]}{x}{C_2[a/y]}$; hence
    $\unfold{K}[a/y]=\unfold{K[a/y]}$.
  \item\label{it:c3-ws-subst} If $\WF(C)$ then $\WF(C[a/y])$.
  \item\label{it:c3-ws-mbind} If $C_1$ is a source choreography and $\WF(C_2)$,
    then $\WF(\mbind{C_1}{x}{C_2})$.
  \end{enumerate}
\end{lemma}
\begin{proof}
  \ref{it:c3-ws-bind} By induction on $C_1$.  For unfolding, the only free
  variables of $C_f$ are its parameters, so
  $C_f[\vec a/\vec x][a/y]=C_f[\vec a[a/y]/\vec x]$.
  \ref{it:c3-ws-subst} By induction on $C$.  For a runtime term, conditions
  \ref{it:c3-wf-sub} and \ref{it:c3-wf-part} follow from the induction
  hypothesis and the fact that substitution does not change the function
  called, and for \ref{it:c3-wf-agree}, by Lemma~\ref{lem:c1-psubst} and
  part~\ref{it:c3-ws-bind},
  $\proj{D[a/y]}{\qq}=\proj{D}{\qq}[a/y]=\proj{\unfold{K}}{\qq}[a/y]=\proj{\unfold{K[a/y]}}{\qq}$.
  \ref{it:c3-ws-mbind} The runtime terms of $\mbind{C_1}{x}{C_2}$ are those of
  $C_2[b/x]$, where $\cret{\sr.b}$ ends the spine of $C_1$; apply
  part~\ref{it:c3-ws-subst}.
\end{proof}

\begin{lemma}[Unfolding]\label{lem:c3-unfold}
  Let $\Gamma\vdash K:\sr$ with $K=\bind{x}{f(\vec\alpha)}{C}$.
  \begin{enumerate}[label=(\roman*),nosep]
  \item\label{it:c3-unf-in} If $\rr\in\rho_f$, then
    $\proj{K}{\rr}=x\leftarrow f(\vec\alpha|_\rr)\,;\,\proj{C}{\rr}$, and
    \rn{N-Call} replaces it by $\proj{\unfold{K}}{\rr}$.
  \item\label{it:c3-unf-out} If $\rr\notin\rho_f$, then $\proj{\unfold{K}}{\rr}=\proj{K}{\rr}$.
  \end{enumerate}
\end{lemma}
\begin{proof}
  By typing, $\Gamma\vdash C_f[\vec a/\vec x]:\sr_f$ (Lemma~\ref{lem:c1-subst},
  after weakening) and $\Gamma,x{:}\sr_f\vdash C:\sr$, so
  Lemma~\ref{lem:c2-pbind} gives
  $\proj{\unfold{K}}{\rr}=\mbind{\proj{C_f[\vec a/\vec x]}{\rr}}{x}{\proj{C}{\rr}}$.

  \ref{it:c3-unf-in} \rn{N-Call} produces
  $\mbind{\proj{C_f}{\rr}[\vec\alpha|_\rr/\vec x|_\rr]}{x}{\proj{C}{\rr}}$.
  By Lemma~\ref{lem:c1-psubst},
  $\proj{C_f[\vec a/\vec x]}{\rr}=\proj{C_f}{\rr}[\vec a/\vec x]$, and by
  Lemma~\ref{lem:c1-local}, applied to $\vec x{:}\vec\pp\vdash C_f:\sr_f$, the
  only parameters free in $\proj{C_f}{\rr}$ are those in $\vec x|_\rr$.  So
  the two agree.

  \ref{it:c3-unf-out} The prefix of the call at $\rr$ is empty, so
  $\proj{K}{\rr}=\proj{C}{\rr}$.  Since
  $\rr\notin\rho_f\supseteq\roles{C_f}=\roles{C_f[\vec a/\vec x]}$,
  Lemma~\ref{lem:c2-erase} gives
  $\proj{\unfold{K}}{\rr}=\mbind{(\kw{ret}\;\bot)}{x}{\proj{C}{\rr}}=\proj{C}{\rr}$.
\end{proof}

\begin{lemma}[View]\label{lem:c3-view}
  Let $\Gamma\vdash K:\sr$ with $K=\bind{x}{f(\vec\alpha)}{C}$, and let
  $\widehat{K}=\rt{\rho_f}{K}{\unfold{K}}$.
  \begin{enumerate}[label=(\roman*),nosep]
  \item\label{it:c3-view-proj} $\Gamma\vdash \widehat{K}:\sr$ and $\projn{\widehat{K}}=\projn{K}$.
  \item\label{it:c3-view-wf} If $\WF(K)$ then $\WF(\widehat{K})$.
  \item\label{it:c3-view-step} $K\lto{\tauat{\rr}}C'$ by \rn{Call} if and only
    if $\widehat{K}\lto{\tauat{\rr}}C'$ by \rn{Enter}.
  \end{enumerate}
\end{lemma}
\begin{proof}
  \ref{it:c3-view-proj} Typing follows from Lemma~\ref{lem:c2-tbind}.  At
  $\rr\in\rho_f$ both projections are $\proj{K}{\rr}$, and at
  $\rr\notin\rho_f$, $\proj{\widehat{K}}{\rr}=\proj{\unfold{K}}{\rr}=\proj{K}{\rr}$ by
  Lemma~\ref{lem:c3-unfold}\ref{it:c3-unf-out}.
  \ref{it:c3-view-wf} $\WF(\unfold{K})$ by
  Lemma~\ref{lem:c3-wfsubst}\ref{it:c3-ws-mbind}, since $C_f[\vec a/\vec x]$
  is a source choreography and $\WF(C)$.  Condition \ref{it:c3-wf-part} is
  trivial, and \ref{it:c3-wf-agree} holds because the current state is
  $\unfold{K}$ itself.
  \ref{it:c3-view-step} Compare the two rules.
\end{proof}

\begin{lemma}[Entering]\label{lem:c3-entering}
  If $\Gamma\vdash\rt{S}{K}{D}:\sr$, $\WF(\rt{S}{K}{D})$ and $\rr\in S$, then
  \[
    \projn{\rt{S}{K}{D}} \lto{\tauat{\rr}} \projn{\rt{S\setminus\{\rr\}}{K}{D}}.
  \]
\end{lemma}
\begin{proof}
  Let $K=\bind{x}{f(\vec\alpha)}{C}$.  Since $\rr\in S\subseteq\rho_f$ by
  \ref{it:c3-wf-part}, the projection at $\rr$ is
  $\proj{K}{\rr}=x\leftarrow f(\vec\alpha|_\rr)\,;\,\proj{C}{\rr}$, and by
  Lemma~\ref{lem:c3-unfold}\ref{it:c3-unf-in} \rn{N-Call} replaces it by
  $\proj{\unfold{K}}{\rr}$, which is $\proj{D}{\rr}$ by \ref{it:c3-wf-agree}.
  That is the projection of $\rt{S\setminus\{\rr\}}{K}{D}$ at $\rr$, since
  $\rr\notin S\setminus\{\rr\}$.  At every other role $\tr$, $\tr\in S$ if
  and only if $\tr\in S\setminus\{\rr\}$, so the two projections agree.  This
  covers the case $S=\{\rr\}$, where the target is $D$.
\end{proof}

\paragraph{New cases of the supporting lemmas.}
\begin{itemize}
\item \emph{Substitution} and \emph{subject reduction}
  (Lemmas~\ref{lem:c1-subst} and~\ref{lem:c1-sr}): arguments are checked at
  their roles like any atom, and runtime terms componentwise.  For \rn{Call},
  $\Gamma\vdash\unfold{K}:\sr$ as in the proof of
  Lemma~\ref{lem:c3-unfold}; \rn{Enter} does not change the components; and
  \rn{Inside} follows from the induction hypothesis.
\item \emph{Projection is local} (Lemma~\ref{lem:c1-local}): a call's prefix
  at $\rr$ mentions only the arguments $\vec\alpha|_\rr$, held at $\rr$; a
  runtime term projects to the projection of one of its components.
\item \emph{Projection commutes with substitution} (Lemma~\ref{lem:c1-psubst}):
  $(\vec\alpha|_\rr)[a/y]=(\vec\alpha[a/y])|_\rr$, and runtime terms are
  handled componentwise, since substitution does not change $S$.
\item \emph{Choreographic frame} (Lemma~\ref{lem:c1-cframe}), for a role
  $\tr\notin\roles{\ell}$.  \rn{Call}: by Lemma~\ref{lem:c3-view},
  $\proj{K}{\tr}=\proj{\widehat{K}}{\tr}$, and \rn{Enter} at $\rr\neq\tr$ does not change
  the projection at $\tr$, as in the proof of Lemma~\ref{lem:c3-entering}.
  \rn{Enter}: the same.  \rn{Inside}: if $\tr\in S$ both sides are
  $\proj{K}{\tr}$; otherwise apply the induction hypothesis to $D$.
\item \emph{Network frame} and \emph{inversion}
  (Lemmas~\ref{lem:c1-nframe} and~\ref{lem:c1-inv}): \rn{N-Call} reads and
  changes only $N(\rr)$, and inversion gains the item
  \begin{enumerate}[label=(\alph*),start=6,nosep]
  \item\label{it:c3-inv-call} if $N(\rr)=x\leftarrow f(\vec a)\,;\,P$, then the
    rule is \rn{N-Call}, $\ell=\tauat{\rr}$, and
    $N'=\upd{N}{\rr}{\mbind{P_f[\vec a/\vec y]}{x}{P}}$ where
    $\projn{\Defs}_\rr(f)=(\vec y).\,P_f$.
  \end{enumerate}
\end{itemize}

\begin{lemma}[Preservation of well-formedness]\label{lem:c3-wfpres}
  If $\Gamma\vdash C:\sr$, $\WF(C)$ and $C\lto{\ell}C'$, then $\WF(C')$.
\end{lemma}
\begin{proof}
  By induction on the step.
  \begin{itemize}[nosep]
  \item \rn{Local} and \rn{Com}: $C'=C_0[v/x]$ with $\WF(C_0)$; use
    Lemma~\ref{lem:c3-wfsubst}\ref{it:c3-ws-subst}.
  \item \rn{Cast}: $C'=\mbind{C_v}{x}{C_0}$ with $C_v$ a source choreography;
    use Lemma~\ref{lem:c3-wfsubst}\ref{it:c3-ws-mbind}.
  \item \rn{Delay}: by the induction hypothesis.
  \item \rn{Call}: $C'$ is the \rn{Enter} successor of the term $\widehat{K}$ of
    Lemma~\ref{lem:c3-view}, which is well-formed, so this is the next case.
  \item \rn{Enter}: condition \ref{it:c3-wf-part} is preserved by removing an
    element from $S$, and \ref{it:c3-wf-agree} by restricting to fewer
    roles.  If $S\setminus\{\rr\}=\emptyset$, then $C'=D$, which is
    well-formed by \ref{it:c3-wf-sub}.
  \item \rn{Inside}: $C=\rt{S}{K}{D}$ and $C'=\rt{S}{K}{D'}$ with
    $D\lto{\ell}D'$ and $S\cap\roles{\ell}=\emptyset$.  $\WF(D')$ holds by the
    induction hypothesis.  Every $\qq\in S$ is outside $\roles{\ell}$, so
    Lemma~\ref{lem:c1-cframe} gives
    $\proj{D'}{\qq}=\proj{D}{\qq}$, which is $\proj{\unfold{K}}{\qq}$.\qedhere
  \end{itemize}
\end{proof}

\begin{theorem}[EPP is lockstep]\label{thm:c3-epp}
  Suppose $\Gamma\vdash C:\sr$ and $\WF(C)$.
  \begin{enumerate}[label=(\arabic*),nosep]
  \item\label{it:c3-complete} If $C\lto{\ell}C'$, then $\projn{C}\lto{\ell}\projn{C'}$.
  \item\label{it:c3-sound} If $\projn{C}\lto{\ell}N$, then $C\lto{\ell}C'$ for
    some $C'$ with $N=\projn{C'}$.
  \end{enumerate}
\end{theorem}
\begin{proof}
  As for Theorem~\ref{thm:c2-epp}.  Subterms of well-formed choreographies
  that the inductions reach (continuations and the current state $D$) are
  well-formed, so the induction hypotheses apply.  The new cases follow.

  \case{\rn{Enter}, completeness}
  This is Lemma~\ref{lem:c3-entering}.

  \case{\rn{Call}, completeness}
  By Lemma~\ref{lem:c3-view}, $\projn{K}=\projn{\widehat{K}}$, $\widehat{K}$ is well-typed and
  well-formed, and $C'$ is the \rn{Enter} successor of $\widehat{K}$.
  Lemma~\ref{lem:c3-entering} applied to $\widehat{K}$ gives
  $\projn{K}=\projn{\widehat{K}}\lto{\tauat{\rr}}\projn{C'}$.

  \case{\rn{Inside}, completeness}
  $C=\rt{S}{K}{D}$, $C'=\rt{S}{K}{D'}$, $D\lto{\ell}D'$ and
  $S\cap\roles{\ell}=\emptyset$.  The induction hypothesis gives
  $\projn{D}\lto{\ell}\projn{D'}$.  For $\rr\in\roles{\ell}$ we have $\rr\notin S$,
  so $\proj{C}{\rr}=\proj{D}{\rr}$ and $\proj{C'}{\rr}=\proj{D'}{\rr}$; for
  $\rr\notin\roles{\ell}$, $\proj{C'}{\rr}=\proj{C}{\rr}$ by
  Lemma~\ref{lem:c1-cframe}.  Corollary~\ref{cor:c1-lift} concludes.

  \case{Call, soundness}
  $C=\bind{x}{f(\vec\alpha)}{C_0}$, and the step involves some
  $\rr\in\rho_f$.  Then $\projn{C}(\rr)$ is the call $x\leftarrow
  f(\vec\alpha|_\rr)\,;\,\proj{C_0}{\rr}$, so by
  Lemma~\ref{lem:c1-inv}\ref{it:c3-inv-call} (as extended above) the step is \rn{N-Call} at $\rr$,
  $\ell=\tauat{\rr}$, and $N$ changes only $\rr$'s program, to
  $\proj{\unfold{C}}{\rr}$ by Lemma~\ref{lem:c3-unfold}\ref{it:c3-unf-in}.
  Rule \rn{Call} gives $C\lto{\tauat{\rr}}C'=\rt{\rho_f\setminus\{\rr\}}{C}{\unfold{C}}$.
  By the completeness case for \rn{Call}, $\projn{C}\lto{\tauat{\rr}}\projn{C'}$
  by \rn{N-Call} at $\rr$, and \rn{N-Call} has a unique target, so
  $\projn{C'}=N$.

  \case{$C=\rt{S}{K}{D}$, soundness}
  We distinguish whether the step involves a role that has not entered.
  \begin{itemize}
  \item If some $\rr\in S\cap\roles{\ell}$, then $\rr\in\rho_f$ by
    \ref{it:c3-wf-part}, so $\projn{C}(\rr)=\proj{K}{\rr}$ is a call, and by
    Lemma~\ref{lem:c1-inv}\ref{it:c3-inv-call} (as extended above) the step is \rn{N-Call} at
    $\rr$ with $\ell=\tauat{\rr}$.  Rule \rn{Enter} gives
    $C\lto{\tauat{\rr}}\rt{S\setminus\{\rr\}}{K}{D}$, whose projection is the
    target of this step, by Lemma~\ref{lem:c3-entering}.
  \item If $S\cap\roles{\ell}=\emptyset$, then $\projn{C}$ agrees with
    $\projn{D}$ on $\roles{\ell}$.  By
    Lemma~\ref{lem:c1-nframe}\ref{it:c1-nframe-move},
    $\projn{D}\lto{\ell}N''$, where $N''$ agrees with $N$ on $\roles{\ell}$
    and with $\projn{D}$ elsewhere.  The induction hypothesis gives
    $D\lto{\ell}D'$ with $\projn{D'}=N''$, and \rn{Inside} gives
    $C\lto{\ell}C'=\rt{S}{K}{D'}$.  For $\rr\in\roles{\ell}$,
    $\proj{C'}{\rr}=\proj{D'}{\rr}=N(\rr)$; for $\rr\notin\roles{\ell}$,
    $N(\rr)=\proj{C}{\rr}=\proj{C'}{\rr}$ by
    Lemma~\ref{lem:c1-nframe}\ref{it:c1-nframe-loc} and
    Lemma~\ref{lem:c1-cframe}.\qedhere
  \end{itemize}
\end{proof}

\begin{corollary}[Bisimulation]\label{cor:c3-bisim}
  $\{(C,\projn{C})\mid\Gamma\vdash C:\sr \text{ and } \WF(C)\}$ is a strong
  bisimulation with identical labels.  In particular this holds from every
  source choreography.
\end{corollary}
\begin{proof}
  Theorem~\ref{thm:c3-epp}, with subject reduction and
  Lemma~\ref{lem:c3-wfpres}.
\end{proof}

\begin{theorem}[Deadlock freedom]\label{thm:c3-df}
  Assume (L2), (L3), $\emptyset\vdash C:\sr$ and $\WF(C)$.  Every network
  reachable from $\projn{C}$ is either terminated or can step.
\end{theorem}
\begin{proof}
  As for Theorem~\ref{thm:c2-df}, using Corollary~\ref{cor:c3-bisim}.  It
  remains to show that every closed, well-formed $C$ that is not a return can
  step.  If its head action is a call to $f$, \rn{Call} fires for any
  $\rr\in\rho_f$, which is nonempty because it contains $\sr_f$.  If $C$ is a
  runtime term $\rt{S}{K}{D}$, then $S\neq\emptyset$ and \rn{Enter} fires for
  any $\rr\in S$.  Neither rule has a premise on the arguments.
\end{proof}

\begin{remark}[Confluence]\label{rem:c3-diamond}
  The network lemmas of \S\ref{sec:c1-confluence} hold with \rn{N-Call}, whose
  label $\tauat{\rr}$ is forced by a call at the head of $N(\rr)$.  So under
  (L1)--(L3), for a closed, well-typed source program, if one maximal run is
  finite then all are, and they have the same length, end in the same
  return, and show the same labels at each role.  The corresponding diamond property fails syntactically for
  choreographies.  Let $\Sig(f)=\farr{}{\{\qq,\rr\}}{\rr}$ with
  $f()=\bind{y}{\qq.1}{\bind{z}{\com{\qq}{y}{\rr}}{\cret{\rr.z}}}$, and
  $C=\bind{x}{f()}{\bind{w}{\tr.1}{\cret{\rr.x}}}$.  Taking $\tauat{\qq}$
  (\rn{Call}) and then $\tauat{\tr}$ freezes $K=C$; taking them in the other
  order freezes $K=\bind{x}{f()}{\cret{\rr.x}}$.  The two results have the
  same projection, but differ as terms.  This is why confluence is proved on
  networks.
\end{remark}

\section{Worked example}
\label{sec:c3-example}

The following function is a loop.  Role $\pp$ doubles a number until it
exceeds $10$, $\qq$ counts the rounds that did not end the loop, and when the
loop ends the number is forwarded to $\rr$ via $\qq$.  Let
$\Sig(\mathit{loop})=\farr{\pp,\qq}{\{\pp,\qq,\rr\}}{\rr}$ and
\[
\begin{array}{r@{\;\;}l}
  \mathit{loop}(x,k) \;=\; & \bind{y}{\pp.(2\cdot x)}{}\\
    & \bind{b}{\pp.(y>10)}{}\\
    & \bind{z}{\cif{\pp}{b}{T(y)}{E(y,k)}}{}\\
    & \cret{\rr.z}\\[1ex]
  T(y) \;=\; & \bind{w}{\com{\pp}{y}{\qq}}{\bind{u}{\com{\qq}{w}{\rr}}{\cret{\rr.u}}}\\
  E(y,k) \;=\; & \bind{k'}{\qq.(k+1)}{\bind{v}{\mathit{loop}(\pp.y,\qq.k')}{\cret{\rr.v}}}
\end{array}
\]
The conditional's receivers are $\{\qq,\rr\}$.  Role $\rr$ occurs in $T$, and
it occurs in $E$ both as a participant of the recursive call and as the holder
of the result.  The projected local definitions are
\[
\begin{array}{r@{\;\;}l}
  \mathit{loop}_\pp(x) \;=\; & y\leftarrow 2\cdot x\,;\; b\leftarrow y>10\,;\\
     & z\leftarrow\sendbr{\{\qq,\rr\}}{b}{(\nsend{\qq}{y}\,;\,\kw{ret}\;\bot)}{(v\leftarrow\mathit{loop}(y)\,;\,\kw{ret}\;\bot)}\,;\;\kw{ret}\;\bot\\[1ex]
  \mathit{loop}_\qq(k) \;=\; & z\leftarrow\pp?\big((w\leftarrow\nrecv{\pp}\,;\,\nsend{\rr}{w}\,;\,\kw{ret}\;\bot),\\
     & \hphantom{z\leftarrow\pp?\big(}(k'\leftarrow k+1\,;\,v\leftarrow\mathit{loop}(k')\,;\,\kw{ret}\;\bot)\big)\,;\;\kw{ret}\;\bot\\[1ex]
  \mathit{loop}_\rr() \;=\; & z\leftarrow\recvbr{\pp}{(u\leftarrow\nrecv{\qq}\,;\,\kw{ret}\;u)}{(v\leftarrow\mathit{loop}()\,;\,\kw{ret}\;v)}\,;\;\kw{ret}\;z
\end{array}
\]
The program calls the loop and uses its result:
\[
  C \;=\; \bind{\mathit{res}}{\mathit{loop}(\pp.3,\qq.0)}{\bind{m}{\rr.(\mathit{res}+1)}{\cret{\rr.m}}},
\]
which projects to $\mathit{res}\leftarrow\mathit{loop}(3)\,;\,\kw{ret}\;\bot$
at $\pp$, $\mathit{res}\leftarrow\mathit{loop}(0)\,;\,\kw{ret}\;\bot$ at $\qq$,
and $\mathit{res}\leftarrow\mathit{loop}()\,;\,m\leftarrow\mathit{res}+1\,;\,\kw{ret}\;m$
at $\rr$.  Write $B(y,k)$ for the conditional
$\cif{\pp}{b}{T(y)}{E(y,k)}$, $R(z)$ for the continuation
$\bind{m}{\rr.(z+1)}{\cret{\rr.m}}$, and $L_2$ for the call
$\mathit{loop}(\pp.6,\qq.k')$.  One run begins
\begin{align*}
  C &\lto{\tauat{\pp}} \rt{\{\qq,\rr\}}{C}{D_1}
    && \text{\rn{Call}}\\
  &\lto{\tauat{\pp}}\;\lto{\tauat{\pp}} \rt{\{\qq,\rr\}}{C}{(\bind{z}{B(6,0)[\ffalse/b]}{R(z)})}
    && \text{\rn{Inside}, \rn{Local}, twice}\\
  &\lto{\tauat{\qq}}\;\lto{\tauat{\rr}} \bind{z}{B(6,0)[\ffalse/b]}{R(z)}
    && \text{\rn{Enter}, twice}\\
  &\lto{\castlbl{\pp}{\ffalse}{\{\qq,\rr\}}} \bind{k'}{\qq.(0+1)}{\bind{v}{L_2}{R(v)}}
    && \text{\rn{Cast}}\\
  &\lto{\tauat{\pp}} \bind{k'}{\qq.(0+1)}{\rt{\{\qq,\rr\}}{K_2}{D_2}}
    && \text{\rn{Delay}, \rn{Call}}
\end{align*}
where $D_1=\unfold{C}=\bind{y}{\pp.(2\cdot 3)}{\bind{b}{\pp.(y>10)}{\bind{z}{B(y,0)}{R(z)}}}$,
$K_2=\bind{v}{L_2}{R(v)}$ and
$D_2=\bind{y}{\pp.(2\cdot 6)}{\bind{b}{\pp.(y>10)}{\bind{z}{B(y,k')}{R(z)}}}$.

Some points of this run:
\begin{itemize}
\item In the first three steps $\pp$ runs ahead inside the call while $\qq$
  and $\rr$ have not entered.  The multicast is not enabled until both have,
  because \rn{Inside} forbids a step involving a role in $S$.
\item After the multicast the continuation $R$ follows the selected branch
  directly.  The recursive call is not in tail position: its result $v$ flows
  into $R(v)$.
\item In the last step $\pp$ enters the recursive call by \rn{Delay}, past
  $\qq$'s pending increment, while $\qq$'s argument $k'$ is still a variable
  (Choice~\ref{cho:c3-novalue}).  When $\qq$ computes $k'=1$, \rn{Local}
  substitutes $1$ into both $K_2$ and $D_2$.  Meanwhile $\rr$ may enter the
  recursive call at any time.  In the network, $\rr$'s program at this point
  is $v\leftarrow\mathit{loop}()\,;\,m\leftarrow v+1\,;\,\kw{ret}\;m$, which
  is $\proj{K_2}{\rr}$.
\end{itemize}
In the second round, $2\cdot 6=12>10$, the multicast announces $\ttrue$, the
value $12$ travels from $\pp$ to $\qq$ to $\rr$, and the run ends in
$\cret{\rr.13}$.

%======================================================================
\chapter{Higher-Order Functions}
\label{ch:ho}
%======================================================================

This chapter removes named functions altogether: signatures, definitions,
declared role sets, and the local definitions of the network semantics.  In
their place, functions become values, written as recursive
$\lambda$-abstractions in the choreography:
\[
  \cfun{f}{x_1{:}T_1,\ldots,x_n{:}T_n}{\rho}{T}{C}.
\]
A function value is \emph{located at the set of roles} $\rho$ that run it.  It
can be passed as an argument, returned as the result of a call or of a
conditional, and bound by the continuation.  Located data is the special case
of a value held at a single role.

Two decisions keep the proofs close to those of Chapter~\ref{ch:fun}.  A
function body is part of a value, and nothing evaluates under it: \rn{Delay}
steps only under bind.  And an application carries its role set,
$\rho.F(\vec\alpha)$, so $\roles{\cdot}$ never looks inside a function.
Applications are then entered one role at a time as calls were in
Chapter~\ref{ch:fun}, with the same runtime term and the same
well-formedness invariant, and the proofs of Chapter~\ref{ch:fun} go through
with a new unfolding lemma.

\section{Design choices}
\label{sec:c4-choices}

\begin{choice}[Functions are syntax, located at role sets]\label{cho:c4-located}
  A function value is a literal or a variable.  Its type
  $\farr{\vec T}{\rho}{T}$ records the roles $\rho$ that run it, and a function
  is held at exactly those roles: $\loc{\farr{\vec T}{\rho}{T}}=\rho$, just as
  $\loc{\pp}=\{\pp\}$ for data held at $\pp$.  Every role in $\rho$ has its
  own projection of every literal it mentions, so no code is ever sent.

  \emph{Proof saved.}  Substituting a function for a variable commutes with
  projection, because each role receives its own projection of the function
  (Lemma~\ref{lem:c4-psubst}).  One locality lemma covers data and functions
  alike: a role mentions a variable only if it belongs to the variable's
  location set (Lemma~\ref{lem:c4-local}).  The \emph{ret} case of the bind
  lemma becomes a split on whether $\rr$ is in that set, with the same proof
  as before.
\end{choice}

\begin{choice}[Applications carry their role set]\label{cho:c4-annot}
  An application is written $\rho.F(\vec\alpha)$.  Its set of roles is
  $\rho$, read off the annotation and never off the body of $F$.  The
  annotation agrees with $F$'s type, so a type checker can fill it in.

  \emph{Proof saved.}  $\roles{\cdot}$ is stable under substitution, including
  substitution of a function for a variable.  Receivers of conditionals,
  erasure decisions and annotations on projected send-and-branches therefore
  do not change as functions flow through a program.  The annotation does
  the job that the declared role set $\rho_f$ did in Chapter~\ref{ch:fun}:
  typing a literal at $\rho$ checks that its body's roles are included in
  $\rho$, and no fixed point is needed.
\end{choice}

\begin{choice}[No evaluation under $\lambda$]\label{cho:c4-nolambda}
  A function body does not step until some role enters an application of the
  function.  \rn{Delay} steps only under bind, and no rule steps inside a
  value.

  \emph{Proof saved.}  Function bodies, like branches, remain source terms,
  so the erasure and bind lemmas apply to them.  Choice~\ref{cho:c4-annot}
  depends on this: since a body cannot have advanced, its roles are bounded
  by its type.
\end{choice}

\begin{choice}[Applications are entered one role at a time]\label{cho:c4-enter}
  The runtime term, the rules \rn{Call}, \rn{Enter} and \rn{Inside}, and the
  well-formedness invariant are those of Chapter~\ref{ch:fun}, with a literal
  in place of a function name and the annotation $\rho$ in place of $\rho_f$.

  \emph{Proof saved.}  The proofs of Chapter~\ref{ch:fun} apply verbatim.
  Apart from the lemmas restated for location sets (\S\ref{sec:c4-correct}),
  only the unfolding lemma is re-proved (Lemma~\ref{lem:c4-unfold}).
\end{choice}

\begin{choice}[A literal need not be closed to be applied]\label{cho:c4-open}
  Following Choice~\ref{cho:c3-novalue}, nothing in an application must be a
  value.  In particular a literal may still have free variables, held at
  other roles and bound by pending actions, when a role enters it.

  \emph{Proof saved.}  Requiring literals to be closed would break
  soundness: in Example~\ref{ex:c4-mk}, $\qq$'s projection of a literal is closed and
  $\qq$ can enter it in the network while the literal still mentions a
  variable held at $\pp$.  Without the requirement, a free variable held at
  $\rr$ blocks $\rr$ through \rn{Delay}, as in
  Remark~\ref{rem:c1-delay}, and a free variable held elsewhere does not appear
  in $\rr$'s projection.
\end{choice}

\begin{choice}[Exact types]\label{cho:c4-exact}
  There is no subtyping between function types with different role sets.

  \emph{Proof saved.}  Projection keeps a parameter at $\rr$ exactly when
  $\rr$ is in the location of its type, and an argument exactly when $\rr$ is
  in the location of the argument.  With exact types the two agree, so at
  every role the projected function and the projected call have the same
  arity.  Subtyping on parameter types would break this.
\end{choice}

\begin{choice}[Functions are not communicated]\label{cho:c4-nocomm}
  Only data atoms can be sent, and a function's location set is fixed by its
  type.

  \emph{Proof saved.}  A role learns which function a variable denotes only
  through its own projection of the program, so the local language needs no representation of foreign code and the network
  semantics needs no rule for it.
\end{choice}

\section{Syntax and typing}
\label{sec:c4-syntax}

\begin{definition}[Types]\label{def:c4-types}
  Types are
  \[
    T \;::=\; \pp \;\mid\; \farr{T_1,\ldots,T_n}{\rho}{T}
    \qquad
    \loc{\pp}=\{\pp\}
    \qquad
    \loc{\farr{\vec T}{\rho}{T}}=\rho.
  \]
  A type is \emph{valid} if it is a role, or a function type
  $\farr{\vec T}{\rho}{T}$ whose components are valid and whose result is held
  by participants, $\loc{T}\subseteq\rho$.  Every valid type has a nonempty
  location.  All types below are valid.
\end{definition}

\begin{definition}[Choreographies with functions]\label{def:c4-syntax}
  \[
  \begin{array}{r@{\;\;}c@{\;\;}l@{\qquad}l}
    F & ::= & g & \text{function variable}\\
      & \mid & \cfun{f}{x_1{:}T_1,\ldots,x_n{:}T_n}{\rho}{T}{C} & \text{recursive function literal}\\[1ex]
    \alpha & ::= & \pp.a \;\mid\; \rho.F & \text{located argument}\\[1ex]
    A & ::= & \pp.e \;\mid\; \com{\pp}{a}{\qq} \;\mid\; \cif{\pp}{a}{C_1}{C_2} & \text{as in Chapter~\ref{ch:cond}}\\
      & \mid & \rho.F(\alpha_1,\ldots,\alpha_n) & \text{application, run at } \rho\\[1ex]
    C & ::= & \cret{\alpha} \;\mid\; \bind{x}{A}{C} \;\mid\; \rt{S}{K}{D}\\[1ex]
    K & ::= & \bind{x}{\rho.V(\vec\alpha)}{C} & V \text{ a literal}
  \end{array}
  \]
  In a literal, $f$ (its own name) and $x_1,\ldots,x_n$ are bound in its body.
  Literals are ranged over by $V$, and the type of
  $V=\cfun{f}{\vec x{:}\vec T}{\rho}{T}{C_V}$ is $T_V=\farr{\vec T}{\rho}{T}$.
  An argument has a location and a payload:
  $\loc{\pp.a}=\{\pp\}$, $\loc{\rho.F}=\rho$, $\pay{\pp.a}=a$ and
  $\pay{\rho.F}=F$.  Source choreographies contain no runtime term; branches
  and function bodies are source choreographies.
\end{definition}

A single set of variables serves data and functions; typing tells them apart.
Substitution $C[u/x]$ replaces a variable by an atom or by a function value,
and goes into function bodies, avoiding capture.  Derived bind becomes
$\mbind{(\cret{\alpha})}{x}{C}=C[\pay{\alpha}/x]$.

\paragraph{Roles.}
$\roles{\cret{\alpha}}=\loc{\alpha}$ and $\roles{\rho.F(\vec\alpha)}=\rho$;
the other cases are as in Definition~\ref{def:c2-roles}.  The roles of a
choreography never depend on the contents of a function value, so
$\roles{C[u/x]}=\roles{C}$ for every substitution.

\paragraph{Typing.}
A context $\Gamma$ maps variables to types.  $\Gamma\vdash_\pp e$ now means
$\Gamma(x)=\pp$ for every $x\in\fv(e)$, so local expressions and sent atoms
mention only data.  The judgements are $\Gamma\vdash F:T$ for function values,
$\Gamma\vdash\alpha:T$ for arguments, and $\Gamma\vdash A:T$ and $\Gamma\vdash
C:T$ for actions and choreographies.
\begin{mathpar}
  \inferrule{\Gamma(g)=T}{\Gamma\vdash g:T}
  \and
  \inferrule{\Gamma,\,f{:}T_V,\,\vec x{:}\vec T\vdash C:T \\ \roles{C}\subseteq\rho}
            {\Gamma\vdash\cfun{f}{\vec x{:}\vec T}{\rho}{T}{C}:T_V}
  \\
  \inferrule{\Gamma\vdash_\pp a}{\Gamma\vdash\pp.a:\pp}
  \and
  \inferrule{\Gamma\vdash F:T \\ \loc{T}=\rho}{\Gamma\vdash\rho.F:T}
  \and
  \inferrule{\Gamma\vdash\alpha:T}{\Gamma\vdash\cret{\alpha}:T}
  \and
  \inferrule{\Gamma\vdash F:\farr{T_1,\ldots,T_n}{\rho}{T} \\ \Gamma\vdash\alpha_i:T_i\ \text{for each } i}
            {\Gamma\vdash\rho.F(\alpha_1,\ldots,\alpha_n):T}
\end{mathpar}
together with the rules of \S\ref{sec:c2-syntax} for bind, local computation,
communication and conditionals, with types in place of roles, and the rule
for runtime terms of \S\ref{sec:c3-syntax}.  A program is a choreography $C$
with $\emptyset\vdash C:T$.

The literal rule is where the annotation is checked, and it allows any
$\rho$ that covers the body's roles.  Since $\roles{\cdot}$ uses the
annotations of nested applications, including recursive ones, which carry
$\rho$ itself, checking is a simple inclusion.

\begin{lemma}[Results are held by participants]\label{lem:c4-locroles}
  If $\Gamma\vdash A:T$ then $\loc{T}\subseteq\roles{A}$, and if $C$ is a source
  choreography with $\Gamma\vdash C:T$, then $\loc{T}\subseteq\roles{C}$.
\end{lemma}
\begin{proof}
  As for Lemma~\ref{lem:c2-locroles}.  For a return,
  $\loc{T}=\loc{\alpha}=\roles{\cret{\alpha}}$.  For an application,
  $\loc{T}\subseteq\rho$ because the function's type is valid.
\end{proof}

\section{Choreographic semantics}
\label{sec:c4-sem}

\begin{definition}[Unfolding]\label{def:c4-unfold}
  For $K=\bind{x}{\rho.V(\vec\alpha)}{C}$ with
  $V=\cfun{f}{\vec x{:}\vec T}{\rho}{T}{C_V}$,
  \[
    \unfold{K} \;=\; \mbind{C_V[V/f][\pay{\vec\alpha}/\vec x]}{x}{C}.
  \]
\end{definition}

The semantics consists of \rn{Local}, \rn{Com}, \rn{Delay} and \rn{Cast},
unchanged, and the rules of Figure~\ref{fig:c3-sem} with
$K=\bind{x}{\rho.V(\vec\alpha)}{C}$ and $\rho$ in place of $\rho_f$:
\begin{mathpar}
  \inferrule*[right=\rn{Call}]
    {K=\bind{x}{\rho.V(\vec\alpha)}{C} \\ \rr\in\rho}
    {K \lto{\tauat{\rr}} \rt{\rho\setminus\{\rr\}}{K}{\unfold{K}}}
\end{mathpar}
together with \rn{Enter} and \rn{Inside} as before.  Labels are unchanged.
\rn{Call} requires the function position to be a literal.  For a closed
program this is no restriction.  A variable in function position is bound
either by an enclosing literal, as its name or a parameter, and then the
application cannot run before \rn{Call} substitutes the variable; or by a
pending bind $\bind{g}{A}{\cdots}$, and then the application's role set is
$\loc{T}\subseteq\roles{A}$, where $T$ is the type of $g$
(Lemma~\ref{lem:c4-locroles}), so \rn{Delay} keeps all of its roles waiting
until $A$ has fired.  The body of a literal is not touched by any rule until
\rn{Call} unfolds it.

A function returned from a call or a conditional is bound by the continuation
like any other result.  For instance, in
$\bind{g}{\cif{\pp}{b}{(\cret{\rho.V_1})}{(\cret{\rho.V_2})}}{\bind{y}{\rho.g(\vec\alpha)}{C}}$,
every role of $\rho$ other than $\pp$ is a receiver of the conditional,
because $\roles{\cret{\rho.V_i}}=\rho$, so every role that will run $g$
learns which function it is.

\section{Networks}
\label{sec:c4-net}

\begin{definition}[Networks with local functions]\label{def:c4-net}
  Local programs gain function values and application:
  \[
  \begin{array}{r@{\;\;}c@{\;\;}l@{\qquad}l}
    W & ::= & g \;\mid\; \nfun{f}{\vec y}{P} & \text{local function}\\
    w & ::= & a \;\mid\; W & \text{local value or atom}\\
    P & ::= & \kw{ret}\;w \;\mid\; \kw{ret}\;\bot \;\mid\; \pi\,P\\
    \pi & ::= & \cdots \;\mid\; x\leftarrow W(\vec w)\,; & \text{application}
  \end{array}
  \]
  Network bind is as in Definition~\ref{def:c2-net} with
  $\mbind{(\kw{ret}\;w)}{x}{P}=P[w/x]$.  The network semantics consists of the
  rules of \S\ref{sec:c2-net} and
  \begin{mathpar}
    \inferrule*[right=\rn{N-Call}]
      {N(\rr)=x\leftarrow W(\vec w)\,;\,P \\ W=\nfun{f}{\vec y}{P_W} \\ |\vec y|=|\vec w|}
      {N \lto{\tauat{\rr}} \upd{N}{\rr}{\mbind{P_W[W/f][\vec w/\vec y]}{x}{P}}}
  \end{mathpar}
\end{definition}

There is no environment of local definitions.  Each role's program contains
its own projection of every function it may run.

\section{Endpoint projection}
\label{sec:c4-epp}

\begin{definition}[Projection with functions]\label{def:c4-epp}
  Function values and arguments project as follows, where an argument is
  projected only at roles in its location:
  \[
  \begin{array}{r@{\;\;=\;\;}l@{\qquad\qquad}r@{\;\;=\;\;}l}
    \proj{g}{\rr} & g & \proj{\pp.a}{\rr} & a\\
    \proj{\cfun{f}{\vec x{:}\vec T}{\rho}{T}{C}}{\rr} & \nfun{f}{\vec x|_\rr}{\proj{C}{\rr}} & \proj{\rho.F}{\rr} & \proj{F}{\rr}
  \end{array}
  \]
  Here $\vec x|_\rr$ keeps the parameters $x_i$ with $\rr\in\loc{T_i}$, and for
  arguments, $\vec\alpha|_\rr$ is the sequence of $\proj{\alpha_i}{\rr}$ for the
  $\alpha_i$ with $\rr\in\loc{\alpha_i}$.  Returns and applications project as
  \[
  \proj{\cret{\alpha}}{\rr} =
    \begin{cases}
      \kw{ret}\;\proj{\alpha}{\rr} & \rr\in\loc{\alpha}\\
      \kw{ret}\;\bot & \text{otherwise}
    \end{cases}
  \qquad
  \pre{\rho.F(\vec\alpha)}{x}{\rr} =
    \begin{cases}
      x\leftarrow\proj{F}{\rr}(\vec\alpha|_\rr)\,; & \rr\in\rho\\
      \varepsilon & \text{otherwise}
    \end{cases}
  \]
  and all other cases, including runtime terms, are as in
  Definitions~\ref{def:c2-epp} and~\ref{def:c3-epp}.  For substitutions we set
  $\proj{a}{\rr}=a$ for atoms.
\end{definition}

A literal projects to a local function at every role, but it is only ever
projected at roles in its location: as an argument or a result at $\rr\in\rho$,
and in function position at a participant of the application.
Lemma~\ref{lem:c1-prefix} still holds, now in the form: if
$\pre{A}{x}{\rr}$ does not bind $x$, then $\rr\notin\loc{T}$, where $T$ is the
type of $A$.  For an application, the prefix binds $x$ at every role of
$\rho\supseteq\loc{T}$.

\section{Correctness}
\label{sec:c4-correct}

The development of Chapter~\ref{ch:fun} carries over.  Four lemmas change
their statements, because a variable may now be held at several roles
(substitution, locality, and the two commutation lemmas), and the unfolding
lemma is re-proved.  Everything else, including
well-formedness (Definition~\ref{def:c3-wf}, with $S\subseteq\rho$ in
condition~\ref{it:c3-wf-part}), the View and Entering lemmas, preservation of
well-formedness, and both halves of the main theorem, is proved as in
Chapter~\ref{ch:fun}.

\begin{lemma}[Substitution]\label{lem:c4-subst}
  If $\Gamma,x{:}T\vdash C:T'$, and either $T=\pp$ and $\Gamma\vdash_\pp a$ with
  $u=a$, or $\Gamma\vdash F:T$ with $u=F$, then $\Gamma\vdash C[u/x]:T'$.  The
  same holds for function values, arguments and actions.
\end{lemma}
\begin{proof}
  By induction on the typing derivation.  For a data variable the argument is
  that of Lemma~\ref{lem:c1-subst}.  For a function variable, the only rules
  that inspect it are the variable rule, which is replaced by the derivation
  of $\Gamma\vdash F:T$ after weakening, and the literal rule, whose side
  condition $\roles{C}\subseteq\rho$ is unaffected because roles are stable
  under substitution.
\end{proof}

\begin{lemma}[Projection is local]\label{lem:c4-local}
  If $\Gamma\vdash C:T$, then every free variable $x$ of $\proj{C}{\rr}$
  satisfies $\rr\in\loc{\Gamma(x)}$.  The same holds for $\proj{F}{\rr}$ when
  $\Gamma\vdash F:T$ and $\rr\in\loc{T}$.
\end{lemma}
\begin{proof}
  By simultaneous induction on derivations.
  For a variable $g$ with $\rr\in\loc{\Gamma(g)}$, the claim is immediate.
  For a literal $V$ with body $C_V$, the free variables of
  $\proj{V}{\rr}$ are those of $\proj{C_V}{\rr}$ other than $f$ and
  $\vec x|_\rr$.  By the induction hypothesis, those of $\proj{C_V}{\rr}$ are
  held at $\rr$ under $\Gamma,f{:}T_V,\vec x{:}\vec T$.  Among the newly bound
  variables, the ones held at $\rr$ are $f$, if $\rr\in\rho$, and exactly the
  parameters in $\vec x|_\rr$, and these are bound by $\proj{V}{\rr}$.
  For a return, $\proj{\alpha}{\rr}$ is taken only when
  $\rr\in\loc{\alpha}$.  For an application at $\rr\in\rho$, $F$ has a type
  with location $\rho\ni\rr$, and each projected argument has $\rr$ in its
  location.  For a bind, the bound variable is handled by the form of
  Lemma~\ref{lem:c1-prefix} stated above.  The remaining cases are as in
  Lemma~\ref{lem:c1-local}.
\end{proof}

\begin{lemma}[Projection commutes with substitution]\label{lem:c4-psubst}
  For a substitution that respects sorts, an atom for a data variable or a
  function value for a function variable, we have
  $\proj{C[u/x]}{\rr}=\proj{C}{\rr}[\proj{u}{\rr}/x]$, and likewise for
  function values.  Hence, if $\Gamma,x{:}T\vdash C:T'$ and $\rr\notin\loc{T}$,
  then $\proj{C[u/x]}{\rr}=\proj{C}{\rr}$.
\end{lemma}
\begin{proof}
  By induction.  Every choice that projection makes depends on roles,
  annotations, locations of arguments, or parameter types, and substitution
  changes none of them.  The second claim follows by
  Lemma~\ref{lem:c4-local}.
\end{proof}

\begin{lemma}[Projection commutes with bind]\label{lem:c4-pbind}
  If $C_1$ is a source choreography, $\Gamma\vdash C_1:T'$ and
  $\Gamma,x{:}T'\vdash C_2:T$, then
  $\proj{\mbind{C_1}{x}{C_2}}{\rr}=\mbind{\proj{C_1}{\rr}}{x}{\proj{C_2}{\rr}}$.
\end{lemma}
\begin{proof}
  As for Lemma~\ref{lem:c2-pbind}.  In the case $C_1=\cret{\alpha}$, the left
  side is $\proj{C_2}{\rr}[\proj{\pay{\alpha}}{\rr}/x]$.  If
  $\rr\in\loc{\alpha}$, the right side is
  $\mbind{(\kw{ret}\;\proj{\alpha}{\rr})}{x}{\proj{C_2}{\rr}}$, which is the
  same, since $\proj{\alpha}{\rr}=\proj{\pay{\alpha}}{\rr}$.  Otherwise
  $\rr\notin\loc{\alpha}=\loc{T'}$, the right side is $\proj{C_2}{\rr}$, and
  $x$ is not free in it by Lemma~\ref{lem:c4-local}.
\end{proof}

Erasure (Lemma~\ref{lem:c2-erase}) holds for source choreographies with the
same proof, the return case now reading $\rr\notin\loc{\alpha}$.

\begin{lemma}[Unfolding]\label{lem:c4-unfold}
  Let $\Gamma\vdash K:T$ with $K=\bind{x}{\rho.V(\vec\alpha)}{C}$ and
  $V=\cfun{f}{\vec x{:}\vec T}{\rho}{T'}{C_V}$.
  \begin{enumerate}[label=(\roman*),nosep]
  \item\label{it:c4-unf-in} If $\rr\in\rho$, then
    $\proj{K}{\rr}=x\leftarrow\proj{V}{\rr}(\vec\alpha|_\rr)\,;\,\proj{C}{\rr}$,
    and \rn{N-Call} replaces it by $\proj{\unfold{K}}{\rr}$.
  \item\label{it:c4-unf-out} If $\rr\notin\rho$, then $\proj{\unfold{K}}{\rr}=\proj{K}{\rr}$.
  \end{enumerate}
\end{lemma}
\begin{proof}
  Write $C_V^\ast$ for $C_V[V/f][\pay{\vec\alpha}/\vec x]$.  By
  Lemma~\ref{lem:c4-subst}, $\Gamma\vdash C_V^\ast:T'$, and by
  Lemma~\ref{lem:c4-pbind},
  $\proj{\unfold{K}}{\rr}=\mbind{\proj{C_V^\ast}{\rr}}{x}{\proj{C}{\rr}}$.

  \ref{it:c4-unf-in} Since $\Gamma\vdash\alpha_i:T_i$, we have
  $\loc{\alpha_i}=\loc{T_i}$, so $\vec\alpha|_\rr$ and $\vec x|_\rr$ have the
  same length, and \rn{N-Call} fires on
  $\proj{V}{\rr}=\nfun{f}{\vec x|_\rr}{\proj{C_V}{\rr}}$, producing
  $\mbind{\proj{C_V}{\rr}[\proj{V}{\rr}/f][\vec\alpha|_\rr/\vec x|_\rr]}{x}{\proj{C}{\rr}}$.
  By Lemma~\ref{lem:c4-psubst},
  $\proj{C_V^\ast}{\rr}=\proj{C_V}{\rr}[\proj{V}{\rr}/f][\proj{\pay{\alpha_i}}{\rr}/x_i]_i$.
  A parameter $x_i$ with $\rr\notin\loc{T_i}$ is not free in
  $\proj{C_V}{\rr}$, by Lemma~\ref{lem:c4-local}, and for the others,
  $\proj{\pay{\alpha_i}}{\rr}=\proj{\alpha_i}{\rr}$.  So the two agree.

  \ref{it:c4-unf-out} The prefix of the application at $\rr$ is empty, and
  $\rr\notin\rho\supseteq\roles{C_V}=\roles{C_V^\ast}$, so by erasure
  \[
    \proj{\unfold{K}}{\rr}=\mbind{(\kw{ret}\;\bot)}{x}{\proj{C}{\rr}}=\proj{C}{\rr}=\proj{K}{\rr}.\qedhere
  \]
\end{proof}

Lemma~\ref{lem:c3-wfsubst} holds with the same proof; in part
\ref{it:c3-ws-bind}, $\unfold{K}[u/y]=\unfold{K[u/y]}$ now also uses the fact
that substitution commutes with the substitutions for $f$ and $\vec x$, since
a literal may have free variables.  Network inversion replaces item (f)
of Chapter~\ref{ch:fun} by
\begin{enumerate}[label=(\alph*),start=6,nosep]
\item if $\rr\in\roles{\ell}$ and $N(\rr)=x\leftarrow W(\vec w)\,;\,P$, then
  $W$ is a literal $\nfun{f}{\vec y}{P_W}$ with $|\vec y|=|\vec w|$, the rule is
  \rn{N-Call}, $\ell=\tauat{\rr}$, and
  $N'=\upd{N}{\rr}{\mbind{P_W[W/f][\vec w/\vec y]}{x}{P}}$.
\end{enumerate}
In the soundness proof, a network step at a participant $\rr$ of an
application $\rho.F(\vec\alpha)$ is therefore an \rn{N-Call} on
$\proj{F}{\rr}$, which is a literal only if $F$ is, so \rn{Call} applies.  For
an application whose function position is a variable, neither side can step
at $\rr$.

\begin{theorem}[EPP is lockstep]\label{thm:c4-epp}
  Suppose $\Gamma\vdash C:T$ and $\WF(C)$.
  \begin{enumerate}[label=(\arabic*),nosep]
  \item If $C\lto{\ell}C'$, then $\projn{C}\lto{\ell}\projn{C'}$.
  \item If $\projn{C}\lto{\ell}N$, then $C\lto{\ell}C'$ for some $C'$ with $N=\projn{C'}$.
  \end{enumerate}
  Consequently $\{(C,\projn{C})\mid\Gamma\vdash C:T\text{ and }\WF(C)\}$ is a
  strong bisimulation with identical labels.
\end{theorem}
\begin{proof}
  As for Theorem~\ref{thm:c3-epp} and Corollary~\ref{cor:c3-bisim}, with
  Lemma~\ref{lem:c4-unfold} in place of Lemma~\ref{lem:c3-unfold}.
\end{proof}

\begin{theorem}[Deadlock freedom]\label{thm:c4-df}
  Assume (L2), (L3), $\emptyset\vdash C:T$ and $\WF(C)$.  Every network
  reachable from $\projn{C}$ is either terminated or can step.
\end{theorem}
\begin{proof}
  As for Theorem~\ref{thm:c3-df}.  A closed application has a closed function
  value in function position, which is a literal, and its role set $\rho$
  contains the nonempty location of its result type, so \rn{Call} fires.
\end{proof}

Confluence holds on networks as in Remark~\ref{rem:c3-diamond},
under the same assumptions.

\section{Worked example}
\label{sec:c4-example}

\begin{example}[Returning a closure]\label{ex:c4-mk}
  A function $\mathit{mk}$ takes a configuration held at $\pp$ and returns a
  request handler run by $\pp$ and $\qq$.  The handler captures a value $c$
  that $\pp$ computes from the configuration.  Let
  $T_h=\farr{\pp}{\{\pp,\qq\}}{\qq}$ and
  \[
  \begin{array}{r@{\;\;}l}
    H \;=\; & \cfun{h}{\mathit{req}{:}\pp}{\{\pp,\qq\}}{\qq}{\bind{z}{\pp.(\mathit{req}\cdot c)}{\bind{w}{\com{\pp}{z}{\qq}}{\cret{\qq.w}}}}\\[0.5ex]
    M \;=\; & \cfun{\mathit{mk}}{\mathit{cfg}{:}\pp}{\{\pp,\qq\}}{T_h}{\bind{c}{\pp.(\mathit{cfg}+1)}{\cret{\{\pp,\qq\}.H}}}\\[0.5ex]
    C \;=\; & \bind{g}{\{\pp,\qq\}.M(\pp.5)}{\bind{y}{\{\pp,\qq\}.g(\pp.2)}{\cret{\qq.y}}}
  \end{array}
  \]
  The type of $M$ is valid because $\loc{T_h}=\{\pp,\qq\}$, and the result $g$
  is held at both $\pp$ and $\qq$, so both can run it.  The projections of
  the two literals are
  \[
  \begin{array}{r@{\;=\;}l@{\qquad}r@{\;=\;}l}
    \proj{H}{\pp} & \nfun{h}{\mathit{req}}{z\leftarrow\mathit{req}\cdot c\,;\,\nsend{\qq}{z}\,;\,\kw{ret}\;\bot} &
    \proj{H}{\qq} & \nfun{h}{}{w\leftarrow\nrecv{\pp}\,;\,\kw{ret}\;w}\\
    \proj{M}{\pp} & \nfun{\mathit{mk}}{\mathit{cfg}}{c\leftarrow\mathit{cfg}+1\,;\,\kw{ret}\;\proj{H}{\pp}} &
    \proj{M}{\qq} & \nfun{\mathit{mk}}{}{\kw{ret}\;\proj{H}{\qq}}
  \end{array}
  \]
  and the program projects to
  $g\leftarrow\proj{M}{\pp}(5)\,;\,y\leftarrow g(2)\,;\,\kw{ret}\;\bot$ at $\pp$
  and $g\leftarrow\proj{M}{\qq}()\,;\,y\leftarrow g()\,;\,\kw{ret}\;y$ at $\qq$.
  Role $\qq$ never sees $c$; its version of the handler comes from its own
  projection of $M$.

  Suppose $\pp$ enters $M$, $\qq$ enters $M$, and then $\qq$ enters the
  handler, all before $\pp$ computes $c$:
  \begin{align*}
    C &\lto{\tauat{\pp}} \rt{\{\qq\}}{C}{(\bind{c}{\pp.(5+1)}{\bind{y}{\{\pp,\qq\}.H(\pp.2)}{\cret{\qq.y}}})}
      && \text{\rn{Call}}\\
      &\lto{\tauat{\qq}} \bind{c}{\pp.(5+1)}{\bind{y}{\{\pp,\qq\}.H(\pp.2)}{\cret{\qq.y}}}
      && \text{\rn{Enter}}\\
      &\lto{\tauat{\qq}} \bind{c}{\pp.(5+1)}{\rt{\{\pp\}}{K}{(\bind{z}{\pp.(2\cdot c)}{\bind{w}{\com{\pp}{z}{\qq}}{\cret{\qq.w}}})}}
      && \text{\rn{Delay}, \rn{Call}}
  \end{align*}
  where $K=\bind{y}{\{\pp,\qq\}.H(\pp.2)}{\cret{\qq.y}}$.  The last step
  applies the literal $H$ while it still has $c$ free, which
  Choice~\ref{cho:c4-open} allows.  In the network, $\qq$'s program after
  entering $\mathit{mk}$ is
  $\mbind{(\kw{ret}\;\proj{H}{\qq})}{g}{(y\leftarrow g()\,;\,\kw{ret}\;y)}=y\leftarrow\proj{H}{\qq}()\,;\,\kw{ret}\;y$,
  and $\proj{H}{\qq}$ is closed, so $\qq$ can enter it.  After the
  step, $\qq$'s program is $w\leftarrow\nrecv{\pp}\,;\,\kw{ret}\;w$, which is
  the projection of the last choreography at $\qq$.  The run continues with
  $\pp$ computing $c=6$, which substitutes into both $K$ and the current state;
  $\pp$ entering the handler; $\pp$ computing $z=12$; and the communication
  $\com{\pp}{12}{\qq}$.  It ends in $\cret{\qq.12}$ after seven steps, as does
  every maximal run.
\end{example}

\begin{remark}[Named functions as literals]
  The loop of \S\ref{sec:c3-example} is written in this chapter as the literal
  \[
    \mathit{Loop}=\cfun{\mathit{loop}}{x{:}\pp,\,k{:}\qq}{\{\pp,\qq,\rr\}}{\rr}{C_{\mathit{loop}}},
  \]
  where the recursive call in $C_{\mathit{loop}}$ becomes
  $\{\pp,\qq,\rr\}.\mathit{loop}(\pp.y,\qq.k')$, an application of the bound
  self-name, and the program is
  $\bind{\mathit{res}}{\{\pp,\qq,\rr\}.\mathit{Loop}(\pp.3,\qq.0)}{\cdots}$.
  In general, a set of definitions from Chapter~\ref{ch:fun} in which every
  cycle of calls is a self-call translates by replacing each call with the
  corresponding literal.  Mutual recursion would need recursive bundles of
  literals, which unfold in the same way and have the same proofs.
\end{remark}

%======================================================================
\chapter{A Minimal Higher-Order Language}
\label{ch:min}
%======================================================================

This chapter pares the language of Chapter~\ref{ch:ho} down to a minimal
core.  Functions take a single argument and have no name, and recursion
comes back through a fixpoint.  The syntax has three sorts, values, actions
and choreographies:
\[
\begin{array}{r@{\;\;}c@{\;\;}l}
  u & ::= & x \;\mid\; v \;\mid\; \clam{x}{T}{C} \;\mid\; \cfix{f}{T}{\clam{x}{T'}{C}}\\
  A & ::= & \pp.e \;\mid\; \com{\pp}{u}{\qq} \;\mid\; \cif{\pp}{u}{C_1}{C_2} \;\mid\; \rho.u(\sigma.u')\\
  C & ::= & \cret{\rho.u} \;\mid\; \bind{x}{A}{C}
\end{array}
\]
together with the runtime term of Chapter~\ref{ch:fun}.  Its constructs are
local computation, communication, the conditional, abstraction, the
fixpoint, application, bind and return.  The distinctions that
Chapter~\ref{ch:ho} drew in the grammar, between atoms and function values,
between literals and other values, and between calls and other
choreographies, are drawn here by typing or by the well-formedness
invariant.  The proofs are those of Chapter~\ref{ch:ho} with one parameter
in place of a vector, and the new lemmas, on unique types, canonical forms
and unrolling, are a few lines each.

\section{Design choices}
\label{sec:c5-choices}

\begin{choice}[One sort of values]\label{cho:c5-values}
  Atoms and function values form a single sort $u$: a variable, a value $v$
  of the local language, an abstraction, or a fixpoint.  Typing tells data
  from functions, as it already did for variables in Chapter~\ref{ch:ho}.

  \emph{Proof saved.}  Substitution, projection and their commutation have
  one case for each form of value, instead of one family of cases for atoms
  and another for function values.  Substitution into a local expression is
  still only defined for atoms, but a well-typed substitution replaces a
  variable that occurs in a local expression by an atom
  (Lemma~\ref{lem:c5-subst}), and Chapter~\ref{ch:ho} already restricted
  itself to such substitutions (Lemma~\ref{lem:c4-psubst}).  The cost is a
  canonical-forms lemma (Lemma~\ref{lem:c5-canon}), which recovers the
  shapes of closed values that the grammar used to provide.
\end{choice}

\begin{choice}[Every value is located where it is used]\label{cho:c5-located}
  A value occurs in a choreography only as a \emph{located value} $\rho.u$:
  as the result of a return, $\cret{\rho.u}$; as the payload of a
  communication or the guard of a conditional, $\pp.u$, which stands for
  $\{\pp\}.u$; and as the function and the argument of an application,
  $\rho.u(\sigma.u')$.  Typing requires $\rho$ to be the location of the
  type of $u$.

  \emph{Proof saved.}  A located value projects, at the roles of its
  location, to the projection of its value, by one equation for data and
  functions alike.  The application $\rho.F(\vec\alpha)$ of
  Chapter~\ref{ch:ho} was already a located function applied to located
  arguments, so no annotation is added.  Since payloads and guards must have
  a role type, functions are not communicated (Choice~\ref{cho:c4-nocomm})
  by typing rather than by syntax.
\end{choice}

\begin{choice}[An abstraction has no location of its own]\label{cho:c5-lambda}
  An abstraction $\clam{x}{T}{C}$ carries no name, no role set and no result
  type.  Like a value $v$ of the local language, it gets its location from
  the located value in which it occurs, or, in a fixpoint, from the type of
  the fixpoint.  The result type is that of the body.

  \emph{Proof saved.}  Projection, unrolling and $\roles{\cdot}$ never
  consulted the role set or the result type written on a literal; only
  typing did, and typing now reads them off the context and the body.  The
  only change is that an abstraction alone, like a value $v$, has many
  types, while a located value still has exactly one
  (Lemma~\ref{lem:c5-unique}).
\end{choice}

\begin{choice}[One argument]\label{cho:c5-one}
  Every function takes exactly one located argument, and functions of
  several arguments are curried.

  \emph{Proof saved.}  Parameter and argument vectors, and their
  restrictions to a role, disappear.  At each role a projected abstraction
  has one parameter or none, a projected application has one argument or
  none, and exact types (Choice~\ref{cho:c4-exact}) make the two agree.  The
  cost is in steps, not in proofs.  Applying a curried function to $n$
  arguments takes $n$ applications, each entered by every role of $\rho$,
  because the type of a function that returns a closure is valid only if
  the closure's location is included in its own: every role that will run
  the closure takes part in creating it (\S\ref{sec:c5-derived}).
\end{choice}

\begin{choice}[Recursion by a fixpoint value]\label{cho:c5-fix}
  Abstractions are not recursive.  A recursive function is a fixpoint
  $\cfix{f}{T}{V}$, a value in which $f$ is bound in the abstraction $V$.
  The \emph{fix rule} unrolls it by substituting the fixpoint for $f$
  (Definition~\ref{def:c5-unfold}).  The body of a fixpoint is always an
  abstraction.

  \emph{Proof saved.}  Unrolling substitutes a value for a variable, so
  projection commutes with it by Lemma~\ref{lem:c5-psubst}, and the
  projection of a fixpoint unrolls in the network to the projection of its
  unrolling (Lemma~\ref{lem:c5-unroll}).  Since the body is an abstraction,
  one unrolling yields an abstraction, so every closed function value
  unrolls to an abstraction and every closed application can be entered.
  Fixpoints of arbitrary values would admit $\cfix{f}{T}{f}$, which is
  closed and well-typed but never unrolls to an abstraction, and whose
  applications would be stuck at both levels.
\end{choice}

\begin{choice}[Unrolling is part of entering]\label{cho:c5-enter}
  The fix rule is not a step.  The first role to enter an application of a
  fixpoint unrolls it and substitutes the argument in the same step,
  \rn{Call}, and in the network every participant does the same in its own
  \rn{N-Call}.

  \emph{Proof saved.}  A separate unrolling step would be taken by each
  participant on its own in the network, so the choreography would have to
  record which participants have unrolled, with a second round of per-role
  entries for every recursive call.  As part of entering, unrolling adds no
  step, and the runtime term, the well-formedness invariant, and the View and
  Entering lemmas of Chapter~\ref{ch:fun} are unchanged.
\end{choice}

\begin{choice}[The shape of a runtime term is an invariant]\label{cho:c5-rt}
  The first component of a runtime term $\rt{S}{K}{D}$ is any choreography.
  Well-formedness, not a syntactic sort, requires it to be an application of
  a literal followed by its continuation.

  \emph{Proof saved.}  No rule inspects $K$ except through $\unfold{K}$, and
  only well-formed terms are related to networks, so a sort for $K$ would
  repeat a condition that well-formedness states anyway.  Substitution maps
  a literal to a literal, so the condition is preserved.
\end{choice}

Choices~\ref{cho:c4-located}--\ref{cho:c4-nocomm} carry over: functions are
located at role sets, applications carry their role set, nothing evaluates
under an abstraction, applications are entered one role at a time, a literal
need not be closed to be applied, types are exact, and functions are not
communicated.

\section{Syntax and typing}
\label{sec:c5-syntax}

\begin{definition}[Types]\label{def:c5-types}
  Types are
  \[
    T \;::=\; \pp \;\mid\; \farr{T_1}{\rho}{T_2}
    \qquad
    \loc{\pp}=\{\pp\}
    \qquad
    \loc{\farr{T_1}{\rho}{T_2}}=\rho.
  \]
  A type is \emph{valid} if it is a role, or a function type
  $\farr{T_1}{\rho}{T_2}$ whose components are valid and whose result is
  held by participants, $\loc{T_2}\subseteq\rho$.  All types below are
  valid.
\end{definition}

\begin{definition}[Minimal choreographies]\label{def:c5-syntax}
  \[
  \begin{array}{r@{\;\;}c@{\;\;}l@{\qquad}l}
    u & ::= & x & \text{variable}\\
      & \mid & v & \text{value of the local language}\\
      & \mid & \clam{x}{T}{C} & \text{abstraction}\\
      & \mid & \cfix{f}{T}{\clam{x}{T'}{C}} & \text{fixpoint}\\[1ex]
    A & ::= & \pp.e & \text{local computation at } \pp\\
      & \mid & \com{\pp}{u}{\qq} & \text{communication from $\pp$ to $\qq$, with } \pp\neq\qq\\
      & \mid & \cif{\pp}{u}{C_1}{C_2} & \text{conditional on a guard held at } \pp\\
      & \mid & \rho.u(\alpha) & \text{application, run at } \rho\\[1ex]
    C & ::= & \cret{\alpha} & \text{return}\\
      & \mid & \bind{x}{A}{C} & \text{bind}\\
      & \mid & \rt{S}{K}{D} & \text{runtime term, } \emptyset\neq S\subseteq\Roles
  \end{array}
  \]
  Here $\alpha$ stands for a \emph{located value} $\rho.u$, with
  $\loc{\alpha}=\rho$ and $\pay{\alpha}=u$; we write $\pp.u$ for
  $\{\pp\}.u$; and $K$ and $D$ range over choreographies.  The variable $x$
  is bound in $C$ by an abstraction and by a bind, and $f$ is bound in the
  abstraction of a fixpoint.  Abstractions and fixpoints are
  \emph{literals}; $V$ ranges over abstractions.  A choreography is a
  \emph{source} choreography if it contains no runtime term, and branches
  and bodies of literals are source choreographies.
\end{definition}

Actions remain a sort of their own.  A choreography is the free monad's
normal form over them, and $\roles{\cdot}$, projection prefixes and
\rn{Delay} are all defined by cases on the action.  Substitution $C[u/x]$
replaces a variable by a value, avoiding capture, including inside the
bodies of literals.  Inside a local expression it is the substitution of the
local language, which takes an atom, so $C[u/x]$ is undefined if $u$ is a
literal and $x$ occurs free in a local expression of $C$.  Well-typed
substitutions never do this (Lemma~\ref{lem:c5-subst}).

\paragraph{Roles.}
On source terms,
\[
\begin{array}{r@{\;\;=\;\;}l@{\qquad}r@{\;\;=\;\;}l}
  \roles{\cret{\alpha}} & \loc{\alpha} &
  \roles{\bind{x}{A}{C}} & \roles{A}\cup\roles{C}\\
  \roles{\pp.e} & \{\pp\} &
  \roles{\com{\pp}{u}{\qq}} & \{\pp,\qq\}\\
  \roles{\rho.u(\alpha)} & \rho &
  \roles{\cif{\pp}{u}{C_1}{C_2}} & \{\pp\}\cup\roles{C_1}\cup\roles{C_2}
\end{array}
\]
and $\rcv{\cif{\pp}{u}{C_1}{C_2}}=(\roles{C_1}\cup\roles{C_2})\setminus\{\pp\}$.
Roles never look inside a value, so $\roles{C[u/x]}=\roles{C}$.

\paragraph{Typing.}
A context $\Gamma$ maps variables to types, and $\Gamma\vdash_\pp e$ means
$\Gamma(x)=\pp$ for every $x\in\fv(e)$.  The judgements are $\Gamma\vdash
u:T$ for values, $\Gamma\vdash\alpha:T$ for located values, and $\Gamma\vdash
A:T$ and $\Gamma\vdash C:T$ for actions and choreographies, given in
Figure~\ref{fig:c5-typing}.  A program is a choreography $C$ with
$\emptyset\vdash C:T$.

\begin{figure}[t]
\begin{mathpar}
  \inferrule{ }{\Gamma\vdash x:\Gamma(x)}
  \and
  \inferrule{ }{\Gamma\vdash v:\pp}
  \and
  \inferrule{\Gamma,x{:}T_1\vdash C:T_2 \\ \roles{C}\subseteq\rho}
            {\Gamma\vdash\clam{x}{T_1}{C}:\farr{T_1}{\rho}{T_2}}
  \and
  \inferrule{\Gamma,f{:}T\vdash V:T}
            {\Gamma\vdash\cfix{f}{T}{V}:T}
  \\
  \inferrule{\Gamma\vdash u:T \\ \loc{T}=\rho}{\Gamma\vdash\rho.u:T}
  \and
  \inferrule{\Gamma\vdash\alpha:T}{\Gamma\vdash\cret{\alpha}:T}
  \and
  \inferrule{\Gamma\vdash A:T' \\ \Gamma,x{:}T'\vdash C:T}{\Gamma\vdash\bind{x}{A}{C}:T}
  \\
  \inferrule{\Gamma\vdash_\pp e}{\Gamma\vdash\pp.e:\pp}
  \and
  \inferrule{\pp\neq\qq \\ \Gamma\vdash\pp.u:\pp}{\Gamma\vdash\com{\pp}{u}{\qq}:\qq}
  \and
  \inferrule{\Gamma\vdash\pp.u:\pp \\ \Gamma\vdash C_1:T \\ \Gamma\vdash C_2:T}
            {\Gamma\vdash\cif{\pp}{u}{C_1}{C_2}:T}
  \\
  \inferrule{\Gamma\vdash\rho.u:\farr{T_1}{\rho}{T_2} \\ \Gamma\vdash\alpha:T_1}
            {\Gamma\vdash\rho.u(\alpha):T_2}
  \and
  \inferrule{\Gamma\vdash K:T \\ \Gamma\vdash D:T}{\Gamma\vdash\rt{S}{K}{D}:T}
\end{mathpar}
\caption{Typing of minimal choreographies.}
\label{fig:c5-typing}
\end{figure}

A value of the local language has every role type, and an abstraction has a
function type for every $\rho$ that covers the roles of its body; the
located value around them selects one.  So $\Gamma\vdash\pp.u:\pp$ holds
exactly when $u$ is an atom with $\Gamma\vdash_\pp u$, and the payloads and
guards of Chapter~\ref{ch:cond} are recovered by typing.

\begin{lemma}[Unique types]\label{lem:c5-unique}
  A located value, an action or a choreography has at most one type under
  a given $\Gamma$.
\end{lemma}
\begin{proof}
  By induction.  The type of $\rho.v$ is the role $\pp$ with
  $\{\pp\}=\rho$, and that of $\rho.x$ is $\Gamma(x)$.  The type of
  $\rho.\clam{x}{T_1}{C}$ is $\farr{T_1}{\rho}{T_2}$, where $T_2$ is the
  unique type of $C$, and that of $\rho.\cfix{f}{T}{V}$ is $T$.  The other
  cases are as in Chapter~\ref{ch:cond}.
\end{proof}

As before, $\loc{A}$ denotes the location of the type of $A$.

\begin{lemma}[Results are held by participants]\label{lem:c5-locroles}
  If $\Gamma\vdash A:T$ then $\loc{T}\subseteq\roles{A}$, and if $C$ is a
  source choreography with $\Gamma\vdash C:T$, then
  $\loc{T}\subseteq\roles{C}$.
\end{lemma}
\begin{proof}
  As for Lemma~\ref{lem:c4-locroles}.  For a return,
  $\loc{T}=\loc{\alpha}=\roles{\cret{\alpha}}$, and for an application,
  $\loc{T_2}\subseteq\rho$ because the function's type is valid.
\end{proof}

\begin{lemma}[Canonical forms]\label{lem:c5-canon}
  Let $u$ be closed with $\Gamma\vdash u:T$.  If $T$ is a role, then $u$ is
  a value of the local language.  If $T$ is a function type, then $u$ is a
  literal.
\end{lemma}
\begin{proof}
  A closed value is not a variable.  Values of the local language have only
  role types, and literals only function types.
\end{proof}

\section{Choreographic semantics}
\label{sec:c5-sem}

\begin{definition}[Unrolling and unfolding]\label{def:c5-unfold}
  Unrolling maps a literal to an abstraction:
  \[
    \unroll{V} = V
    \qquad
    \unroll{\cfix{f}{T}{V}} = V[\cfix{f}{T}{V}/f].
  \]
  The second equation is the \emph{fix rule}.  For
  $K=\bind{x}{\rho.u(\alpha)}{C}$ with $u$ a literal and
  $\unroll{u}=\clam{y}{T}{C_u}$, its bound variables renamed apart,
  \[
    \unfold{K} \;=\; \mbind{C_u[\pay{\alpha}/y]}{x}{C}.
  \]
\end{definition}

Derived bind is that of Definition~\ref{def:c1-bind}, with
$\mbind{(\cret{\alpha})}{x}{C}=C[\pay{\alpha}/x]$.  Labels are those of
Chapter~\ref{ch:cond}:
\[
  \ell \;::=\; \tauat{\pp} \;\mid\; \com{\pp}{v}{\qq} \;\mid\; \castlbl{\pp}{v}{Q}.
\]

\begin{figure}[t]
\begin{mathpar}
  \inferrule*[right=\rn{Local}]
    {e\Downarrow v}
    {\bind{x}{\pp.e}{C} \lto{\tauat{\pp}} C[v/x]}
  \and
  \inferrule*[right=\rn{Com}]
    { }
    {\bind{x}{\com{\pp}{v}{\qq}}{C} \lto{\com{\pp}{v}{\qq}} C[v/x]}
  \\
  \inferrule*[right=\rn{Cast}]
    {v\in\{\ttrue,\ffalse\} \\ A=\cif{\pp}{v}{C_\ttrue}{C_\ffalse} \\ Q=\rcv{A}}
    {\bind{x}{A}{C} \lto{\castlbl{\pp}{v}{Q}} \mbind{C_v}{x}{C}}
  \and
  \inferrule*[right=\rn{Delay}]
    {C\lto{\ell}C' \\ \roles{\ell}\cap\roles{A}=\emptyset}
    {\bind{x}{A}{C} \lto{\ell} \bind{x}{A}{C'}}
  \\
  \inferrule*[right=\rn{Call}]
    {K=\bind{x}{\rho.u(\alpha)}{C} \\ u \text{ a literal} \\ \rr\in\rho}
    {K \lto{\tauat{\rr}} \rt{\rho\setminus\{\rr\}}{K}{\unfold{K}}}
  \and
  \inferrule*[right=\rn{Enter}]
    {\rr\in S}
    {\rt{S}{K}{D} \lto{\tauat{\rr}} \rt{S\setminus\{\rr\}}{K}{D}}
  \and
  \inferrule*[right=\rn{Inside}]
    {D\lto{\ell}D' \\ S\cap\roles{\ell}=\emptyset}
    {\rt{S}{K}{D} \lto{\ell} \rt{S}{K}{D'}}
\end{mathpar}
\caption{Semantics of minimal choreographies.  Throughout, $v$ is a value of
  the local language, and a target $\rt{\emptyset}{K}{D}$ stands for $D$.}
\label{fig:c5-sem}
\end{figure}

The semantics is given in Figure~\ref{fig:c5-sem}.  These are the rules of
Chapters~\ref{ch:core} to~\ref{ch:fun}, with \rn{Call} as in
Chapter~\ref{ch:ho}.  \rn{Com} fires only when the payload is a value of the
local language, \rn{Cast} only when the guard is a boolean, and \rn{Call}
only when the function is a literal.  For well-typed terms, a closed payload
or function always has the required shape (Lemma~\ref{lem:c5-canon}).  A
variable in function position is handled as in \S\ref{sec:c4-sem}: it is
bound by an enclosing literal, or by a pending bind whose action involves
every role of the application.

\section{Networks}
\label{sec:c5-net}

\begin{definition}[Minimal networks]\label{def:c5-net}
  Local values, programs and prefixes are
  \[
  \begin{array}{r@{\;\;}c@{\;\;}l@{\qquad}l}
    w & ::= & x \;\mid\; v \;\mid\; \nlam{\vec y}{P} \;\mid\; \nfix{f}{\nlam{\vec y}{P}} & \text{with } |\vec y|\le 1\\
    P & ::= & \kw{ret}\;w \;\mid\; \kw{ret}\;\bot \;\mid\; \pi\,P\\
    \pi & ::= & x\leftarrow e\,; \;\mid\; \nsend{\qq}{w}\,; \;\mid\; x\leftarrow\nrecv{\pp}\,;\\
        & \mid & x\leftarrow\sendbr{Q}{w}{P_1}{P_2}\,; \;\mid\; x\leftarrow\recvbr{\pp}{P_1}{P_2}\,;\\
        & \mid & x\leftarrow w(\vec w)\,; & \text{with } |\vec w|\le 1
  \end{array}
  \]
  As in choreographies, the last two forms of $w$ are \emph{literals}, $W$
  ranges over local abstractions, $\unroll{W}=W$, and
  $\unroll{\nfix{f}{W}}=W[\nfix{f}{W}/f]$.  Network bind is
  \[
  \mbind{(\kw{ret}\;w)}{x}{P} = P[w/x]
  \qquad
  \mbind{(\kw{ret}\;\bot)}{x}{P} = P
  \qquad
  \mbind{(\pi\,P_1)}{x}{P_2} = \pi\,(\mbind{P_1}{x}{P_2})
  \]
  where, in the last equation, the variable bound by $\pi$ is renamed apart
  from $\fv(P_2)\cup\{x\}$.
\end{definition}

\begin{figure}[t]
\begin{mathpar}
  \inferrule*[right=\rn{N-Local}]
    {N(\pp) = x\leftarrow e\,;\,P \\ e\Downarrow v}
    {N \lto{\tauat{\pp}} \upd{N}{\pp}{P[v/x]}}
  \and
  \inferrule*[right=\rn{N-Com}]
    {N(\pp) = \nsend{\qq}{v}\,;\,P \\ N(\qq) = x\leftarrow\nrecv{\pp}\,;\,P'}
    {N \lto{\com{\pp}{v}{\qq}} N[\pp\mapsto P,\ \qq\mapsto P'[v/x]]}
  \\
  \inferrule*[right=\rn{N-Cast}]
    {N(\pp) = x\leftarrow\sendbr{Q}{v}{P_\ttrue}{P_\ffalse}\,;\,P \\
     v\in\{\ttrue,\ffalse\} \\\\
     N(\qq) = x_\qq\leftarrow\recvbr{\pp}{P^\qq_\ttrue}{P^\qq_\ffalse}\,;\,P^\qq \text{ for every } \qq\in Q}
    {N \lto{\castlbl{\pp}{v}{Q}}
     N[\pp\mapsto\mbind{P_v}{x}{P}][\qq\mapsto\mbind{P^\qq_v}{x_\qq}{P^\qq}]_{\qq\in Q}}
  \\
  \inferrule*[right=\rn{N-Call}]
    {N(\rr)=x\leftarrow w(\vec w)\,;\,P \\ w \text{ a literal} \\ \unroll{w}=\nlam{\vec y}{P_w} \\ |\vec y|=|\vec w|}
    {N \lto{\tauat{\rr}} \upd{N}{\rr}{\mbind{P_w[\vec w/\vec y]}{x}{P}}}
\end{mathpar}
\caption{Semantics of minimal networks.  Throughout, $v$ is a value of the
  local language.}
\label{fig:c5-net}
\end{figure}

The network semantics is given in Figure~\ref{fig:c5-net}.  As in the
choreography, the fix rule is applied by the role that enters, in the same
step.

\section{Endpoint projection}
\label{sec:c5-epp}

\begin{definition}[Projection]\label{def:c5-epp}
  Values project to local values:
  \[
    \proj{x}{\rr}=x
    \qquad
    \proj{v}{\rr}=v
    \qquad
    \proj{\clam{x}{T}{C}}{\rr}=\nlam{x|_\rr}{\proj{C}{\rr}}
    \qquad
    \proj{\cfix{f}{T}{V}}{\rr}=\nfix{f}{\proj{V}{\rr}}
  \]
  where $x|_\rr$ is $x$ if $\rr\in\loc{T}$ and empty otherwise.  A located
  value is projected only at the roles of its location, as
  $\proj{\rho.u}{\rr}=\proj{u}{\rr}$, and $\alpha|_\rr$ is
  $\proj{\alpha}{\rr}$ if $\rr\in\loc{\alpha}$ and empty otherwise.
  Choreographies project as
  \[
  \proj{\cret{\alpha}}{\rr} =
    \begin{cases}
      \kw{ret}\;\proj{\alpha}{\rr} & \rr\in\loc{\alpha}\\
      \kw{ret}\;\bot & \text{otherwise}
    \end{cases}
  \qquad
  \qquad
  \proj{\rt{S}{K}{D}}{\rr} =
    \begin{cases}
      \proj{K}{\rr} & \rr\in S\\
      \proj{D}{\rr} & \text{otherwise}
    \end{cases}
  \]
  and $\proj{\bind{x}{A}{C}}{\rr} = \pre{A}{x}{\rr}\,\proj{C}{\rr}$, where
  the prefix of an action is $\varepsilon$ at every role not listed
  below:
  \[
  \begin{array}{r@{\;\;=\;\;}l@{\qquad}l}
    \pre{\pp.e}{x}{\pp} & x\leftarrow e\,;\\
    \pre{\com{\pp}{u}{\qq}}{x}{\pp} & \nsend{\qq}{\proj{u}{\pp}}\,;\\
    \pre{\com{\pp}{u}{\qq}}{x}{\qq} & x\leftarrow\nrecv{\pp}\,;\\
    \pre{A}{x}{\pp} & x\leftarrow\sendbr{Q}{\proj{u}{\pp}}{\proj{C_1}{\pp}}{\proj{C_2}{\pp}}\,;
      & A=\cif{\pp}{u}{C_1}{C_2},\ Q=\rcv{A}\\
    \pre{A}{x}{\rr} & x\leftarrow\recvbr{\pp}{\proj{C_1}{\rr}}{\proj{C_2}{\rr}}\,;
      & \text{the same } A,\ \rr\in Q\\
    \pre{\rho.u(\alpha)}{x}{\rr} & x\leftarrow\proj{u}{\rr}(\alpha|_\rr)\,; & \rr\in\rho.
  \end{array}
  \]
\end{definition}

Lemma~\ref{lem:c1-prefix} holds in the form of Chapter~\ref{ch:ho}: if
$\pre{A}{x}{\rr}$ does not bind $x$, then $\rr\notin\loc{A}$.  Projection
preserves the shape of a value in both directions: $\proj{u}{\rr}$ is a value
of the local language exactly when $u$ is, in which case
$\proj{u}{\rr}=u$, and it is a literal exactly when $u$ is.

\begin{definition}[Well-formedness]\label{def:c5-wf}
  $\WF(\cret{\alpha})$ holds, $\WF(\bind{x}{A}{C})$ holds if $\WF(C)$ does,
  and $\WF(\rt{S}{K}{D})$ holds if $K=\bind{x}{\rho.u(\alpha)}{C}$ with $u$
  a literal (the \emph{shape condition}) and
  \begin{enumerate}[label=(\roman*),nosep]
  \item $\WF(K)$ and $\WF(D)$,
  \item $S\subseteq\rho$, and
  \item $\proj{D}{\qq}=\proj{\unfold{K}}{\qq}$ for every $\qq\in S$.
  \end{enumerate}
\end{definition}

Conditions (i)--(iii) are those of Definition~\ref{def:c3-wf}, and the shape
condition replaces the syntactic sort of calls.  Every source choreography
is well-formed.

\section{Correctness}
\label{sec:c5-correct}

The development of Chapter~\ref{ch:ho} carries over, with a single
parameter in place of a vector.  We give the lemmas whose statements or
proofs change, and the new lemma on unrolling.

\begin{lemma}[Substitution]\label{lem:c5-subst}
  If $\Gamma,x{:}T\vdash C:T'$ and $\Gamma\vdash u:T$, then $C[u/x]$ is
  defined and $\Gamma\vdash C[u/x]:T'$.  The same holds for values, located
  values and actions.
\end{lemma}
\begin{proof}
  By induction on the typing derivation.  If $x$ is free in a local
  expression $e$, checked as $\Gamma,x{:}T\vdash_\pp e$, then $T=\pp$, so $u$
  is a value of the local language or a variable held at $\pp$, and
  $\Gamma\vdash_\pp e[u/x]$ as in Lemma~\ref{lem:c1-subst}.  An occurrence
  of $x$ as a value is replaced by the derivation of $\Gamma\vdash u:T$,
  after weakening, and the located value around it keeps its type.  The side
  condition $\roles{C}\subseteq\rho$ of the abstraction rule is unaffected,
  because roles are stable under substitution.
\end{proof}

\begin{lemma}[Projection is local]\label{lem:c5-local}
  If $\Gamma\vdash C:T$, then every free variable $x$ of $\proj{C}{\rr}$
  satisfies $\rr\in\loc{\Gamma(x)}$.  The same holds for $\proj{u}{\rr}$
  when $\Gamma\vdash u:T$ and $\rr\in\loc{T}$.
\end{lemma}
\begin{proof}
  As for Lemma~\ref{lem:c4-local}, with these value cases.  A variable has
  type $\Gamma(x)$, and a value $v$ is closed.  For an abstraction
  $\clam{x}{T_1}{C}$, the free variables of $\proj{C}{\rr}$ are held at
  $\rr$ under $\Gamma,x{:}T_1$ by the induction hypothesis, and $x$ is one of
  them only if $\rr\in\loc{T_1}$, when the projection binds it.  For a
  fixpoint $\cfix{f}{T}{V}$, $\rr\in\loc{T}$, so the induction hypothesis
  applies to $V$, and $f$ is bound by the projection.  A located value is
  projected only at roles of its location, which is the location of its
  type.
\end{proof}

\begin{lemma}[Projection commutes with substitution]\label{lem:c5-psubst}
  If $C[u/x]$ is defined, then
  $\proj{C[u/x]}{\rr}=\proj{C}{\rr}[\proj{u}{\rr}/x]$, and likewise for
  values.  Hence, if $\Gamma,x{:}T\vdash C:T'$, $\Gamma\vdash u:T$ and
  $\rr\notin\loc{T}$, then $\proj{C[u/x]}{\rr}=\proj{C}{\rr}$.
\end{lemma}
\begin{proof}
  By induction.  Every choice that projection makes depends on roles,
  receivers, locations of located values, or parameter types, and
  substitution changes none of them.  Where $x$ is free in a local
  expression, $u$ is an atom, and $\proj{u}{\rr}=u$.  The second claim
  follows by Lemma~\ref{lem:c5-local}.
\end{proof}

Erasure (Lemma~\ref{lem:c2-erase}), typing of bind
(Lemma~\ref{lem:c2-tbind}) and projection commutes with bind
(Lemma~\ref{lem:c4-pbind}) hold with the same proofs.  In the return case of
the last, $\proj{\alpha}{\rr}=\proj{\pay{\alpha}}{\rr}$ by definition.

\begin{lemma}[Unrolling]\label{lem:c5-unroll}
  Let $u$ be a literal.
  \begin{enumerate}[label=(\roman*),nosep]
  \item\label{it:c5-unroll-type} If $\Gamma\vdash u:T$ then $\Gamma\vdash\unroll{u}:T$.
  \item\label{it:c5-unroll-proj} $\proj{\unroll{u}}{\rr}=\unroll{\proj{u}{\rr}}$ for every role $\rr$.
  \item\label{it:c5-unroll-subst} $\unroll{u}[u'/y]=\unroll{u[u'/y]}$.
  \end{enumerate}
\end{lemma}
\begin{proof}
  All three are immediate for an abstraction.  Let $u=\cfix{f}{T}{V}$.
  \ref{it:c5-unroll-type} From $\Gamma,f{:}T\vdash V:T$ and
  $\Gamma\vdash u:T$, Lemma~\ref{lem:c5-subst} gives $\Gamma\vdash V[u/f]:T$.
  \ref{it:c5-unroll-proj} By Lemma~\ref{lem:c5-psubst},
  $\proj{V[u/f]}{\rr}=\proj{V}{\rr}[\nfix{f}{\proj{V}{\rr}}/f]$, which is
  $\unroll{\nfix{f}{\proj{V}{\rr}}}$.
  \ref{it:c5-unroll-subst} The two substitutions commute, since $f$ is bound
  in $u$, so it is distinct from $y$ and not free in $u'$.
\end{proof}

\begin{lemma}[Unfolding]\label{lem:c5-unfold}
  Let $\Gamma\vdash K:T$ with $K=\bind{x}{\rho.u(\alpha)}{C}$ and $u$ a
  literal.
  \begin{enumerate}[label=(\roman*),nosep]
  \item\label{it:c5-unf-in} If $\rr\in\rho$, then
    $\proj{K}{\rr}=x\leftarrow\proj{u}{\rr}(\alpha|_\rr)\,;\,\proj{C}{\rr}$,
    and \rn{N-Call} replaces it by $\proj{\unfold{K}}{\rr}$.
  \item\label{it:c5-unf-out} If $\rr\notin\rho$, then $\proj{\unfold{K}}{\rr}=\proj{K}{\rr}$.
  \end{enumerate}
\end{lemma}
\begin{proof}
  By typing, $\Gamma\vdash u:\farr{T_1}{\rho}{T_2}$ and
  $\Gamma\vdash\alpha:T_1$.  Let $\unroll{u}=\clam{y}{T_1}{C_u}$ and
  $C_u^\ast=C_u[\pay{\alpha}/y]$.  By
  Lemma~\ref{lem:c5-unroll}\ref{it:c5-unroll-type} and the abstraction
  rule, $\Gamma,y{:}T_1\vdash C_u:T_2$ and $\roles{C_u}\subseteq\rho$, so
  $\Gamma\vdash C_u^\ast:T_2$ by Lemma~\ref{lem:c5-subst}, and by
  Lemma~\ref{lem:c4-pbind},
  $\proj{\unfold{K}}{\rr}=\mbind{\proj{C_u^\ast}{\rr}}{x}{\proj{C}{\rr}}$.

  \ref{it:c5-unf-in} By Lemma~\ref{lem:c5-unroll}\ref{it:c5-unroll-proj},
  $\proj{u}{\rr}$ is a literal that unrolls to
  $\nlam{y|_\rr}{\proj{C_u}{\rr}}$.  Since $\loc{\alpha}=\loc{T_1}$, the
  parameter $y|_\rr$ and the argument $\alpha|_\rr$ are both present or both
  absent, so \rn{N-Call} fires and produces
  $\mbind{\proj{C_u}{\rr}[\alpha|_\rr/y|_\rr]}{x}{\proj{C}{\rr}}$.  If
  $\rr\in\loc{T_1}$, then
  $\proj{C_u}{\rr}[\proj{\alpha}{\rr}/y]=\proj{C_u^\ast}{\rr}$ by
  Lemma~\ref{lem:c5-psubst}.  Otherwise
  $\proj{C_u}{\rr}=\proj{C_u^\ast}{\rr}$ by the second claim of that lemma.

  \ref{it:c5-unf-out} The prefix of the application at $\rr$ is empty, and
  $\rr\notin\rho\supseteq\roles{C_u}=\roles{C_u^\ast}$, so by erasure
  \[
    \proj{\unfold{K}}{\rr}=\mbind{(\kw{ret}\;\bot)}{x}{\proj{C}{\rr}}=\proj{C}{\rr}=\proj{K}{\rr}.\qedhere
  \]
\end{proof}

\paragraph{The rest of the development.}
\begin{itemize}
\item \emph{Subject reduction}: as before, with $\Gamma\vdash\unfold{K}:T$
  for \rn{Call} as in the proof of Lemma~\ref{lem:c5-unfold}.
\item \emph{Substitution and well-formedness}
  (Lemma~\ref{lem:c3-wfsubst}): in part~\ref{it:c3-ws-bind}, the
  commutation of unfolding with substitution also uses
  Lemma~\ref{lem:c5-unroll}\ref{it:c5-unroll-subst}; in
  part~\ref{it:c3-ws-subst}, the shape condition is preserved because
  substitution maps a literal to a literal.
\item \emph{View}, \emph{Entering} and \emph{preservation of
  well-formedness} (Lemmas~\ref{lem:c3-view}, \ref{lem:c3-entering}
  and~\ref{lem:c3-wfpres}) hold with the same proofs, using
  Lemma~\ref{lem:c5-unfold} in place of Lemma~\ref{lem:c3-unfold}.
  \rn{Call} establishes the shape condition, and \rn{Enter} and \rn{Inside}
  leave $K$ unchanged.
\item \emph{Frames} (Lemmas~\ref{lem:c1-cframe} and~\ref{lem:c1-nframe}):
  as in Chapter~\ref{ch:fun}.
\item \emph{Network inversion} (Lemma~\ref{lem:c1-inv}): items (a)--(e) are
  as before, where a send or a send-and-branch fires only if its local value
  is a value $v$ of the local language.  Item (f) becomes: if
  $N(\rr)=x\leftarrow w(\vec w)\,;\,P$ and $\rr\in\roles{\ell}$, then $w$ is
  a literal, $\unroll{w}=\nlam{\vec y}{P_w}$ with $|\vec y|=|\vec w|$, the
  rule is \rn{N-Call}, $\ell=\tauat{\rr}$, and
  $N'=\upd{N}{\rr}{\mbind{P_w[\vec w/\vec y]}{x}{P}}$.
\end{itemize}
In the soundness proof, a step at the sender of $\com{\pp}{u}{\qq}$ or at
the guard holder of a conditional forces $\proj{u}{\pp}$ to be a value of
the local language, so $u$ is that value; and a step at a participant $\rr$
of $\rho.u(\alpha)$ forces $\proj{u}{\rr}$ to be a literal, so $u$ is a
literal and \rn{Call} applies.

\begin{theorem}[EPP is lockstep]\label{thm:c5-epp}
  Suppose $\Gamma\vdash C:T$ and $\WF(C)$.
  \begin{enumerate}[label=(\arabic*),nosep]
  \item If $C\lto{\ell}C'$, then $\projn{C}\lto{\ell}\projn{C'}$.
  \item If $\projn{C}\lto{\ell}N$, then $C\lto{\ell}C'$ for some $C'$ with $N=\projn{C'}$.
  \end{enumerate}
  Consequently $\{(C,\projn{C})\mid\Gamma\vdash C:T\text{ and }\WF(C)\}$ is a
  strong bisimulation with identical labels.
\end{theorem}
\begin{proof}
  As for Theorem~\ref{thm:c4-epp}, with Lemma~\ref{lem:c5-unfold} in place
  of Lemma~\ref{lem:c4-unfold}.
\end{proof}

\begin{theorem}[Deadlock freedom]\label{thm:c5-df}
  Assume (L2), (L3), $\emptyset\vdash C:T$ and $\WF(C)$.  Every network
  reachable from $\projn{C}$ is either terminated or can step.
\end{theorem}
\begin{proof}
  As for Theorem~\ref{thm:c4-df}.  The payload of a closed head
  communication and the guard of a closed head conditional have a role type,
  so by Lemma~\ref{lem:c5-canon} they are values of the local language, and
  the guard is a boolean by (L3).  The function of a closed head
  application has a function type, so it is a literal, and its role set
  contains the nonempty location of the result type, so \rn{Call} fires.
\end{proof}

Confluence holds on networks as in Remark~\ref{rem:c3-diamond}, under the
same assumptions.

\section{Derived forms}
\label{sec:c5-derived}

The constructs left out of this chapter can be written in terms of the
remaining ones.

\begin{itemize}
\item \emph{Several arguments.}  A literal
  $\cfun{f}{x_1{:}T_1,\ldots,x_n{:}T_n}{\rho}{T}{C}$ of
  Chapter~\ref{ch:ho}, with $n\ge 1$, becomes the fixpoint
  \[
    \cfix{f}{T^\ast}{\clam{x_1}{T_1}{\cret{\rho.\big(\clam{x_2}{T_2}{\cdots\cret{\rho.(\clam{x_n}{T_n}{C^\ast})}}\big)}}}
  \]
  of type
  $T^\ast=\farr{T_1}{\rho}{\farr{T_2}{\rho}{\cdots\farr{T_n}{\rho}{T}}}$,
  and an application $\bind{x}{\rho.F(\alpha_1,\ldots,\alpha_n)}{C}$
  becomes
  $\bind{g_1}{\rho.F^\ast(\alpha_1)}{\cdots\bind{g_{n-1}}{\rho.g_{n-2}(\alpha_{n-1})}{\bind{x}{\rho.g_{n-1}(\alpha_n)}{C^\ast}}}$,
  where ${}^\ast$ translates nested literals and applications in the same
  way.  Every intermediate type has location $\rho$, so all are valid.
  The translation is not step for step: each of the $n$ applications is
  entered by every role of $\rho$, so a call takes $n\cdot|\rho|$ entries
  instead of $|\rho|$.
\item \emph{No arguments.}  A function of no arguments is an abstraction
  $\clam{z}{\pp}{C}$ of type $\farr{\pp}{\rho}{T}$ with an unused parameter,
  applied to a dummy argument $\pp.v$.  If $\pp\notin\rho$, no role keeps
  the parameter or the argument,
  and every role's local function has no parameter, as in
  Chapter~\ref{ch:fun}.
\item \emph{Naming a value.}  $\kw{let}\;x=\alpha\;\kw{in}\;C$ stands for
  $C[\pay{\alpha}/x]=\mbind{(\cret{\alpha})}{x}{C}$ and takes no step.  No
  action can do this for a function held at several roles: it would be one
  step of all of them with no network counterpart.
\item \emph{Returning an action.}  $\cret{A}$ stands for
  $\bind{x}{A}{\cret{\loc{A}.x}}$.  It adds no step, because derived bind
  has a right unit: $\mbind{C}{x}{\cret{\loc{T}.x}}=C$ for a source
  choreography $C$ with $\Gamma\vdash C:T$, by induction on its spine, since
  a return $\cret{\alpha}$ of type $T$ has $\loc{\alpha}=\loc{T}$.  For
  instance, \rn{Cast} takes
  $\bind{x}{\cif{\pp}{v}{C_\ttrue}{C_\ffalse}}{\cret{\loc{A}.x}}$ directly
  to $C_v$.
\end{itemize}

\section{Worked example}
\label{sec:c5-example}

The loop of \S\ref{sec:c3-example}, curried.  Let $\rho=\{\pp,\qq,\rr\}$,
$T_k=\farr{\qq}{\rho}{\rr}$ and $T_L=\farr{\pp}{\rho}{T_k}$, and
\[
\begin{array}{r@{\;\;}l}
  \mathit{Loop} \;=\; & \cfix{\mathit{loop}}{T_L}{\clam{x}{\pp}{\cret{\rho.(\clam{k}{\qq}{B(x,k)})}}}\\[1ex]
  B(x,k) \;=\; & \bind{y}{\pp.(2\cdot x)}{\bind{b}{\pp.(y>10)}{\bind{z}{\cif{\pp}{b}{T(y)}{E(y,k)}}{\cret{\rr.z}}}}\\
  T(y) \;=\; & \bind{w}{\com{\pp}{y}{\qq}}{\bind{u}{\com{\qq}{w}{\rr}}{\cret{\rr.u}}}\\
  E(y,k) \;=\; & \bind{k'}{\qq.(k+1)}{\bind{g}{\rho.\mathit{loop}(\pp.y)}{\bind{v}{\rho.g(\qq.k')}{\cret{\rr.v}}}}
\end{array}
\]
The inner abstraction mentions $x$, which is held at $\pp$ and substituted
when the outer one is entered (Choice~\ref{cho:c4-open}).  Its type $T_k$
has location $\rho$, so $T_L$ is valid and the outer function returns a
closure that all three roles can run.  The program is
\[
  C \;=\; \bind{h}{\rho.\mathit{Loop}(\pp.3)}{\bind{\mathit{res}}{\rho.h(\qq.0)}{\bind{m}{\rr.(\mathit{res}+1)}{\cret{\rr.m}}}}.
\]
The projections of $\mathit{Loop}$ are
\[
\begin{array}{r@{\;\;}l}
  \pp: & \nfix{\mathit{loop}}{\nlam{x}{\kw{ret}\;\nlam{}{y\leftarrow 2\cdot x\,;\;b\leftarrow y>10\,;}}}\\
       & \quad z\leftarrow\sendbr{\{\qq,\rr\}}{b}{(\nsend{\qq}{y}\,;\,\kw{ret}\;\bot)}{(g\leftarrow\mathit{loop}(y)\,;\,v\leftarrow g()\,;\,\kw{ret}\;\bot)}\,;\;\kw{ret}\;\bot\\[1ex]
  \qq: & \nfix{\mathit{loop}}{\nlam{}{\kw{ret}\;\nlam{k}{z\leftarrow\pp?\big((w\leftarrow\nrecv{\pp}\,;\,\nsend{\rr}{w}\,;\,\kw{ret}\;\bot),}}}\\
       & \quad (k'\leftarrow k+1\,;\,g\leftarrow\mathit{loop}()\,;\,v\leftarrow g(k')\,;\,\kw{ret}\;\bot)\big)\,;\;\kw{ret}\;\bot\\[1ex]
  \rr: & \nfix{\mathit{loop}}{\nlam{}{\kw{ret}\;\nlam{}{z\leftarrow\recvbr{\pp}{(u\leftarrow\nrecv{\qq}\,;\,\kw{ret}\;u)}{(g\leftarrow\mathit{loop}()\,;\,v\leftarrow g()\,;\,\kw{ret}\;v)}\,;\;\kw{ret}\;z}}}
\end{array}
\]
and the program projects to
$h\leftarrow\proj{\mathit{Loop}}{\pp}(3)\,;\,\mathit{res}\leftarrow h()\,;\,\kw{ret}\;\bot$
at $\pp$,
$h\leftarrow\proj{\mathit{Loop}}{\qq}()\,;\,\mathit{res}\leftarrow h(0)\,;\,\kw{ret}\;\bot$
at $\qq$, and
$h\leftarrow\proj{\mathit{Loop}}{\rr}()\,;\,\mathit{res}\leftarrow h()\,;\,m\leftarrow\mathit{res}+1\,;\,\kw{ret}\;m$
at $\rr$.  Each role keeps exactly the parameters it holds: $x$ at $\pp$,
$k$ at $\qq$, and none at $\rr$.

Write $H=\clam{k}{\qq}{B(3,k)[\mathit{Loop}/\mathit{loop}]}$ and
$R(z)=\bind{m}{\rr.(z+1)}{\cret{\rr.m}}$.  One run begins
\begin{align*}
  C &\lto{\tauat{\pp}} \rt{\{\qq,\rr\}}{C}{D_1}
    && \text{\rn{Call}, unrolling } \mathit{Loop}\\
  &\lto{\tauat{\pp}} \rt{\{\qq,\rr\}}{C}{\rt{\{\qq,\rr\}}{D_1}{D_2}}
    && \text{\rn{Inside}, \rn{Call}}\\
  &\lto{\tauat{\qq}}\;\lto{\tauat{\qq}} \rt{\{\rr\}}{C}{\rt{\{\rr\}}{D_1}{D_2}}
    && \text{\rn{Enter}; \rn{Inside}, \rn{Enter}}\\
  &\lto{\tauat{\rr}}\;\lto{\tauat{\rr}} D_2
    && \text{\rn{Enter}, twice}
\end{align*}
where
\[
\begin{array}{r@{\;\;}l}
  D_1 \;=\; \unfold{C} \;=\; & \bind{\mathit{res}}{\rho.H(\qq.0)}{R(\mathit{res})}\\
  D_2 \;=\; \unfold{D_1} \;=\; & \bind{y}{\pp.(2\cdot 3)}{\bind{b}{\pp.(y>10)}{}}\\
    & \bind{z}{\cif{\pp}{b}{T(y)}{E(y,0)[\mathit{Loop}/\mathit{loop}]}}{R(z)}.
\end{array}
\]
In the first step, unrolling $\mathit{Loop}$ and substituting $3$ for $x$
gives $\cret{\rho.H}$, and derived bind passes $H$ to the continuation, so
$D_1$ applies $H$ directly.  After the second step, both runtime terms are
well-formed.  Their first components apply the literals $\mathit{Loop}$ and
$H$.  At $\qq$ and $\rr$, which have entered neither, the inner term
projects as $D_1=\unfold{C}$ does, as condition (iii) of the outer term
requires, and $D_2=\unfold{D_1}$ satisfies that condition for the inner
term.  As in \S\ref{sec:c3-example}, $\pp$ may instead run ahead
and compute $y$ and $b$ through both runtime terms before the others enter.

From $D_2$ the run proceeds as in \S\ref{sec:c3-example}, with each round
taking six entries, two applications times three roles, where
Chapter~\ref{ch:fun} took three.  In the first round the guard is $\ffalse$,
and the else branch applies $\mathit{Loop}$ to $\pp.6$ and the resulting
closure to $\qq.k'$.  In the second, $2\cdot 6=12>10$, the value $12$
travels from $\pp$ to $\qq$ to $\rr$, and the run ends in $\cret{\rr.13}$.

%======================================================================
\chapter*{Summary}
\addcontentsline{toc}{chapter}{Summary}
%======================================================================

Each of the first four chapters added one feature, and the fifth removed
what the fourth did not need.  All of them kept the correspondence between
a choreography and its projection exact.  The table lists what each chapter
cost in the proofs.

\begin{center}
\renewcommand{\arraystretch}{1.25}
\begin{tabular}{p{0.2\textwidth}p{0.36\textwidth}p{0.34\textwidth}}
  \textbf{Chapter} & \textbf{New syntax and rules} & \textbf{New lemmas}\\\hline
  \ref{ch:core} Core & $\bind{x}{A}{C}$; \rn{Local}, \rn{Com}, \rn{Delay} &
    substitution, locality, commutation with substitution, frames,
    inversion, network confluence\\
  \ref{ch:cond} Conditionals & conditional action; \rn{Cast}, \rn{N-Cast} &
    results held by participants; erasure; typing of bind; projection
    commutes with bind\\
  \ref{ch:fun} Functions & calls; runtime term; \rn{Call}, \rn{Enter}, \rn{Inside}, \rn{N-Call} &
    well-formedness and its preservation; substitution and well-formedness;
    unfolding; view; entering\\
  \ref{ch:ho} Higher-order & literals located at role sets; application $\rho.F(\vec\alpha)$ &
    substitution, locality and both commutation lemmas restated for
    location sets; unfolding re-proved\\
  \ref{ch:min} Minimal & one sort of values; one-argument $\lambda$;
    $\kw{fix}$, unrolled on entry &
    unique types; canonical forms; unrolling; unfolding restated
\end{tabular}
\end{center}

Some decisions come back in every chapter.  A role that is not involved in
an action gets nothing for it, and with the free monad's normal form this
makes one action one step at each level and the frame lemmas trivial.  Only
what appears in a label has to be a value.  Communication and multicast need
values because their labels carry them, local computation needs a closed
expression because it evaluates it, and entering a call needs nothing; \rn{Delay}
alone then blocks variables bound by pending actions.  Finally, branches and
function bodies never step before they are selected or entered.  They stay
source syntax, so roles, receivers and annotations are stable, bind is
applied only to source terms, and the one runtime term needs only one
invariant.  Chapter~\ref{ch:related} discusses what these decisions cost
compared with other approaches, chiefly the coarser knowledge of choice
that the multicast provides.

\begin{thebibliography}{99}
\bibitem{cc}
Lu\'is Cruz-Filipe and Fabrizio Montesi.
\newblock A core model for choreographic programming.
\newblock \emph{Theoretical Computer Science}, 802:38--66, 2020.

\bibitem{pc17}
Lu\'is Cruz-Filipe and Fabrizio Montesi.
\newblock Procedural choreographic programming.
\newblock In \emph{Formal Techniques for Distributed Objects, Components,
  and Systems (FORTE 2017)}, volume 10321 of \emph{LNCS}, pages 92--107.
  Springer, 2017.

\bibitem{montesi-book}
Fabrizio Montesi.
\newblock \emph{Introduction to Choreographies}.
\newblock Cambridge University Press, 2023.

\bibitem{cmp-jar23}
Lu\'is Cruz-Filipe, Fabrizio Montesi, and Marco Peressotti.
\newblock A formal theory of choreographic programming.
\newblock \emph{Journal of Automated Reasoning}, 67(2):21, 2023.

\bibitem{chorlambda22}
Lu\'is Cruz-Filipe, Eva Graversen, Lovro Lugovi\'c, Fabrizio Montesi, and
  Marco Peressotti.
\newblock Functional choreographic programming.
\newblock In \emph{Theoretical Aspects of Computing (ICTAC 2022)}, volume
  13572 of \emph{LNCS}, pages 212--237. Springer, 2022.

\bibitem{chorlambda-arxiv}
Lu\'is Cruz-Filipe, Eva Graversen, Lovro Lugovi\'c, Fabrizio Montesi, and
  Marco Peressotti.
\newblock Functional choreographic programming.
\newblock Full version, arXiv:2111.03701, 2022.

\bibitem{chorlambda23}
Lu\'is Cruz-Filipe, Eva Graversen, Lovro Lugovi\'c, Fabrizio Montesi, and
  Marco Peressotti.
\newblock Modular compilation for higher-order functional choreographies.
\newblock In \emph{37th European Conference on Object-Oriented Programming
  (ECOOP 2023)}, volume 263 of \emph{LIPIcs}, pages 7:1--7:37, 2023.

\bibitem{polychor24}
Eva Graversen, Andrew~K. Hirsch, and Fabrizio Montesi.
\newblock Alice or {B}ob?: Process polymorphism in choreographies.
\newblock \emph{Journal of Functional Programming}, 34:e1, 2024.

\bibitem{pirouette22}
Andrew~K. Hirsch and Deepak Garg.
\newblock Pirouette: Higher-order typed functional choreographies.
\newblock \emph{Proceedings of the ACM on Programming Languages}, 6(POPL):23:1--23:27, 2022.

\bibitem{relax25}
Dan Plyukhin, Xueying Qin, and Fabrizio Montesi.
\newblock Relax! {T}he semilenient core of choreographic programming
  (functional pearl).
\newblock \emph{Proceedings of the ACM on Programming Languages},
  9(ICFP):947--973, 2025.

\bibitem{ozone24}
Dan Plyukhin, Marco Peressotti, and Fabrizio Montesi.
\newblock Ozone: Fully out-of-order choreographies.
\newblock In \emph{38th European Conference on Object-Oriented Programming
  (ECOOP 2024)}, volume 313 of \emph{LIPIcs}, pages 31:1--31:28, 2024.

\bibitem{haschor23}
Gan Shen, Shun Kashiwa, and Lindsey Kuper.
\newblock {HasChor}: Functional choreographic programming for all
  (functional pearl).
\newblock \emph{Proceedings of the ACM on Programming Languages},
  7(ICFP):541--565, 2023.

\bibitem{enclaves23}
Shun Kashiwa, Gan Shen, Soroush Zare, and Lindsey Kuper.
\newblock Portable, efficient, and practical library-level choreographic
  programming.
\newblock arXiv:2311.11472, 2023.

\bibitem{multichor25}
Mako Bates, Shun Kashiwa, Syed Jafri, Gan Shen, Lindsey Kuper, and
  Joseph~P. Near.
\newblock Efficient, portable, census-polymorphic choreographic programming.
\newblock \emph{Proceedings of the ACM on Programming Languages},
  9(PLDI):1143--1166, 2025.

\bibitem{ebsmall24}
Mako Bates and Joseph~P. Near.
\newblock We know {I} know you know; choreographic programming with
  multicast and multiply located values.
\newblock arXiv:2403.05417, 2024.

\bibitem{shenkuper24}
Gan Shen and Lindsey Kuper.
\newblock Toward verified library-level choreographic programming with
  algebraic effects.
\newblock arXiv:2407.06509, 2024.

\bibitem{bohler26}
Timon B\"ohler, Simon Daniel, David Richter, Pascal Weisenburger, and Mira
  Mezini.
\newblock Mechanizing choreographic programs and {H}oare logic with state
  transformers.
\newblock In \emph{Type-Driven Development (TyDe 2026)}, to appear.
  arXiv:2608.16346.
\end{thebibliography}

\end{document}
